% =============================================================================
% ΚΕΦΑΛΑΙΟ 4: ΥΛΟΠΟΙΗΣΗ
% =============================================================================
% Αυτό το αρχείο περιέχει το πλήρες Κεφάλαιο 4 της διπλωματικής εργασίας
% "Μοντελοποίηση Αγορών Ηλεκτρικής Ενέργειας με Julia/JuMP"
% 
% Για ενσωμάτωση στο thesis.tex, χρησιμοποιήστε: % =============================================================================
% ΚΕΦΑΛΑΙΟ 4: ΥΛΟΠΟΙΗΣΗ
% =============================================================================
% Αυτό το αρχείο περιέχει το πλήρες Κεφάλαιο 4 της διπλωματικής εργασίας
% "Μοντελοποίηση Αγορών Ηλεκτρικής Ενέργειας με Julia/JuMP"
% 
% Για ενσωμάτωση στο thesis.tex, χρησιμοποιήστε: % =============================================================================
% ΚΕΦΑΛΑΙΟ 4: ΥΛΟΠΟΙΗΣΗ
% =============================================================================
% Αυτό το αρχείο περιέχει το πλήρες Κεφάλαιο 4 της διπλωματικής εργασίας
% "Μοντελοποίηση Αγορών Ηλεκτρικής Ενέργειας με Julia/JuMP"
% 
% Για ενσωμάτωση στο thesis.tex, χρησιμοποιήστε: % =============================================================================
% ΚΕΦΑΛΑΙΟ 4: ΥΛΟΠΟΙΗΣΗ
% =============================================================================
% Αυτό το αρχείο περιέχει το πλήρες Κεφάλαιο 4 της διπλωματικής εργασίας
% "Μοντελοποίηση Αγορών Ηλεκτρικής Ενέργειας με Julia/JuMP"
% 
% Για ενσωμάτωση στο thesis.tex, χρησιμοποιήστε: \input{chapter4_implementation}
% =============================================================================

\chapter{Υλοποίηση}
\label{ch:implementation}

Στο παρόν κεφάλαιο παρουσιάζεται η υλοποίηση του συστήματος προσομοίωσης της Προημερήσιας Αγοράς ηλεκτρικής ενέργειας. Η ανάλυση ακολουθεί τη ροή επεξεργασίας δεδομένων από τα μοντέλα οντοτήτων, μέσω του επιλυτή Unit Commitment και της μηχανής στρατηγικής προσφορών, έως την εκκαθάριση της αγοράς με τη μέθοδο MPCC. Για κάθε συνιστώσα παρέχονται αντιπροσωπευτικά αποσπάσματα κώδικα που επιδεικνύουν τις βασικές αλγοριθμικές επιλογές. Ο πλήρης κώδικας είναι διαθέσιμος στο αποθετήριο του έργου\footnote{\url{https://github.com/silentech-inc/energy-markets}}.


% =============================================================================
\section{Μοντελοποίηση Οντοτήτων}
\label{sec:entity-modeling}

Η υλοποίηση βασίζεται σε σαφώς ορισμένες δομές δεδομένων για την αναπαράσταση των οντοτήτων του συστήματος ηλεκτρικής ενέργειας. Οι δομές αυτές ορίζονται ως αμετάβλητα \texttt{struct} στη Julia, εξασφαλίζοντας ακεραιότητα δεδομένων και αποτελεσματικότητα.

\subsection{Μονάδα Παραγωγής (Generator)}
\label{subsec:generator-struct}

Η δομή \texttt{Generator} αναπαριστά μια θερμική μονάδα παραγωγής με όλα τα τεχνικά και οικονομικά χαρακτηριστικά που απαιτούνται για το πρόβλημα Unit Commitment:

\begin{samepage}
\begin{minted}[fontsize=\small,linenos,frame=lines]{julia}
struct Generator
    code::String           # Μοναδικός κωδικός (EIC code)
    name::String           # Όνομα μονάδας
    fuel_type::Symbol      # Τύπος καυσίμου (:lignite, :gas, :hydro, κλπ)
    location::String       # Τοποθεσία
    p_max::Float64         # Μέγιστη ισχύς (MW)
    p_min::Float64         # Ελάχιστη ισχύς (MW)
    bidding_zone::String   # Ζώνη προσφοράς (π.χ. "GR")
    marginal_cost::Float64 # Οριακό κόστος (€/MWh)
    # Προαιρετικές παράμετροι από εξαγωγή ιστορικών δεδομένων
    # nothing = χρήση προεπιλεγμένων τιμών ανά τύπο καυσίμου
    ramp_up::Union{Float64, Nothing}    # Ρυθμός αύξησης (κλάσμα p_max/ώρα)
    ramp_down::Union{Float64, Nothing}  # Ρυθμός μείωσης (κλάσμα p_max/ώρα)
    min_uptime::Union{Int, Nothing}     # Ελάχιστος χρόνος λειτουργίας (ώρες)
    min_downtime::Union{Int, Nothing}   # Ελάχιστος χρόνος αδράνειας (ώρες)

    # Κατασκευαστής με προεπιλεγμένες τιμές nothing
    function Generator(code, name, fuel_type, location, p_max, p_min,
                       bidding_zone, marginal_cost,
                       ramp_up=nothing, ramp_down=nothing,
                       min_uptime=nothing, min_downtime=nothing)
        new(code, name, fuel_type, location, p_max, p_min,
            bidding_zone, marginal_cost,
            ramp_up, ramp_down, min_uptime, min_downtime)
    end
end
\end{minted}
\end{samepage}

Ο τύπος καυσίμου (\texttt{fuel\_type}) αποθηκεύεται ως \texttt{Symbol} για αποτελεσματική σύγκριση και χρησιμοποιείται για την ανάκτηση παραμέτρων ειδικών ανά τύπο καυσίμου. Τα πεδία \texttt{ramp\_up}, \texttt{ramp\_down}, \texttt{min\_uptime} και \texttt{min\_downtime} ορίζονται ως \texttt{Union\{T, Nothing\}}, γεγονός που επιτρέπει τη χρήση της τιμής \texttt{nothing} ως σημαίας (sentinel) για την ενεργοποίηση των προεπιλεγμένων τιμών ανά τύπο καυσίμου. Ο εσωτερικός κατασκευαστής (inner constructor) ορίζει αυτά τα πεδία με προεπιλεγμένη τιμή \texttt{nothing}, διατηρώντας τη συμβατότητα με τον υπάρχοντα κώδικα. Οι τεχνικές παράμετροι εξάγονται αυτόματα από ιστορικά δεδομένα παραγωγής, όπως περιγράφεται στην Ενότητα~\ref{subsec:parameter-inference}.

Η ανάκτηση μονάδων παραγωγής από τη βάση δεδομένων πραγματοποιείται με τη συνάρτηση \texttt{get\_generators}, η οποία υποστηρίζει φιλτράρισμα μη διαθεσιμοτήτων:

\begin{samepage}
\begin{minted}[fontsize=\small,linenos,frame=lines]{julia}
function get_generators(bidding_zone::String, day::Date; 
                        exclude_unavailable::Bool=true)
    # Ερώτημα SQL με LEFT JOIN στον πίνακα unavailability
    # για αποκλεισμό ή μείωση ισχύος μονάδων με διακοπές
    
    query = """
    WITH active_outages AS (
        SELECT asset_code, MIN(available_capacity_mw) as available_capacity_mw
        FROM entsoe.unavailability_of_production_and_generation_units
        WHERE status = 'Active'
          AND start_outage_utc::timestamp <= \$2::date + interval '1 day'
          AND end_outage_utc::timestamp >= \$2::date
        GROUP BY asset_code
    )
    SELECT g.*, COALESCE(o.available_capacity_mw, g.generation_unit_installed_capacity_mw) as p_max
    FROM entsoe.production_and_generation_units g
    LEFT JOIN active_outages o ON g.generation_unit_code = o.asset_code
    WHERE g.area_map_code = \$1
      AND g.production_unit_status = 'COMMISSIONED'
      AND (o.asset_code IS NULL OR o.available_capacity_mw > 0)
    """
    # ... επεξεργασία αποτελεσμάτων
end
\end{minted}
\end{samepage}

\subsection{Παράμετροι Τύπου Καυσίμου}
\label{subsec:fuel-type-parameters}

Οι τεχνικοί περιορισμοί των μονάδων παραγωγής διαφέρουν σημαντικά ανάλογα με τον τύπο καυσίμου. Η δομή \texttt{FuelTypeParameters} συγκεντρώνει αυτές τις παραμέτρους:

\begin{samepage}
\begin{minted}[fontsize=\small,linenos,frame=lines]{julia}
struct FuelTypeParameters
    # Χαρακτηριστικά εκκίνησης (σε ΩΡΕΣ)
    cold_startup_time::Int          # Ώρες για ψυχρή εκκίνηση
    warm_startup_time::Int          # Ώρες για θερμή εκκίνηση
    hot_startup_time::Int           # Ώρες για καυτή εκκίνηση
    # Ελάχιστοι χρόνοι λειτουργίας (σε ΩΡΕΣ)
    min_uptime::Int                 # Ελάχιστες ώρες λειτουργίας
    min_downtime::Int               # Ελάχιστες ώρες αδράνειας
    # Ρυθμοί μεταβολής (κλάσμα δυναμικότητας ανά ΩΡΑ)
    ramp_up_rate::Float64           # Ρυθμός αύξησης (κλάσμα/ώρα)
    ramp_down_rate::Float64         # Ρυθμός μείωσης (κλάσμα/ώρα)
    # Λειτουργική ευελιξία
    min_load_factor::Float64        # Ελάχιστο φορτίο (κλάσμα)
    part_load_efficiency::Float64   # Απόδοση στο ελάχιστο φορτίο
    # Κατώφλια θερμοκρασιακής μετάβασης (ΩΡΕΣ εκτός)
    warm_threshold::Int             # Ώρες αδράνειας → θερμή κατάσταση
    cold_threshold::Int             # Ώρες αδράνειας → ψυχρή κατάσταση
    # Οικονομικές παράμετροι
    startup_cost_multiplier::Float64  # Πολλαπλασιαστής κόστους εκκίνησης
    no_load_cost_fraction::Float64    # Κόστος χωρίς φορτίο (κλάσμα)
end
\end{minted}
\end{samepage}

Σημαντικό σχεδιαστικό χαρακτηριστικό είναι ότι όλες οι χρονικές παράμετροι αποθηκεύονται σε \textbf{ώρες} (όχι σε περιόδους). Ο επιλυτής UC μετατρέπει αυτόματα σε περιόδους βάσει της πραγματικής χρονικής ανάλυσης των δεδομένων ($\text{periods\_per\_hour} = 60 / \text{resolution\_minutes}$). Τα πεδία \texttt{warm\_threshold} και \texttt{cold\_threshold} χρησιμοποιούνται τόσο στον προσδιορισμό θερμοκρασιακής κατάστασης εκκίνησης (Ενότητα~\ref{subsec:initial-conditions}) όσο και στους αντίστοιχους περιορισμούς του μοντέλου UC. Τα πεδία \texttt{startup\_cost\_multiplier} και \texttt{no\_load\_cost\_fraction} τροφοδοτούν τον υπολογισμό κόστους εκκίνησης και κόστους χωρίς φορτίο στην αντικειμενική συνάρτηση (Ενότητα~\ref{subsec:uc-balance-objective}).

Η συνάρτηση \texttt{get\_fuel\_type\_parameters} επιστρέφει τις κατάλληλες παραμέτρους για κάθε τύπο καυσίμου. Για παράδειγμα, οι λιγνιτικές μονάδες έχουν μεγαλύτερους χρόνους εκκίνησης και χαμηλότερους ρυθμούς μεταβολής σε σύγκριση με τις μονάδες συνδυασμένου κύκλου φυσικού αερίου.

\subsection{Φορτίο και Ανανεώσιμες Πηγές}
\label{subsec:load-renewables}

Η ζήτηση ηλεκτρικής ενέργειας αναπαρίσταται με τη δομή \texttt{Load}:

\begin{samepage}
\begin{minted}[fontsize=\small,linenos,frame=lines]{julia}
struct Load
    timeslot::String           # Χρονοθυρίδα (π.χ. "20240615-0000")
    resolution_code::String    # Χρονική ανάλυση ("PT15M", "PT60M")
    bidding_zone::String       # Ζώνη προσφοράς
    value::Float64             # Κατανάλωση (MW)
end
\end{minted}
\end{samepage}

Οι προβλέψεις παραγωγής από ΑΠΕ αναπαρίστανται με τη δομή \texttt{RenewablesGenerationForecast}:

\begin{samepage}
\begin{minted}[fontsize=\small,linenos,frame=lines]{julia}
struct RenewablesGenerationForecast
    date_time::String                    # Χρονοσφραγίδα
    resolution_code::String              # Χρονική ανάλυση
    bidding_zone::String                 # Ζώνη προσφοράς
    production_type::String              # "Solar" ή "Wind Onshore/Offshore"
    aggregated_generation_forecast::Float64  # Πρόβλεψη (MW)
end
\end{minted}
\end{samepage}

\subsection{Εξαγωγή Τεχνικών Παραμέτρων από Ιστορικά Δεδομένα}
\label{subsec:parameter-inference}

Οι τεχνικές παράμετροι των μονάδων παραγωγής (ρυθμοί μεταβολής, ελάχιστη ισχύς, ελάχιστοι χρόνοι λειτουργίας/αδράνειας) δεν είναι διαθέσιμες στα δημόσια δεδομένα του ENTSO-E. Για την αντιμετώπιση αυτού του περιορισμού, αναπτύχθηκε ένας μηχανισμός αυτόματης εξαγωγής (inference) αυτών των παραμέτρων από ιστορικά δεδομένα πραγματικής παραγωγής.

\paragraph{Πηγή Δεδομένων}
Τα ιστορικά δεδομένα παραγωγής ανακτώνται από τον πίνακα \texttt{actual\_generation\_output\_per\_generation\_unit} του ENTSO-E, ο οποίος περιέχει τη μετρούμενη παραγωγή κάθε μονάδας ανά 15 λεπτά ή ανά ώρα. Η συνάρτηση \texttt{get\_historical\_generation} ανακτά δεδομένα για παράθυρο 3 μηνών πριν από την επιλεγμένη ημερομηνία.

\paragraph{Εξαγωγή Ρυθμών Μεταβολής}
Οι ρυθμοί μεταβολής (ramp rates) υπολογίζονται από τη συνάρτηση \texttt{infer\_ramp\_rates}:

\begin{samepage}
\begin{minted}[fontsize=\small,linenos,frame=lines]{julia}
function infer_ramp_rates(historical_data::DataFrame, p_max::Float64;
                          percentile::Float64=0.95)
    if nrow(historical_data) < MIN_DATA_POINTS_FOR_RAMP_INFERENCE
        return (ramp_up=nothing, ramp_down=nothing)
    end

    # Ταξινόμηση κατά χρόνο και υπολογισμός μεταβολών
    sort!(historical_data, :date_time_utc)
    gen_values = historical_data.actual_generation_output_mw
    deltas = diff(gen_values)

    # Διαχωρισμός σε αυξήσεις και μειώσεις
    ramp_ups = filter(x -> x > 0, deltas)
    ramp_downs = map(abs, filter(x -> x < 0, deltas))

    # Κανονικοποίηση σε ωριαίο ρυθμό βάσει ανάλυσης δεδομένων
    resolution_code = historical_data.resolution_code[1]
    resolution_minutes = parse_resolution_to_minutes(resolution_code)
    hourly_factor = 60.0 / resolution_minutes

    # 95ο εκατοστημόριο, μετατροπή σε κλάσμα p_max/ώρα
    ramp_up = isempty(ramp_ups) ? nothing :
        (quantile(ramp_ups, percentile) * hourly_factor) / p_max
    ramp_down = isempty(ramp_downs) ? nothing :
        (quantile(ramp_downs, percentile) * hourly_factor) / p_max

    return (ramp_up=ramp_up, ramp_down=ramp_down)
end
\end{minted}
\end{samepage}

Οι ρυθμοί αποθηκεύονται ως κλάσμα της μέγιστης ισχύος ανά ώρα (π.χ. 0.20 σημαίνει ότι η μονάδα μπορεί να μεταβάλει την παραγωγή της κατά 20\% της $P^{\max}$ ανά ώρα), επιτρέποντας την εύκολη μετατροπή σε οποιαδήποτε χρονική ανάλυση.

\paragraph{Εξαγωγή Ελάχιστης Ισχύος}
Η ελάχιστη ισχύς ($P^{\min}$) εξάγεται αναλύοντας τις περιόδους σταθερής λειτουργίας, εξαιρώντας τις μεταβατικές καταστάσεις (εκκίνηση/τερματισμός):

\begin{samepage}
\begin{minted}[fontsize=\small,linenos,frame=lines]{julia}
function infer_p_min(historical_data::DataFrame, p_max::Float64,
                     fuel_type::Symbol; percentile::Float64=0.05)
    if nrow(historical_data) < MIN_DATA_POINTS_FOR_RAMP_INFERENCE
        return nothing
    end

    sort!(historical_data, :date_time_utc)
    gen_values = historical_data.actual_generation_output_mw
    deltas = diff(gen_values)

    # Ανίχνευση μεταβατικών καταστάσεων (|Δ| > 5% της p_max)
    ramp_threshold = 0.05 * p_max

    # Συλλογή σταθερών μη-μηδενικών τιμών
    stable_values = Float64[]
    for i in 2:(length(gen_values) - 1)
        value = gen_values[i]
        if value > 0 && abs(deltas[i-1]) < ramp_threshold &&
                        abs(deltas[i]) < ramp_threshold
            push!(stable_values, value)
        end
    end

    if length(stable_values) < 50
        return nothing  # Ανεπαρκή δεδομένα
    end

    # 5ο εκατοστημόριο της σταθερής λειτουργίας
    inferred_p_min = quantile(stable_values, percentile)

    # Εφαρμογή ορίων ανά τύπο καυσίμου (clamp σε MW)
    (min_frac, max_frac) = P_MIN_BOUNDS[get_p_min_bounds_category(fuel_type)]
    return clamp(inferred_p_min, min_frac * p_max, max_frac * p_max)
end
\end{minted}
\end{samepage}

\paragraph{Όρια ανά Τύπο Καυσίμου}
Για την αποφυγή μη ρεαλιστικών τιμών, εφαρμόζονται όρια (bounds) ανά τύπο καυσίμου. Τα όρια αυτά βασίζονται στα τεχνικά χαρακτηριστικά κάθε τεχνολογίας παραγωγής:

\begin{table}[htbp]
\centering
\caption{Όρια ελάχιστης ισχύος ($P^{\min}$) ανά τύπο καυσίμου}
\label{tab:p-min-bounds}
\begin{tabular}{@{}lcc@{}}
\toprule
\textbf{Τύπος Καυσίμου} & \textbf{Ελάχιστο (\%)} & \textbf{Μέγιστο (\%)} \\
\midrule
Λιγνίτης/Άνθρακας & 45 & 65 \\
Φυσικό αέριο (CCGT) & 35 & 55 \\
Φυσικό αέριο (OCGT) & 20 & 45 \\
Υδροηλεκτρικά/Μπαταρίες & 0 & 0 \\
Προεπιλογή & 30 & 60 \\
\bottomrule
\end{tabular}
\end{table}

Οι ευέλικτοι πόροι (υδροηλεκτρικά, μπαταρίες) λαμβάνουν $P^{\min} = 0$ χωρίς εξαγωγή, καθώς μπορούν να λειτουργήσουν σε οποιοδήποτε επίπεδο ισχύος. Η αντίστοιχη σταθερά \texttt{FLEXIBLE\_FUEL\_TYPES} περιλαμβάνει τους τύπους \texttt{Hydro Pumped Storage}, \texttt{Hydro Run-of-river}, \texttt{Hydro Water Reservoir}, \texttt{Battery} και \texttt{Other}. Αξίζει να σημειωθεί ότι οι προεπιλεγμένες τιμές \texttt{min\_load\_factor} στη δομή \texttt{FuelTypeParameters} ευθυγραμμίστηκαν με τα όρια εξαγωγής (π.χ. λιγνίτης: 0.45, ευέλικτοι πόροι: 0.0).

\paragraph{Εξαγωγή Ελάχιστων Χρόνων Λειτουργίας/Αδράνειας}
Οι ελάχιστοι χρόνοι υπολογίζονται αναλύοντας τους κύκλους λειτουργίας (on/off) της μονάδας:

\begin{samepage}
\begin{minted}[fontsize=\small,linenos,frame=lines]{julia}
function infer_uptime_downtime(historical_data::DataFrame, fuel_type::Symbol;
                               percentile::Float64=0.05)
    if nrow(historical_data) < MIN_DATA_POINTS_FOR_RAMP_INFERENCE
        return (nothing, nothing)
    end

    sort!(historical_data, :date_time_utc)
    gen_values = historical_data.actual_generation_output_mw

    # Υπολογισμός διάρκειας περιόδου σε ώρες
    resolution_minutes = parse_resolution_to_minutes(
        historical_data.resolution_code[1])
    period_hours = resolution_minutes / 60.0

    # Προσδιορισμός κατάστασης (on: > 1 MW, off: = 0)
    is_on = gen_values .> 1.0

    # Εντοπισμός κύκλων λειτουργίας
    on_durations, off_durations = Float64[], Float64[]
    current_state = is_on[1]
    current_duration = 1

    for i in 2:length(is_on)
        if is_on[i] == current_state
            current_duration += 1
        else
            duration_hours = current_duration * period_hours
            if current_state
                push!(on_durations, duration_hours)
            else
                push!(off_durations, duration_hours)
            end
            current_state = is_on[i]
            current_duration = 1
        end
    end

    # Φιλτράρισμα σύντομων glitches (< 2 περιόδων)
    min_dur = 2 * period_hours
    filter!(d -> d >= min_dur, on_durations)
    filter!(d -> d >= min_dur, off_durations)

    # Απαίτηση τουλάχιστον 5 κύκλων για αξιόπιστη εξαγωγή
    # Εφαρμογή ορίων ανά τύπο καυσίμου (βλ. Πίνακα~\ref{tab:uptime-downtime-bounds})
    bounds = UPTIME_DOWNTIME_BOUNDS[get_p_min_bounds_category(fuel_type)]
    min_uptime = length(on_durations) >= 5 ?
        round(Int, clamp(quantile(on_durations, percentile),
                         bounds[1]...)) : nothing
    min_downtime = length(off_durations) >= 5 ?
        round(Int, clamp(quantile(off_durations, percentile),
                         bounds[2]...)) : nothing

    return (min_uptime, min_downtime)
end
\end{minted}
\end{samepage}

Ομοίως με τα όρια $P^{\min}$, εφαρμόζονται όρια ανά τύπο καυσίμου για τους χρόνους λειτουργίας και αδράνειας, όπως παρουσιάζονται στον Πίνακα~\ref{tab:uptime-downtime-bounds}.

\begin{table}[htbp]
\centering
\caption{Όρια ελάχιστων χρόνων λειτουργίας/αδράνειας ανά τύπο καυσίμου (ώρες)}
\label{tab:uptime-downtime-bounds}
\begin{tabular}{@{}lcccc@{}}
\toprule
\textbf{Τύπος Καυσίμου} & \multicolumn{2}{c}{\textbf{Χρόνος Λειτουργίας}} & \multicolumn{2}{c}{\textbf{Χρόνος Αδράνειας}} \\
\cmidrule(lr){2-3} \cmidrule(lr){4-5}
& \textbf{Ελάχ.} & \textbf{Μέγ.} & \textbf{Ελάχ.} & \textbf{Μέγ.} \\
\midrule
Λιγνίτης/Άνθρακας & 8 & 48 & 4 & 24 \\
Φυσικό αέριο (CCGT) & 2 & 12 & 1 & 8 \\
Φυσικό αέριο (OCGT) & 1 & 4 & 1 & 4 \\
Προεπιλογή & 2 & 24 & 1 & 12 \\
\bottomrule
\end{tabular}
\end{table}

Η εφαρμογή ορίων διασφαλίζει ότι οι εξαχθείσες τιμές παραμένουν εντός τεχνικά ρεαλιστικού εύρους. Για μονάδες βάσης (π.χ. λιγνιτικές), η εξαγωγή χρόνων συχνά δεν είναι εφικτή λόγω σπάνιων κύκλων λειτουργίας/αδράνειας, οπότε χρησιμοποιούνται οι προεπιλεγμένες τιμές τύπου καυσίμου.

\paragraph{Ενοποίηση Εξαγωγής Παραμέτρων}
Η κεντρική συνάρτηση \texttt{infer\_parameters\_for\_generators} ενορχηστρώνει την εξαγωγή για ένα σύνολο μονάδων, εφαρμόζοντας διαφορετική λογική ανάλογα με τον τύπο καυσίμου:

\begin{samepage}
\begin{minted}[fontsize=\small,linenos,frame=lines]{julia}
function infer_parameters_for_generators(generators::Vector{Generator},
                                         day::Date)
    updated = Generator[]
    for gen in generators
        historical = get_historical_generation(gen.code, day)
        rates = infer_ramp_rates(historical, gen.p_max)

        if gen.fuel_type in FLEXIBLE_FUEL_TYPES
            # Ευέλικτοι πόροι: p_min = 0, χωρίς uptime/downtime
            push!(updated, Generator(gen.code, gen.name, gen.fuel_type,
                gen.location, gen.p_max, 0.0, gen.bidding_zone,
                gen.marginal_cost, rates.ramp_up, rates.ramp_down,
                nothing, nothing))
        else
            # Θερμικές μονάδες: εξαγωγή όλων των παραμέτρων
            p_min = infer_p_min(historical, gen.p_max, gen.fuel_type)
            (uptime, downtime) = infer_uptime_downtime(historical,
                                                        gen.fuel_type)
            push!(updated, Generator(gen.code, gen.name, gen.fuel_type,
                gen.location, gen.p_max,
                p_min !== nothing ? p_min : gen.p_min,
                gen.bidding_zone, gen.marginal_cost,
                rates.ramp_up, rates.ramp_down, uptime, downtime))
        end
    end
    return updated
end
\end{minted}
\end{samepage}

\paragraph{Αποθήκευση σε Κρυφή Μνήμη (Caching)}
Η εξαγωγή παραμέτρων από ιστορικά δεδομένα είναι υπολογιστικά ακριβή (περίπου 17 λεπτά για 39 μονάδες). Για τη βελτίωση της απόδοσης, οι εξαχθείσες παράμετροι αποθηκεύονται στον πίνακα \texttt{simulations.generator\_inferred\_parameters}:

\begin{samepage}
\begin{minted}[fontsize=\small,linenos,frame=lines]{julia}
# Ανάκτηση μονάδων με αυτόματη εξαγωγή και caching
generators = get_generators_with_inferred_params("GR", Date(2024, 6, 15);
    use_cache = true,       # Χρήση cached τιμών αν διαθέσιμες
    max_age_days = 7        # Ανανέωση cache κάθε 7 ημέρες
)
\end{minted}
\end{samepage}

Η χρήση της κρυφής μνήμης μειώνει τον χρόνο ανάκτησης από περίπου 17 λεπτά σε περίπου 2 δευτερόλεπτα---βελτίωση κατά 525 φορές.


\subsection{Αρχικές Συνθήκες Μονάδων Παραγωγής}
\label{subsec:initial-conditions}

Για τη ρεαλιστική μοντελοποίηση του προβλήματος Unit Commitment, η κατάσταση κάθε μονάδας κατά τη χρονική στιγμή $t=0$ (έναρξη του ορίζοντα βελτιστοποίησης) πρέπει να είναι γνωστή. Οι αρχικές συνθήκες καθορίζουν αν η μονάδα λειτουργούσε, σε ποιο επίπεδο ισχύος, και για πόσο χρόνο βρισκόταν στην τρέχουσα κατάστασή της. Αυτή η πληροφορία είναι απαραίτητη για τη σωστή επιβολή περιορισμών στις πρώτες περιόδους του ορίζοντα.

\paragraph{Δομή Αρχικών Συνθηκών}
Η δομή \texttt{InitialConditions} αναπαριστά την κατάσταση μιας μονάδας στο $t=0$:

\begin{samepage}
\begin{minted}[fontsize=\small,linenos,frame=lines]{julia}
struct InitialConditions
    is_on::Bool           # Λειτουργούσε η μονάδα στο t=0;
    output::Float64       # Ισχύς εξόδου σε MW (0 αν εκτός λειτουργίας)
    hours_on::Int         # Συνεχόμενες ώρες λειτουργίας
    hours_off::Int        # Συνεχόμενες ώρες αδράνειας
    thermal_state::Symbol # Θερμοκρασιακή κατάσταση (:hot, :warm, :cold)
end
\end{minted}
\end{samepage}

Η θερμοκρασιακή κατάσταση προσδιορίζεται από τον χρόνο αδράνειας: \texttt{:hot} αν $T_{\text{off}} \leq 8$ ώρες, \texttt{:warm} αν $8 < T_{\text{off}} \leq 48$ ώρες, και \texttt{:cold} αν $T_{\text{off}} > 48$ ώρες. Τα κατώφλια αυτά προσαρμόζονται ανά τύπο καυσίμου μέσω των αντίστοιχων πεδίων \texttt{warm\_threshold} και \texttt{cold\_threshold} της δομής \texttt{FuelTypeParameters}.

\paragraph{Εξαγωγή από Ιστορικά Δεδομένα}
Η κεντρική συνάρτηση \texttt{get\_initial\_conditions} εξάγει τις αρχικές συνθήκες για ένα σύνολο μονάδων χρησιμοποιώντας ένα μαζικό ερώτημα SQL (batch query) αντί $N$ ξεχωριστών ερωτημάτων. Η τεχνική αυτή επιτυγχάνει βελτίωση ταχύτητας κατά 36 φορές (από $\sim$9 λεπτά σε $\sim$15 δευτερόλεπτα για 40 μονάδες):

\begin{samepage}
\begin{minted}[fontsize=\small,linenos,frame=lines]{julia}
function get_initial_conditions(generators::Vector{Generator},
                                market_day::Date;
                                use_historical::Bool=true)
    # Μαζική ανάκτηση 72 ωρών ιστορικού για όλες τις μονάδες
    # σε ένα ερώτημα SQL (WHERE generation_unit_code = ANY($1))
    day_before = market_day - Day(1)
    start_utc = DateTime(day_before, Time(23, 0, 0))  # ~00:00 CET
    codes = [gen.code for gen in generators]
    all_historical = get_recent_generation_batch(codes, start_utc)

    conditions = Dict{String, InitialConditions}()
    for gen in generators
        gen_data = filter(r -> r.generation_unit_code == gen.code,
                          all_historical)
        ic = infer_initial_conditions_from_data(gen_data, gen.fuel_type)
        conditions[gen.code] = ic
    end
    return conditions
end
\end{minted}
\end{samepage}

Ο αλγόριθμος εξαγωγής αναλύει τα δεδομένα σε φθίνουσα χρονική σειρά: η πιο πρόσφατη μέτρηση $> 1$ MW δηλώνει ότι η μονάδα ήταν σε λειτουργία. Στη συνέχεια, μετρούνται οι συνεχόμενες περίοδοι στην ίδια κατάσταση για τον υπολογισμό του $T_{\text{on}}$ ή $T_{\text{off}}$. Σε περίπτωση απουσίας δεδομένων, χρησιμοποιούνται ευρετικές προεπιλεγμένες τιμές ανά τύπο καυσίμου: οι μονάδες βάσης (λιγνίτης, πυρηνικά) θεωρούνται σε λειτουργία, οι μονάδες μέσου φορτίου (CCGT) εκτός λειτουργίας αλλά θερμές, οι μονάδες αιχμής (OCGT) εκτός λειτουργίας και ψυχρές, και οι ευέλικτοι πόροι (υδροηλεκτρικά, μπαταρίες) εκτός λειτουργίας αλλά έτοιμοι.


\subsection{Τύποι Εντολών Αγοράς}
\label{subsec:market-orders}

Οι εντολές αγοράς υλοποιούνται ως ιεραρχία τύπων με αφηρημένο βασικό τύπο \texttt{MarketOrder}. Ο κύριος τύπος που χρησιμοποιείται είναι η απλή ωριαία εντολή \texttt{SimpleOrder}:

\begin{samepage}
\begin{minted}[fontsize=\small,linenos,frame=lines]{julia}
abstract type MarketOrder end

struct SimpleOrder <: MarketOrder
    type::Symbol           # :supply ή :demand
    price::Float64         # Τιμή (€/MWh)
    quantity::Float64      # Ποσότητα (MWh)
    zone::Symbol           # Ζώνη προσφοράς
    date_time::DateTime    # Χρόνος παράδοσης
    resolution_code::Int   # Ανάλυση σε λεπτά (60, 30, 15)
end
\end{minted}
\end{samepage}

Πέρα από τις απλές εντολές, υλοποιούνται και σύνθετοι τύποι όπως οι εντολές μπλοκ (\texttt{BlockOrder}), οι συνδεδεμένες εντολές μπλοκ (\texttt{LinkedBlockOrder}) και οι αποκλειστικές εντολές μπλοκ (\texttt{ExclusiveBlockOrder}), ακολουθώντας την προδιαγραφή του αλγορίθμου EUPHEMIA.


% =============================================================================
\section{Επιλυτής Unit Commitment}
\label{sec:uc-solver}

Ο επιλυτής Unit Commitment υλοποιείται στο module \texttt{UnitCommitment.jl} και αποτελεί το πρώτο βήμα της αλυσίδας υπολογισμών. Η κύρια συνάρτηση \texttt{solve\_unit\_commitment} δέχεται ως είσοδο τη ζώνη προσφοράς και την ημερομηνία, ανακτά τα απαραίτητα δεδομένα από τη βάση, διαμορφώνει το πρόβλημα βελτιστοποίησης και επιστρέφει τη βέλτιστη λύση.

\subsection{Δημιουργία Μοντέλου}
\label{subsec:uc-model-creation}

Η δημιουργία του μοντέλου JuMP ξεκινά με την επιλογή του επιλυτή και τον ορισμό των μεταβλητών απόφασης:

\begin{samepage}
\begin{minted}[fontsize=\small,linenos,frame=lines]{julia}
function solve_unit_commitment(bidding_zone::String, day::Date;
                               optimizer::String="auto",
                               mip_gap::Float64=0.01,
                               time_limit::Float64=600.0)
    # Φόρτωση δεδομένων
    generators = get_generators(bidding_zone, day)
    loads = get_loads(bidding_zone, day)
    renewables = get_generation_forecast_for_wind_and_solar(bidding_zone, day)
    initial_conditions = get_initial_conditions(generators, day)

    N = length(generators)  # Αριθμός μονάδων
    T = 24                  # Χρονικές περίοδοι
    periods_per_hour = 60 / resolution_minutes

    # Δημιουργία μοντέλου με ρύθμιση παραμέτρων επιλυτή
    model = Model(optimizer_func)
    set_silent(model)
    set_optimizer_attribute(model, MOI.RelativeGapTolerance(), mip_gap)
    set_optimizer_attribute(model, MOI.TimeLimitSec(), time_limit)

    # Εξαγωγή αρχικών συνθηκών σε πίνακες
    u0 = zeros(Int, N)      # Αρχική δέσμευση (0 ή 1)
    g0 = zeros(Float64, N)  # Αρχική παραγωγή (MW)
    T_on0 = zeros(Int, N)   # Ώρες ήδη σε λειτουργία
    T_off0 = zeros(Int, N)  # Ώρες ήδη εκτός λειτουργίας
    for (i, gen) in enumerate(generators)
        ic = initial_conditions[gen.code]
        u0[i] = ic.is_on ? 1 : 0
        g0[i] = ic.output
        T_on0[i] = ic.hours_on
        T_off0[i] = ic.hours_off
    end

    # Μεταβλητές παραγωγής (συνεχείς)
    @variable(model, 0 <= g[i=1:N, t=1:T] <= generators[i].p_max)

    # Μεταβλητές δέσμευσης (δυαδικές)
    @variable(model, u[i=1:N, t=1:T], Bin)  # Κατάσταση λειτουργίας
    @variable(model, v[i=1:N, t=1:T], Bin)  # Εκκίνηση
    @variable(model, z[i=1:N, t=1:T], Bin)  # Τερματισμός

    # Μεταβλητές για θερμοκρασιακές καταστάσεις εκκίνησης
    Θ = [:cold, :warm, :hot]
    @variable(model, v_θ[i=1:N, θ in Θ, t=1:T], Bin)

    # ... περαιτέρω μεταβλητές
end
\end{minted}
\end{samepage}

\subsection{Περιορισμοί Ισχύος}
\label{subsec:impl-uc-power-constraints}

Οι περιορισμοί ισχύος εξασφαλίζουν ότι η παραγωγή κάθε μονάδας βρίσκεται εντός των τεχνικών της ορίων:

\begin{samepage}
\begin{minted}[fontsize=\small,linenos,frame=lines]{julia}
# Ανώτατο όριο παραγωγής
@constraint(model, [i = 1:N, t = 1:T], 
    g[i, t] <= generators[i].p_max * u[i, t])

# Κατώτατο όριο παραγωγής με παραμέτρους ανά τύπο καυσίμου
for (i, gen) in enumerate(generators)
    params = fuel_params[i]
    # Το ελάχιστο είναι το μέγιστο μεταξύ p_min και του ελάχιστου 
    # φορτίου βάσει τύπου καυσίμου
    min_gen = max(gen.p_min, params.min_load_factor * gen.p_max)
    @constraint(model, [t = 1:T], g[i, t] >= min_gen * u[i, t])
end
\end{minted}
\end{samepage}

\subsection{Περιορισμοί Δέσμευσης}
\label{subsec:impl-uc-commitment-constraints}

Η σχέση μεταξύ κατάστασης λειτουργίας, εκκίνησης και τερματισμού υλοποιείται ως εξής, με ενσωμάτωση αρχικών συνθηκών στην περίοδο $t=1$:

\begin{samepage}
\begin{minted}[fontsize=\small,linenos,frame=lines]{julia}
# Σύνδεση t=1 με αρχική κατάσταση u0
if initial_conditions !== nothing
    @constraint(model, [i in 1:N],
        u[i, 1] == u0[i] + v[i, 1] - z[i, 1])
end

# Σύνδεση μεταβλητών δέσμευσης (t ≥ 2)
@constraint(model, [i in 1:N, t in 2:T],
    u[i, t] == u[i, t-1] + v[i, t] - z[i, t])

# Αποκλεισμός ταυτόχρονης εκκίνησης και τερματισμού
@constraint(model, [i = 1:N, t = 1:T], v[i, t] + z[i, t] <= 1)

# Ελάχιστος χρόνος λειτουργίας (ώρες → περίοδοι)
for (i, gen) in enumerate(generators)
    UT_hours = gen.min_uptime !== nothing ?
        gen.min_uptime : fuel_params[i].min_uptime
    UT = max(1, ceil(Int, UT_hours * periods_per_hour))
    if UT > 1
        @constraint(model, [t = UT:T],
            sum(v[i, τ] for τ in t-UT+1:t) <= u[i, t])
        # Αρχική συνθήκη: εναπομένων χρόνος λειτουργίας
        if initial_conditions !== nothing && u0[i] == 1
            remaining = max(0, UT - ceil(Int, T_on0[i] * periods_per_hour))
            for t in 1:min(remaining, T)
                @constraint(model, u[i, t] == 1)
            end
        end
    end
end

# Ελάχιστος χρόνος αδράνειας (ώρες → περίοδοι)
for (i, gen) in enumerate(generators)
    DT_hours = gen.min_downtime !== nothing ?
        gen.min_downtime : fuel_params[i].min_downtime
    DT = max(1, ceil(Int, DT_hours * periods_per_hour))
    if DT > 1
        @constraint(model, [t = DT:T],
            sum(z[i, τ] for τ in t-DT+1:t) <= 1 - u[i, t])
        # Αρχική συνθήκη: εναπομένων χρόνος αδράνειας
        if initial_conditions !== nothing && u0[i] == 0
            remaining = max(0, DT - ceil(Int, T_off0[i] * periods_per_hour))
            for t in 1:min(remaining, T)
                @constraint(model, u[i, t] == 0)
            end
        end
    end
end
\end{minted}
\end{samepage}

Η μετατροπή $\texttt{UT\_hours} \times \texttt{periods\_per\_hour}$ διασφαλίζει ότι οι περιορισμοί εφαρμόζονται σωστά ανεξάρτητα από τη χρονική ανάλυση (15, 30 ή 60 λεπτά). Οι αρχικές συνθήκες επιβάλλουν τη συνέχιση της τρέχουσας κατάστασης κάθε μονάδας: αν μια μονάδα λειτουργούσε αλλά δεν έχει συμπληρώσει τον ελάχιστο χρόνο λειτουργίας, εξαναγκάζεται να παραμείνει σε λειτουργία για τις εναπομένουσες περιόδους.

\subsection{Περιορισμοί Ρυθμού Μεταβολής}
\label{subsec:impl-uc-ramping}

Οι περιορισμοί ρυθμού μεταβολής (ramping) χρησιμοποιούν την τεχνική Big-M για την αποδέσμευση κατά τις μεταβατικές καταστάσεις:

\begin{samepage}
\begin{minted}[fontsize=\small,linenos,frame=lines]{julia}
# Υπολογισμός ρυθμών μεταβολής ανά μονάδα:
# εξαχθείσα τιμή (fraction/hour) ή προεπιλογή τύπου καυσίμου
# Μετατροπή σε MW/περίοδο: rate × p_max × period_hours
period_hours = resolution_minutes / 60.0
ramp_up = Float64[]
ramp_down = Float64[]
for (i, gen) in enumerate(generators)
    if gen.ramp_up !== nothing
        push!(ramp_up, gen.ramp_up * gen.p_max * period_hours)
    else
        push!(ramp_up, fuel_params[i].ramp_up_rate * gen.p_max * period_hours)
    end
    if gen.ramp_down !== nothing
        push!(ramp_down, gen.ramp_down * gen.p_max * period_hours)
    else
        push!(ramp_down, fuel_params[i].ramp_down_rate * gen.p_max * period_hours)
    end
end

# Big-M: κλιμακούμενο βάσει μεγέθους προβλήματος
max_p_max = maximum(gen.p_max for gen in generators)
M = 2.0 * max_p_max

# Περιορισμός αύξησης (t ≥ 2)
@constraint(model, [i in 1:N, t in 2:T],
    g[i, t] - g[i, t-1] <= ramp_up[i] + M * u_SU[i, t])

# Περιορισμός μείωσης (t ≥ 2)
@constraint(model, [i in 1:N, t in 2:T],
    g[i, t-1] - g[i, t] <= ramp_down[i] + M * u_SD[i, t])
\end{minted}
\end{samepage}

\subsection{Ισορροπία Ισχύος και Αντικειμενική Συνάρτηση}
\label{subsec:uc-balance-objective}

Ο περιορισμός ισορροπίας ισχύος εξασφαλίζει την κάλυψη της καθαρής ζήτησης:

\begin{samepage}
\begin{minted}[fontsize=\small,linenos,frame=lines]{julia}
# Περιορισμός ισορροπίας
@constraint(model, [t in 1:T],
    sum(g[i, t] for i in 1:N) == net_demand[t])
\end{minted}
\end{samepage}

Η αντικειμενική συνάρτηση αποτελείται από τρεις συνιστώσες κόστους: κόστος παραγωγής, κόστος εκκίνησης και κόστος χωρίς φορτίο. Πριν τη διαμόρφωση, υπολογίζονται οι παράμετροι κόστους:

\begin{samepage}
\begin{minted}[fontsize=\small,linenos,frame=lines]{julia}
# Κόστος εκκίνησης ανά μονάδα και θερμοκρασιακή κατάσταση (€)
# Βασικό κόστος = startup_cost_multiplier × marginal_cost × p_max
# Πολλαπλασιαστές: hot=1.0, warm=1.5, cold=2.5
C_SU = Dict{Tuple{Int,Symbol},Float64}()
for i in 1:N
    params = fuel_params[i]
    gen = generators[i]
    base = params.startup_cost_multiplier * gen.marginal_cost * gen.p_max
    C_SU[(i, :hot)]  = base * 1.0
    C_SU[(i, :warm)] = base * 1.5
    C_SU[(i, :cold)] = base * 2.5
end

# Κόστος χωρίς φορτίο ανά μονάδα (€/περίοδο)
# Πάγιο κόστος δέσμευσης ανεξάρτητο από παραγωγή
period_hours = resolution_minutes / 60.0
C_NL = Float64[]
for i in 1:N
    params = fuel_params[i]
    gen = generators[i]
    push!(C_NL, params.no_load_cost_fraction *
                gen.marginal_cost * gen.p_min * period_hours)
end

# Αντικειμενική: ελαχιστοποίηση συνολικού κόστους
@objective(model, Min,
    # 1. Κόστος παραγωγής (μεταβλητό)
    sum(generators[i].marginal_cost * g[i, t] for i in 1:N, t in 1:T)
    # 2. Κόστος εκκίνησης (εξαρτώμενο από θερμοκρασία)
    + sum(C_SU[(i, θ)] * v_θ[i, θ, t] for i in 1:N, θ in Θ, t in 1:T)
    # 3. Κόστος χωρίς φορτίο (πάγιο κόστος δέσμευσης)
    + sum(C_NL[i] * u[i, t] for i in 1:N, t in 1:T)
)
\end{minted}
\end{samepage}

Η τρισδιάστατη δομή κόστους αποτυπώνει ρεαλιστικά τα κίνητρα λειτουργίας: το κόστος εκκίνησης αποθαρρύνει τις συχνές εκκινήσεις/τερματισμούς (ιδιαίτερα τις ψυχρές εκκινήσεις, που κοστίζουν 2,5 φορές περισσότερο από τις θερμές), ενώ το κόστος χωρίς φορτίο αντανακλά τα πάγια λειτουργικά κόστη δέσμευσης μιας μονάδας ανεξάρτητα από το επίπεδο παραγωγής της.

\subsection{Επίλυση και Επιστροφή Αποτελεσμάτων}
\label{subsec:uc-solution}

Μετά την επίλυση, τα αποτελέσματα επιστρέφονται σε δομημένη μορφή:

\begin{samepage}
\begin{minted}[fontsize=\small,linenos,frame=lines]{julia}
optimize!(model)

if termination_status(model) == OPTIMAL
    return (
        status = OPTIMAL,
        generators = generators,
        generation = value.(g),       # Πίνακας παραγωγής N×T
        commitment = value.(u),       # Πίνακας δέσμευσης N×T
        startup = value.(v),          # Πίνακας εκκινήσεων N×T
        shutdown = value.(z),         # Πίνακας τερματισμών N×T
        objective_value = objective_value(model),
        solve_time = solve_time
    )
else
    return (status = termination_status(model),)
end
\end{minted}
\end{samepage}

\subsection{Αποθήκευση Αποτελεσμάτων σε Κρυφή Μνήμη}
\label{subsec:uc-caching}

Η επίλυση του προβλήματος Unit Commitment απαιτεί 10--15 λεπτά ανά ζώνη. Για την αποφυγή επαναληπτικών υπολογισμών, τα αποτελέσματα αποθηκεύονται σε τρεις πίνακες PostgreSQL (\texttt{uc\_results}, \texttt{uc\_generation}, \texttt{uc\_net\_demand}) με κλειδί \texttt{(bidding\_zone, market\_date, code\_version)}. Η συνάρτηση \texttt{solve\_unit\_commitment} ελέγχει αυτόματα τη διαθεσιμότητα αποθηκευμένων αποτελεσμάτων:

\begin{samepage}
\begin{minted}[fontsize=\small,linenos,frame=lines]{julia}
function solve_unit_commitment(bidding_zone::String, day::Date;
                               optimizer::String="auto",
                               use_cache::Bool=true,
                               force_rerun::Bool=false,
                               # ... λοιπές παράμετροι
                               )
    # Έλεγχος κρυφής μνήμης (εκτός αν απενεργοποιημένη)
    if use_cache && !force_rerun
        cached = load_uc_results(bidding_zone, day)
        if cached !== nothing
            return cached   # Επιστροφή σε ~2 δευτερόλεπτα
        end
    end

    # ... κανονική επίλυση (10-15 λεπτά) ...

    # Αποθήκευση αποτελεσμάτων (delete-before-insert)
    if use_cache
        save_uc_results(solution, bidding_zone, day)
    end
    return solution
end
\end{minted}
\end{samepage}

Η παράμετρος \texttt{force\_rerun} διαδίδεται μέσω ολόκληρης της αλυσίδας κλήσεων (\texttt{generate\_energy\_prices} $\to$ \texttt{create\_typed\_order\_book} $\to$ \texttt{generate\_market\_orders\_from\_uc} $\to$ \texttt{solve\_unit\_commitment}), επιτρέποντας τον επανυπολογισμό αποτελεσμάτων όταν αλλάξει ο κώδικας ή τα δεδομένα εισόδου. Η αποθήκευση ανέρχεται σε περίπου 195 KB ανά ζώνη/ημέρα---κόστος αμελητέο σε σύγκριση με την εξοικονόμηση χρόνου.


% =============================================================================
\section{Μηχανή Στρατηγικής Προσφορών}
\label{sec:bidding-strategy-engine}

Η μηχανή στρατηγικής προσφορών υλοποιείται στο module \texttt{BiddingStrategy.jl} και μετατρέπει τα αποτελέσματα του Unit Commitment σε εντολές αγοράς. Η κύρια συνάρτηση \texttt{generate\_market\_orders\_from\_uc} υποστηρίζει τρεις διακριτές στρατηγικές.

\subsection{Δομή Αποτελέσματος}
\label{subsec:bidding-result}

Τα αποτελέσματα της μετατροπής επιστρέφονται στη δομή \texttt{UCToBidsResult}:

\begin{samepage}
\begin{minted}[fontsize=\small,linenos,frame=lines]{julia}
struct UCToBidsResult
    supply_orders::Vector{SimpleOrder}   # Εντολές πώλησης
    demand_orders::Vector{SimpleOrder}   # Εντολές αγοράς
    bidding_zone::Symbol                 # Ζώνη προσφοράς
    day::Date                            # Ημερομηνία
    total_supply_quantity::Float64       # Συνολική προσφερόμενη ποσότητα
    total_demand_quantity::Float64       # Συνολική ζητούμενη ποσότητα
    success::Bool                        # Επιτυχία μετατροπής
    message::String                      # Μήνυμα κατάστασης
end
\end{minted}
\end{samepage}

\subsection{Στρατηγική Δεσμευμένων Μονάδων}
\label{subsec:committed-only-strategy}

Η συντηρητική στρατηγική \texttt{:committed\_only} δημιουργεί προσφορές μόνο για τις μονάδες που έχουν δεσμευτεί στη λύση του UC:

\begin{samepage}
\begin{minted}[fontsize=\small,linenos,frame=lines]{julia}
function apply_committed_only_strategy!(orders, uc_solution, generators, 
                                        day, markup_factor)
    for (i, gen) in enumerate(generators)
        for t in 1:24
            # Έλεγχος αν η μονάδα είναι δεσμευμένη
            if uc_solution.commitment[i, t] > COMMITMENT_THRESHOLD
                generation = uc_solution.generation[i, t]
                if generation > GENERATION_THRESHOLD
                    # Δημιουργία εντολής πώλησης
                    price = gen.marginal_cost * markup_factor
                    dt = DateTime(day, Time(t-1, 0))
                    push!(orders, SimpleOrder(
                        :supply, price, generation, 
                        Symbol(gen.bidding_zone), dt, 60
                    ))
                end
            end
        end
    end
end
\end{minted}
\end{samepage}

\subsection{Στρατηγική Μέγιστης Διαθεσιμότητας}
\label{subsec:all-available-max}

Η στρατηγική \texttt{:all\_available\_max} προσφέρει όλες τις μονάδες στη μέγιστη δυναμικότητά τους, ανεξάρτητα από τη λύση του UC. Αν και μη ρεαλιστική, χρησιμεύει ως βάση σύγκρισης:

\begin{samepage}
\begin{minted}[fontsize=\small,linenos,frame=lines]{julia}
function apply_all_available_max_strategy!(orders, generators, day, markup_factor)
    for gen in generators
        for t in 1:24
            price = gen.marginal_cost * markup_factor
            dt = DateTime(day, Time(t-1, 0))
            push!(orders, SimpleOrder(
                :supply, price, gen.p_max,
                Symbol(gen.bidding_zone), dt, 60
            ))
        end
    end
end
\end{minted}
\end{samepage}

\subsection{Στρατηγική Ρεαλιστικής Διαθεσιμότητας}
\label{subsec:all-available-realistic}

Η στρατηγική \texttt{:all\_available\_realistic} συνδυάζει πληροφορίες από τη λύση UC με τη διαθεσιμότητα όλων των μονάδων:

\begin{samepage}
\begin{minted}[fontsize=\small,linenos,frame=lines]{julia}
function apply_all_available_realistic_strategy!(orders, uc_solution, 
                                                  generators, day, markup_factor)
    for (i, gen) in enumerate(generators)
        for t in 1:24
            # Καθορισμός ποσότητας βάσει κατάστασης δέσμευσης
            if uc_solution.commitment[i, t] > COMMITMENT_THRESHOLD
                # Δεσμευμένη μονάδα: χρήση πραγματικής παραγωγής
                quantity = uc_solution.generation[i, t]
            elseif gen.p_min < 0.1 * gen.p_max
                # Μη δεσμευμένη με χαμηλό ελάχιστο: 20% της δυναμικότητας
                quantity = DEFAULT_UNCOMMITTED_UNIT_FRACTION * gen.p_max
            else
                # Μη δεσμευμένη: ελάχιστη ισχύς
                quantity = gen.p_min
            end
            
            if quantity > GENERATION_THRESHOLD
                price = gen.marginal_cost * markup_factor
                dt = DateTime(day, Time(t-1, 0))
                push!(orders, SimpleOrder(
                    :supply, price, quantity,
                    Symbol(gen.bidding_zone), dt, 60
                ))
            end
        end
    end
end
\end{minted}
\end{samepage}

\subsection{Δημιουργία Εντολών Ζήτησης}
\label{subsec:demand-orders-creation}

Οι εντολές ζήτησης δημιουργούνται με βάση τα δεδομένα φορτίου, με τιμή ίση με το ανώτατο όριο της αγοράς (price-taking demand):

\begin{samepage}
\begin{minted}[fontsize=\small,linenos,frame=lines]{julia}
function create_demand_orders!(orders, loads, day, bidding_zone, price_cap)
    for load in loads
        if load.value > DEMAND_THRESHOLD
            hour = parse(Int, load.timeslot[end-3:end-2])
            dt = DateTime(day, Time(hour, 0))
            push!(orders, SimpleOrder(
                :demand, price_cap, load.value,
                Symbol(bidding_zone), dt, 60
            ))
        end
    end
end
\end{minted}
\end{samepage}


% =============================================================================
\section{Εκκαθάριση Αγοράς με MPCC}
\label{sec:mpcc-market-clearing}

Η εκκαθάριση της αγοράς υλοποιείται στο module \texttt{MPCC.jl} χρησιμοποιώντας τη μέθοδο Mathematical Programming with Complementarity Constraints. Ο αλγόριθμος δέχεται ένα βιβλίο εντολών (\texttt{MPCCOrderBook}) και επιστρέφει τις τιμές εκκαθάρισης και τους βαθμούς αποδοχής των εντολών.

\subsection{Δομή Βιβλίου Εντολών}
\label{subsec:order-book-structure}

Το βιβλίο εντολών οργανώνει τις εντολές αγοράς μαζί με τις παραμέτρους του προβλήματος:

\begin{samepage}
\begin{minted}[fontsize=\small,linenos,frame=lines]{julia}
struct MPCCOrderBook
    orders::Vector{MarketOrder}              # Όλες οι εντολές
    nodes::Vector{String}                    # Ζώνες προσφοράς
    periods::Vector{String}                  # Χρονικές περίοδοι
    price_limits::Tuple{Float64,Float64}     # (min, max) τιμή
    network_topology::Union{Nothing,TransferCapacity}  # Περιορισμοί δικτύου
end
\end{minted}
\end{samepage}

\subsection{Μεταβλητές Απόφασης}
\label{subsec:mpcc-variables}

Οι κύριες μεταβλητές του μοντέλου MPCC είναι οι τιμές αγοράς, οι βαθμοί αποδοχής και οι ροές μεταφοράς:

\begin{samepage}
\begin{minted}[fontsize=\small,linenos,frame=lines]{julia}
function solve_mpcc_market_clearing(order_book::MPCCOrderBook; 
                                    big_m::Float64=BIG_M_PARAMETER)
    model = Model(HiGHS.Optimizer)
    set_silent(model)
    
    # Τιμές αγοράς ανά ζώνη και περίοδο
    @variable(model, market_price[node in order_book.nodes, 
                                  period in order_book.periods])
    
    # Βαθμός αποδοχής κάθε εντολής [0, 1]
    @variable(model, 0 <= stepwise_acceptance[order_id in order_ids] <= 1)
    
    # Dual μεταβλητή για περιορισμούς συμπληρωματικότητας
    @variable(model, stepwise_dual[order_id in order_ids] >= 0)
    
    # Μεταβλητή έκτακτης αποκοπής φορτίου
    @variable(model, load_shed[node in order_book.nodes, 
                               period in order_book.periods] >= 0)
    
    # Ροές μεταφοράς (αν υπάρχει δίκτυο)
    if order_book.network_topology !== nothing
        zone_pairs = get_zone_pairs(order_book.network_topology)
        @variable(model, flow[(s,d) in zone_pairs, t in order_book.periods])
    end
end
\end{minted}
\end{samepage}

\subsection{Περιορισμοί Συμπληρωματικότητας}
\label{subsec:complementarity-constraints}

Οι περιορισμοί συμπληρωματικότητας εξασφαλίζουν ότι μια εντολή γίνεται αποδεκτή μόνο όταν η τιμή αγοράς είναι ευνοϊκή. Υλοποιούνται με τη μέθοδο Big-M:

\begin{samepage}
\begin{minted}[fontsize=\small,linenos,frame=lines]{julia}
# Υπολογισμός του δεξιού μέλους της dual συνθήκης
for (i, order) in enumerate(simple_orders)
    order_id = order_ids[i]
    node_id = string(order.zone)
    time_period = extract_time_period(order.date_time, order_book.periods)
    
    # Ποσότητα με πρόσημο: αρνητική για πώληση, θετική για αγορά
    quantity = order.type == :supply ? -order.quantity : order.quantity
    
    stepwise_dual_rhs[order_id] = @expression(model,
        stepwise_dual[order_id] +
        quantity * market_price[node_id, time_period] -
        quantity * order.price
    )
end

# Περιορισμός dual
@constraint(model, [order_id in order_ids],
    0 <= stepwise_dual_rhs[order_id])

# Big-M μετασχηματισμός για συμπληρωματικότητα
@variable(model, complementarity_aux[order_id in order_ids], Bin)

@constraint(model, [order_id in order_ids],
    stepwise_acceptance[order_id] <= complementarity_aux[order_id] * big_m)

@constraint(model, [order_id in order_ids],
    stepwise_dual_rhs[order_id] <= (1 - complementarity_aux[order_id]) * big_m)
\end{minted}
\end{samepage}

\subsection{Περιορισμός Ισορροπίας Ισχύος}
\label{subsec:mpcc-power-balance}

Ο περιορισμός ισορροπίας ισχύος διαφοροποιείται ανάλογα με το αν υπάρχουν περιορισμοί δικτύου:

\begin{samepage}
\begin{minted}[fontsize=\small,linenos,frame=lines]{julia}
if has_network_constraints
    # Πολυζωνική ισορροπία με ροές μεταφοράς
    @constraint(model, [node_id in nodes, period in periods],
        # Συνεισφορά εντολών: πώληση (αρνητική) + αγορά (θετική)
        sum(stepwise_acceptance[order_ids[i]] * 
            (orders[i].type == :supply ? -orders[i].quantity : orders[i].quantity)
            for i in orders_at_node_period[(node_id, period)]) +
        # Εισροές από άλλες ζώνες
        sum(flow[(z, node_id), period] for z in connected_zones_to[node_id]) -
        # Εκροές προς άλλες ζώνες
        sum(flow[(node_id, z), period] for z in connected_zones_from[node_id]) +
        # Έκτακτη αποκοπή
        load_shed[node_id, period] == 0
    )
else
    # Μονοζωνική ισορροπία (χωρίς ροές)
    @constraint(model, [node_id in nodes, period in periods],
        sum(stepwise_acceptance[order_ids[i]] * 
            (orders[i].type == :supply ? -orders[i].quantity : orders[i].quantity)
            for i in orders_at_node_period[(node_id, period)]) +
        load_shed[node_id, period] == 0
    )
end
\end{minted}
\end{samepage}

\subsection{Αντικειμενική Συνάρτηση}
\label{subsec:mpcc-objective}

Η αντικειμενική συνάρτηση μεγιστοποιεί το οικονομικό πλεόνασμα, με ποινή για την αποκοπή φορτίου:

\begin{samepage}
\begin{minted}[fontsize=\small,linenos,frame=lines]{julia}
load_shed_penalty = 10000.0  # €/MWh

@objective(model, Max,
    # Πλεόνασμα από εντολές
    sum(stepwise_acceptance[order_ids[i]] *
        (order.type == :supply ? 
            -order.quantity * order.price :  # Κόστος παραγωγής
             order.quantity * order.price)   # Αξία κατανάλωσης
        for (i, order) in enumerate(simple_orders)) -
    # Ποινή αποκοπής φορτίου
    sum(load_shed_penalty * load_shed[n, p] 
        for n in order_book.nodes, p in order_book.periods)
)
\end{minted}
\end{samepage}

\subsection{Δομή Αποτελέσματος}
\label{subsec:mpcc-result}

Τα αποτελέσματα της εκκαθάρισης επιστρέφονται στη δομή \texttt{MPCCResult}:

\begin{samepage}
\begin{minted}[fontsize=\small,linenos,frame=lines]{julia}
struct MPCCResult
    status::Symbol                                # :optimal, :infeasible, κλπ
    objective_value::Float64                      # Τιμή αντικειμενικής
    market_prices::Dict{String,Dict{String,Float64}}  # Τιμές [zone][period]
    stepwise_acceptance::Dict{String,Float64}     # Αποδοχές εντολών
    block_acceptance::Dict{String,Float64}        # Αποδοχές block
    block_activation::Dict{String,Float64}        # Ενεργοποιήσεις block
    transmission_flows::Dict{String,Dict{String,Float64}}  # Ροές
    solve_time::Float64                           # Χρόνος επίλυσης
    total_time::Float64                           # Συνολικός χρόνος
    solver_name::String                           # Όνομα επιλυτή
    message::String                               # Μήνυμα
end
\end{minted}
\end{samepage}


% =============================================================================
\section{Μοντελοποίηση Δικτύου}
\label{sec:network-modeling}

Η μοντελοποίηση των περιορισμών δικτύου υλοποιείται στο module \texttt{Network.jl}. Το module παρέχει δύο προσεγγίσεις: τη βασισμένη σε γραμμές μεταφοράς (\texttt{NetworkTopology}) και τη βασισμένη σε ζώνες προσφοράς (\texttt{TransferCapacity}). Η δεύτερη προσέγγιση είναι η συνιστώμενη για χρήση με δεδομένα ENTSO-E.

\subsection{Δομή Διαθέσιμης Ικανότητας Μεταφοράς}
\label{subsec:transfer-capacity-struct}

Η δομή \texttt{TransferCapacity} αναπαριστά την ικανότητα μεταφοράς μεταξύ ζωνών προσφοράς:

\begin{samepage}
\begin{minted}[fontsize=\small,linenos,frame=lines]{julia}
struct TransferCapacity
    bidding_zones::Vector{String}    # Όλες οι ζώνες
    time_periods::Vector{String}     # Χρονικές περίοδοι
    # Ικανότητα στην κατεύθυνση source → sink
    capacity_forward::Dict{Tuple{String,String,String},Float64}
    # Ικανότητα στην κατεύθυνση sink → source
    capacity_backward::Dict{Tuple{String,String,String},Float64}
end
\end{minted}
\end{samepage}

\subsection{Ανάκτηση Δεδομένων ENTSO-E}
\label{subsec:entsoe-data-retrieval}

Η δημιουργία της δομής \texttt{TransferCapacity} από δεδομένα ENTSO-E γίνεται με την ακόλουθη συνάρτηση:

\begin{samepage}
\begin{minted}[fontsize=\small,linenos,frame=lines]{julia}
function create_transfer_capacity_from_entsoe(date::Date, 
                                              bidding_zones::Vector{String}=String[])
    zone_filter = isempty(bidding_zones) ? "" :
        "AND (out_map_code IN ('$(join(bidding_zones, "','"))') " *
        "OR in_map_code IN ('$(join(bidding_zones, "','"))'))"
    
    query = """
    SELECT
        out_map_code as source_zone,
        in_map_code as sink_zone,
        EXTRACT(HOUR FROM date_time_utc) + 1 as time_period,
        capacity_mw as capacity
    FROM entsoe.offered_transfer_capacities_implicit
    WHERE DATE(date_time_utc) = \$1
    $zone_filter
    ORDER BY out_map_code, in_map_code, date_time_utc
    """
    
    df = sql2df(query, [date])
    
    if nrow(df) == 0
        error("No transfer capacity data found for $date")
    end
    
    return build_transfer_capacity_from_dataframe(df)
end
\end{minted}
\end{samepage}

\subsection{Προσθήκη Περιορισμών στο Μοντέλο}
\label{subsec:adding-network-constraints}

Οι περιορισμοί ικανότητας μεταφοράς προστίθενται στο μοντέλο JuMP ως εξής:

\begin{samepage}
\begin{minted}[fontsize=\small,linenos,frame=lines]{julia}
function add_transfer_capacity_constraints!(model::Model, 
                                            transfer_capacity::TransferCapacity, 
                                            FLOW)
    zones = transfer_capacity.bidding_zones
    periods = transfer_capacity.time_periods
    cap_fwd = transfer_capacity.capacity_forward
    cap_bwd = transfer_capacity.capacity_backward
    
    # Περιορισμοί για κάθε ζεύγος ζωνών και περίοδο
    # Θετική ροή: source → sink, περιορίζεται από capacity_forward
    # Αρνητική ροή: sink → source, περιορίζεται από capacity_backward
    @constraint(model, 
        [source in zones, sink in zones, t in periods; source != sink],
        -get(cap_bwd, (source, sink, t), 0.0) <= 
            FLOW[source, sink, t] <= 
            get(cap_fwd, (source, sink, t), 0.0)
    )
    
    return model
end
\end{minted}
\end{samepage}

\subsection{Βοηθητικές Συναρτήσεις}
\label{subsec:network-utilities}

Το module παρέχει βοηθητικές συναρτήσεις για την εξαγωγή πληροφοριών συνδεσιμότητας:

\begin{samepage}
\begin{minted}[fontsize=\small,linenos,frame=lines]{julia}
# Επιστροφή ζευγών ζωνών με διαθέσιμη ικανότητα μεταφοράς
function get_zone_pairs(tc::TransferCapacity)
    pairs = Set{Tuple{String,String}}()
    for (source, sink, _) in keys(tc.capacity_forward)
        if get(tc.capacity_forward, (source, sink, "1"), 0.0) > 0
            push!(pairs, (source, sink))
        end
    end
    return collect(pairs)
end

# Επιστροφή συνδεδεμένων ζωνών για μια δεδομένη ζώνη
function get_connected_zones(tc::TransferCapacity, zone::String)
    connected = Set{String}()
    for (source, sink, _) in keys(tc.capacity_forward)
        if source == zone
            push!(connected, sink)
        elseif sink == zone
            push!(connected, source)
        end
    end
    return collect(connected)
end
\end{minted}
\end{samepage}


% =============================================================================
\section{Ολοκλήρωση και Κύρια Ροή Εργασιών}
\label{sec:integration}

Η ολοκλήρωση όλων των συνιστωσών επιτυγχάνεται μέσω δύο κύριων συναρτήσεων: η \texttt{generate\_energy\_prices} για μονοζωνική εκκαθάριση και η \texttt{run\_multi\_zone\_market\_clearing} για πολυζωνική εκκαθάριση με περιορισμούς διασυνοριακής μεταφοράς.

\subsection{Μονοζωνική Εκκαθάριση}
\label{subsec:single-zone-clearing}

Η συνάρτηση \texttt{generate\_energy\_prices} υλοποιεί την πλήρη ροή εργασιών για μια ζώνη προσφοράς:

\begin{samepage}
\begin{minted}[fontsize=\small,linenos,frame=lines]{julia}
function generate_energy_prices(bidding_zone::String, date::Date;
                                order_method::Symbol=:uc_based,
                                markup_factor::Float64=1.1,
                                save_to_db::Bool=true)
    
    # Βήμα 1: Δημιουργία βιβλίου εντολών
    if order_method == :uc_based
        # Επίλυση UC και μετατροπή σε εντολές
        uc_result = solve_unit_commitment(bidding_zone, date)
        bids_result = generate_market_orders_from_uc(
            bidding_zone, date;
            markup_factor=markup_factor,
            bidding_strategy=:committed_only
        )
        order_book = create_typed_order_book(bids_result, bidding_zone, date)
    elseif order_method == :alternative
        # Εναλλακτική μέθοδος (ταχύτερη)
        result = create_adjusted_order_book(bidding_zone, date)
        order_book = result.order_book
    end
    
    # Βήμα 2: Εκκαθάριση αγοράς
    mpcc_result = solve_mpcc_market_clearing(order_book)
    
    # Βήμα 3: Αποθήκευση αποτελεσμάτων
    if save_to_db && mpcc_result.status == :optimal
        save_energy_prices(mpcc_result, bidding_zone, date)
    end
    
    return mpcc_result
end
\end{minted}
\end{samepage}

\subsection{Πολυζωνική Εκκαθάριση}
\label{subsec:multi-zone-clearing}

Η συνάρτηση \texttt{run\_multi\_zone\_market\_clearing} εκτελεί ταυτόχρονη εκκαθάριση πολλαπλών ζωνών με περιορισμούς ATC:

\begin{samepage}
\begin{minted}[fontsize=\small,linenos,frame=lines]{julia}
function run_multi_zone_market_clearing(date::Date;
                                        zones::Vector{String}=String[],
                                        order_method::Symbol=:alternative,
                                        save_to_db::Bool=true)
    
    # Ανακάλυψη διαθέσιμων ζωνών αν δεν καθορίστηκαν
    if isempty(zones)
        zones = get_zones_with_transfer_capacity(date)
    end
    
    println("🌍 Running multi-zone clearing for $(length(zones)) zones")
    
    # Δημιουργία πολυζωνικού βιβλίου εντολών
    order_book = create_multi_zone_order_book(date, zones; 
                                              order_method=order_method)
    
    # Εκκαθάριση αγοράς
    result = solve_mpcc_market_clearing(order_book)
    
    # Αποθήκευση αποτελεσμάτων
    if save_to_db && result.status == :optimal
        for zone in zones
            save_energy_prices(result, zone, date)
        end
        save_transmission_flows(result, date)
    end
    
    # Εκτύπωση αποτελεσμάτων
    println("💰 Zonal Clearing Prices (average):")
    for zone in zones
        prices = values(result.market_prices[zone])
        avg = sum(prices) / length(prices)
        println("   $zone: €$(round(avg, digits=2))/MWh")
    end
    
    return result
end
\end{minted}
\end{samepage}

\paragraph{Παράλληλη Εκτέλεση UC}
Κατά την πολυζωνική εκκαθάριση με τη μέθοδο \texttt{:uc\_based}, η επίλυση Unit Commitment εκτελείται ανεξάρτητα για κάθε ζώνη. Η παράμετρος \texttt{parallel=true} ενεργοποιεί την ταυτόχρονη εκτέλεση αυτών των ανεξάρτητων υπολογισμών μέσω της συνάρτησης \texttt{pmap} του module \texttt{Distributed.jl}. Κάθε ζώνη ανατίθεται σε έναν ξεχωριστό εργάτη (worker process), ο οποίος εκτελεί πλήρως τον κύκλο ανάκτησης δεδομένων, επίλυσης UC και μετατροπής σε εντολές αγοράς. Η μέθοδος \texttt{pmap} εξασφαλίζει δυναμική κατανομή εργασιών: μόλις ένας εργάτης ολοκληρώσει μια ζώνη, αναλαμβάνει την επόμενη διαθέσιμη.

Σε περίπτωση που δεν υπάρχουν διαθέσιμοι εργάτες (\texttt{nworkers() == 0}), το σύστημα υποβαθμίζεται αυτόματα σε σειριακή εκτέλεση χωρίς σφάλμα. Η ασφάλεια νημάτων (thread safety) εξασφαλίζεται από το γεγονός ότι κάθε εργάτης λειτουργεί σε ξεχωριστό address space και τα κλειδιά κρυφής μνήμης είναι ειδικά ανά ζώνη, ενώ οι εγγραφές στη βάση δεδομένων χρησιμοποιούν συναλλαγές PostgreSQL. Η εκκαθάριση MPCC εκτελείται πάντα σειριακά μετά την ολοκλήρωση όλων των UC, καθώς απαιτεί το ενοποιημένο βιβλίο εντολών όλων των ζωνών.

\subsection{Αποθήκευση Αποτελεσμάτων}
\label{subsec:saving-results}

Τα αποτελέσματα αποθηκεύονται στη βάση δεδομένων μέσω των συναρτήσεων \texttt{save\_energy\_prices} και \texttt{save\_transmission\_flows}:

\begin{samepage}
\begin{minted}[fontsize=\small,linenos,frame=lines]{julia}
function save_energy_prices(result::MPCCResult, bidding_zone::String, date::Date)
    # Εξασφάλιση ύπαρξης πίνακα
    ensure_energy_prices_table()
    
    # Προετοιμασία δεδομένων
    prices = result.market_prices[bidding_zone]
    
    withdb() do conn
        for (period, price) in prices
            hour = parse(Int, period)
            execute(conn, """
                INSERT INTO simulations.energy_prices 
                    (bidding_zone, date, hour, price, created_at)
                VALUES (\$1, \$2, \$3, \$4, NOW())
                ON CONFLICT (bidding_zone, date, hour) 
                DO UPDATE SET price = EXCLUDED.price, created_at = NOW()
            """, [bidding_zone, date, hour, price])
        end
    end
end
\end{minted}
\end{samepage}

\subsection{Επεξεργασία Κατά Παρτίδες}
\label{subsec:batch-processing}

Για την επεξεργασία πολλαπλών ημερομηνιών, παρέχονται οι συναρτήσεις \texttt{generate\_energy\_prices\_for\_date\_range} και \texttt{run\_multi\_zone\_for\_date\_range}:

\begin{samepage}
\begin{minted}[fontsize=\small,linenos,frame=lines]{julia}
function run_multi_zone_for_date_range(start_date::Date, end_date::Date;
                                       zones::Vector{String}=String[],
                                       order_method::Symbol=:alternative,
                                       force_rerun::Bool=false,
                                       parallel::Bool=false)
    results = Dict{Date,MPCCResult}()

    for date in start_date:Day(1):end_date
        try
            result = run_multi_zone_market_clearing(date;
                zones=zones, order_method=order_method,
                force_rerun=force_rerun, parallel=parallel)
            results[date] = result
        catch e
            @warn "Failed to process $date: $e"
        end
    end
    return results
end
\end{minted}
\end{samepage}

\subsection{Παράδειγμα Χρήσης}
\label{subsec:usage-example}

Ένα πλήρες παράδειγμα χρήσης του συστήματος για τον υπολογισμό τιμών στην Ελλάδα:

\begin{samepage}
\begin{minted}[fontsize=\small,linenos,frame=lines]{julia}
using Euphemia
using Dates

# Μονοζωνική εκκαθάριση για την Ελλάδα
prices = generate_energy_prices("GR", Date(2024, 6, 15);
    order_method = :uc_based,
    markup_factor = 1.1,
    save_to_db = true
)

# Πολυζωνική εκκαθάριση Ελλάδας-Βουλγαρίας-Ιταλίας
result = run_multi_zone_market_clearing(Date(2024, 6, 15);
    zones = ["GR", "BG", "IT"],
    order_method = :alternative,
    save_to_db = true
)

# Εκτύπωση διασυνοριακών ροών
for (flow_id, flows) in result.transmission_flows
    avg_flow = sum(values(flows)) / length(flows)
    println("$flow_id: $(round(avg_flow, digits=1)) MW average")
end
\end{minted}
\end{samepage}


% =============================================================================
\section{Αυτοματοποιημένη Εκτέλεση}
\label{sec:automated-execution}

Για τη συστηματική παραγωγή τιμών σε ημερήσια βάση, αναπτύχθηκε υποδομή αυτοματοποίησης βασισμένη σε GitHub Actions. Δύο ροές εργασιών (workflows) εκτελούνται σε αυτο-φιλοξενούμενο μηχάνημα (self-hosted runner):

\begin{itemize}
\item \texttt{generate-energy-prices.yml}: Μονοζωνική εκκαθάριση, προγραμματισμένη στις 02:00 UTC ημερησίως.
\item \texttt{generate-multi-zone-prices.yml}: Πολυζωνική εκκαθάριση με διασυνοριακούς περιορισμούς, προγραμματισμένη στις 03:00 UTC ημερησίως.
\end{itemize}

Η πολυζωνική ροή εργασιών εκτελεί τα εξής βήματα: εγκατάσταση εξαρτήσεων Julia, ρύθμιση σύνδεσης βάσης δεδομένων, εκκίνηση παράλληλων εργατών (workers), εκτέλεση του σεναρίου \texttt{bin/multi\_zone\_main.jl} με τις κατάλληλες παραμέτρους, και αποστολή ειδοποιήσεων μέσω Slack σε περίπτωση επιτυχίας ή αποτυχίας. Κατά τη χρήση του εμπορικού επιλυτή Gurobi, ο αριθμός παράλληλων εργατών περιορίζεται αυτόματα σε 2, λόγω του ορίου ταυτόχρονων συνεδριών της άδειας WLS (Web License Service). Κάθε ροή εργασιών υποστηρίζει τόσο χειροκίνητη ενεργοποίηση (\texttt{workflow\_dispatch}) με παραμετροποιήσιμο εύρος ημερομηνιών, όσο και αυτόματη εκτέλεση μέσω χρονοπρογραμματισμού (\texttt{cron}).

% =============================================================================

\chapter{Υλοποίηση}
\label{ch:implementation}

Στο παρόν κεφάλαιο παρουσιάζεται η υλοποίηση του συστήματος προσομοίωσης της Προημερήσιας Αγοράς ηλεκτρικής ενέργειας. Η ανάλυση ακολουθεί τη ροή επεξεργασίας δεδομένων από τα μοντέλα οντοτήτων, μέσω του επιλυτή Unit Commitment και της μηχανής στρατηγικής προσφορών, έως την εκκαθάριση της αγοράς με τη μέθοδο MPCC. Για κάθε συνιστώσα παρέχονται αντιπροσωπευτικά αποσπάσματα κώδικα που επιδεικνύουν τις βασικές αλγοριθμικές επιλογές. Ο πλήρης κώδικας είναι διαθέσιμος στο αποθετήριο του έργου\footnote{\url{https://github.com/silentech-inc/energy-markets}}.


% =============================================================================
\section{Μοντελοποίηση Οντοτήτων}
\label{sec:entity-modeling}

Η υλοποίηση βασίζεται σε σαφώς ορισμένες δομές δεδομένων για την αναπαράσταση των οντοτήτων του συστήματος ηλεκτρικής ενέργειας. Οι δομές αυτές ορίζονται ως αμετάβλητα \texttt{struct} στη Julia, εξασφαλίζοντας ακεραιότητα δεδομένων και αποτελεσματικότητα.

\subsection{Μονάδα Παραγωγής (Generator)}
\label{subsec:generator-struct}

Η δομή \texttt{Generator} αναπαριστά μια θερμική μονάδα παραγωγής με όλα τα τεχνικά και οικονομικά χαρακτηριστικά που απαιτούνται για το πρόβλημα Unit Commitment:

\begin{samepage}
\begin{minted}[fontsize=\small,linenos,frame=lines]{julia}
struct Generator
    code::String           # Μοναδικός κωδικός (EIC code)
    name::String           # Όνομα μονάδας
    fuel_type::Symbol      # Τύπος καυσίμου (:lignite, :gas, :hydro, κλπ)
    location::String       # Τοποθεσία
    p_max::Float64         # Μέγιστη ισχύς (MW)
    p_min::Float64         # Ελάχιστη ισχύς (MW)
    bidding_zone::String   # Ζώνη προσφοράς (π.χ. "GR")
    marginal_cost::Float64 # Οριακό κόστος (€/MWh)
    # Προαιρετικές παράμετροι από εξαγωγή ιστορικών δεδομένων
    # nothing = χρήση προεπιλεγμένων τιμών ανά τύπο καυσίμου
    ramp_up::Union{Float64, Nothing}    # Ρυθμός αύξησης (κλάσμα p_max/ώρα)
    ramp_down::Union{Float64, Nothing}  # Ρυθμός μείωσης (κλάσμα p_max/ώρα)
    min_uptime::Union{Int, Nothing}     # Ελάχιστος χρόνος λειτουργίας (ώρες)
    min_downtime::Union{Int, Nothing}   # Ελάχιστος χρόνος αδράνειας (ώρες)

    # Κατασκευαστής με προεπιλεγμένες τιμές nothing
    function Generator(code, name, fuel_type, location, p_max, p_min,
                       bidding_zone, marginal_cost,
                       ramp_up=nothing, ramp_down=nothing,
                       min_uptime=nothing, min_downtime=nothing)
        new(code, name, fuel_type, location, p_max, p_min,
            bidding_zone, marginal_cost,
            ramp_up, ramp_down, min_uptime, min_downtime)
    end
end
\end{minted}
\end{samepage}

Ο τύπος καυσίμου (\texttt{fuel\_type}) αποθηκεύεται ως \texttt{Symbol} για αποτελεσματική σύγκριση και χρησιμοποιείται για την ανάκτηση παραμέτρων ειδικών ανά τύπο καυσίμου. Τα πεδία \texttt{ramp\_up}, \texttt{ramp\_down}, \texttt{min\_uptime} και \texttt{min\_downtime} ορίζονται ως \texttt{Union\{T, Nothing\}}, γεγονός που επιτρέπει τη χρήση της τιμής \texttt{nothing} ως σημαίας (sentinel) για την ενεργοποίηση των προεπιλεγμένων τιμών ανά τύπο καυσίμου. Ο εσωτερικός κατασκευαστής (inner constructor) ορίζει αυτά τα πεδία με προεπιλεγμένη τιμή \texttt{nothing}, διατηρώντας τη συμβατότητα με τον υπάρχοντα κώδικα. Οι τεχνικές παράμετροι εξάγονται αυτόματα από ιστορικά δεδομένα παραγωγής, όπως περιγράφεται στην Ενότητα~\ref{subsec:parameter-inference}.

Η ανάκτηση μονάδων παραγωγής από τη βάση δεδομένων πραγματοποιείται με τη συνάρτηση \texttt{get\_generators}, η οποία υποστηρίζει φιλτράρισμα μη διαθεσιμοτήτων:

\begin{samepage}
\begin{minted}[fontsize=\small,linenos,frame=lines]{julia}
function get_generators(bidding_zone::String, day::Date; 
                        exclude_unavailable::Bool=true)
    # Ερώτημα SQL με LEFT JOIN στον πίνακα unavailability
    # για αποκλεισμό ή μείωση ισχύος μονάδων με διακοπές
    
    query = """
    WITH active_outages AS (
        SELECT asset_code, MIN(available_capacity_mw) as available_capacity_mw
        FROM entsoe.unavailability_of_production_and_generation_units
        WHERE status = 'Active'
          AND start_outage_utc::timestamp <= \$2::date + interval '1 day'
          AND end_outage_utc::timestamp >= \$2::date
        GROUP BY asset_code
    )
    SELECT g.*, COALESCE(o.available_capacity_mw, g.generation_unit_installed_capacity_mw) as p_max
    FROM entsoe.production_and_generation_units g
    LEFT JOIN active_outages o ON g.generation_unit_code = o.asset_code
    WHERE g.area_map_code = \$1
      AND g.production_unit_status = 'COMMISSIONED'
      AND (o.asset_code IS NULL OR o.available_capacity_mw > 0)
    """
    # ... επεξεργασία αποτελεσμάτων
end
\end{minted}
\end{samepage}

\subsection{Παράμετροι Τύπου Καυσίμου}
\label{subsec:fuel-type-parameters}

Οι τεχνικοί περιορισμοί των μονάδων παραγωγής διαφέρουν σημαντικά ανάλογα με τον τύπο καυσίμου. Η δομή \texttt{FuelTypeParameters} συγκεντρώνει αυτές τις παραμέτρους:

\begin{samepage}
\begin{minted}[fontsize=\small,linenos,frame=lines]{julia}
struct FuelTypeParameters
    # Χαρακτηριστικά εκκίνησης (σε ΩΡΕΣ)
    cold_startup_time::Int          # Ώρες για ψυχρή εκκίνηση
    warm_startup_time::Int          # Ώρες για θερμή εκκίνηση
    hot_startup_time::Int           # Ώρες για καυτή εκκίνηση
    # Ελάχιστοι χρόνοι λειτουργίας (σε ΩΡΕΣ)
    min_uptime::Int                 # Ελάχιστες ώρες λειτουργίας
    min_downtime::Int               # Ελάχιστες ώρες αδράνειας
    # Ρυθμοί μεταβολής (κλάσμα δυναμικότητας ανά ΩΡΑ)
    ramp_up_rate::Float64           # Ρυθμός αύξησης (κλάσμα/ώρα)
    ramp_down_rate::Float64         # Ρυθμός μείωσης (κλάσμα/ώρα)
    # Λειτουργική ευελιξία
    min_load_factor::Float64        # Ελάχιστο φορτίο (κλάσμα)
    part_load_efficiency::Float64   # Απόδοση στο ελάχιστο φορτίο
    # Κατώφλια θερμοκρασιακής μετάβασης (ΩΡΕΣ εκτός)
    warm_threshold::Int             # Ώρες αδράνειας → θερμή κατάσταση
    cold_threshold::Int             # Ώρες αδράνειας → ψυχρή κατάσταση
    # Οικονομικές παράμετροι
    startup_cost_multiplier::Float64  # Πολλαπλασιαστής κόστους εκκίνησης
    no_load_cost_fraction::Float64    # Κόστος χωρίς φορτίο (κλάσμα)
end
\end{minted}
\end{samepage}

Σημαντικό σχεδιαστικό χαρακτηριστικό είναι ότι όλες οι χρονικές παράμετροι αποθηκεύονται σε \textbf{ώρες} (όχι σε περιόδους). Ο επιλυτής UC μετατρέπει αυτόματα σε περιόδους βάσει της πραγματικής χρονικής ανάλυσης των δεδομένων ($\text{periods\_per\_hour} = 60 / \text{resolution\_minutes}$). Τα πεδία \texttt{warm\_threshold} και \texttt{cold\_threshold} χρησιμοποιούνται τόσο στον προσδιορισμό θερμοκρασιακής κατάστασης εκκίνησης (Ενότητα~\ref{subsec:initial-conditions}) όσο και στους αντίστοιχους περιορισμούς του μοντέλου UC. Τα πεδία \texttt{startup\_cost\_multiplier} και \texttt{no\_load\_cost\_fraction} τροφοδοτούν τον υπολογισμό κόστους εκκίνησης και κόστους χωρίς φορτίο στην αντικειμενική συνάρτηση (Ενότητα~\ref{subsec:uc-balance-objective}).

Η συνάρτηση \texttt{get\_fuel\_type\_parameters} επιστρέφει τις κατάλληλες παραμέτρους για κάθε τύπο καυσίμου. Για παράδειγμα, οι λιγνιτικές μονάδες έχουν μεγαλύτερους χρόνους εκκίνησης και χαμηλότερους ρυθμούς μεταβολής σε σύγκριση με τις μονάδες συνδυασμένου κύκλου φυσικού αερίου.

\subsection{Φορτίο και Ανανεώσιμες Πηγές}
\label{subsec:load-renewables}

Η ζήτηση ηλεκτρικής ενέργειας αναπαρίσταται με τη δομή \texttt{Load}:

\begin{samepage}
\begin{minted}[fontsize=\small,linenos,frame=lines]{julia}
struct Load
    timeslot::String           # Χρονοθυρίδα (π.χ. "20240615-0000")
    resolution_code::String    # Χρονική ανάλυση ("PT15M", "PT60M")
    bidding_zone::String       # Ζώνη προσφοράς
    value::Float64             # Κατανάλωση (MW)
end
\end{minted}
\end{samepage}

Οι προβλέψεις παραγωγής από ΑΠΕ αναπαρίστανται με τη δομή \texttt{RenewablesGenerationForecast}:

\begin{samepage}
\begin{minted}[fontsize=\small,linenos,frame=lines]{julia}
struct RenewablesGenerationForecast
    date_time::String                    # Χρονοσφραγίδα
    resolution_code::String              # Χρονική ανάλυση
    bidding_zone::String                 # Ζώνη προσφοράς
    production_type::String              # "Solar" ή "Wind Onshore/Offshore"
    aggregated_generation_forecast::Float64  # Πρόβλεψη (MW)
end
\end{minted}
\end{samepage}

\subsection{Εξαγωγή Τεχνικών Παραμέτρων από Ιστορικά Δεδομένα}
\label{subsec:parameter-inference}

Οι τεχνικές παράμετροι των μονάδων παραγωγής (ρυθμοί μεταβολής, ελάχιστη ισχύς, ελάχιστοι χρόνοι λειτουργίας/αδράνειας) δεν είναι διαθέσιμες στα δημόσια δεδομένα του ENTSO-E. Για την αντιμετώπιση αυτού του περιορισμού, αναπτύχθηκε ένας μηχανισμός αυτόματης εξαγωγής (inference) αυτών των παραμέτρων από ιστορικά δεδομένα πραγματικής παραγωγής.

\paragraph{Πηγή Δεδομένων}
Τα ιστορικά δεδομένα παραγωγής ανακτώνται από τον πίνακα \texttt{actual\_generation\_output\_per\_generation\_unit} του ENTSO-E, ο οποίος περιέχει τη μετρούμενη παραγωγή κάθε μονάδας ανά 15 λεπτά ή ανά ώρα. Η συνάρτηση \texttt{get\_historical\_generation} ανακτά δεδομένα για παράθυρο 3 μηνών πριν από την επιλεγμένη ημερομηνία.

\paragraph{Εξαγωγή Ρυθμών Μεταβολής}
Οι ρυθμοί μεταβολής (ramp rates) υπολογίζονται από τη συνάρτηση \texttt{infer\_ramp\_rates}:

\begin{samepage}
\begin{minted}[fontsize=\small,linenos,frame=lines]{julia}
function infer_ramp_rates(historical_data::DataFrame, p_max::Float64;
                          percentile::Float64=0.95)
    if nrow(historical_data) < MIN_DATA_POINTS_FOR_RAMP_INFERENCE
        return (ramp_up=nothing, ramp_down=nothing)
    end

    # Ταξινόμηση κατά χρόνο και υπολογισμός μεταβολών
    sort!(historical_data, :date_time_utc)
    gen_values = historical_data.actual_generation_output_mw
    deltas = diff(gen_values)

    # Διαχωρισμός σε αυξήσεις και μειώσεις
    ramp_ups = filter(x -> x > 0, deltas)
    ramp_downs = map(abs, filter(x -> x < 0, deltas))

    # Κανονικοποίηση σε ωριαίο ρυθμό βάσει ανάλυσης δεδομένων
    resolution_code = historical_data.resolution_code[1]
    resolution_minutes = parse_resolution_to_minutes(resolution_code)
    hourly_factor = 60.0 / resolution_minutes

    # 95ο εκατοστημόριο, μετατροπή σε κλάσμα p_max/ώρα
    ramp_up = isempty(ramp_ups) ? nothing :
        (quantile(ramp_ups, percentile) * hourly_factor) / p_max
    ramp_down = isempty(ramp_downs) ? nothing :
        (quantile(ramp_downs, percentile) * hourly_factor) / p_max

    return (ramp_up=ramp_up, ramp_down=ramp_down)
end
\end{minted}
\end{samepage}

Οι ρυθμοί αποθηκεύονται ως κλάσμα της μέγιστης ισχύος ανά ώρα (π.χ. 0.20 σημαίνει ότι η μονάδα μπορεί να μεταβάλει την παραγωγή της κατά 20\% της $P^{\max}$ ανά ώρα), επιτρέποντας την εύκολη μετατροπή σε οποιαδήποτε χρονική ανάλυση.

\paragraph{Εξαγωγή Ελάχιστης Ισχύος}
Η ελάχιστη ισχύς ($P^{\min}$) εξάγεται αναλύοντας τις περιόδους σταθερής λειτουργίας, εξαιρώντας τις μεταβατικές καταστάσεις (εκκίνηση/τερματισμός):

\begin{samepage}
\begin{minted}[fontsize=\small,linenos,frame=lines]{julia}
function infer_p_min(historical_data::DataFrame, p_max::Float64,
                     fuel_type::Symbol; percentile::Float64=0.05)
    if nrow(historical_data) < MIN_DATA_POINTS_FOR_RAMP_INFERENCE
        return nothing
    end

    sort!(historical_data, :date_time_utc)
    gen_values = historical_data.actual_generation_output_mw
    deltas = diff(gen_values)

    # Ανίχνευση μεταβατικών καταστάσεων (|Δ| > 5% της p_max)
    ramp_threshold = 0.05 * p_max

    # Συλλογή σταθερών μη-μηδενικών τιμών
    stable_values = Float64[]
    for i in 2:(length(gen_values) - 1)
        value = gen_values[i]
        if value > 0 && abs(deltas[i-1]) < ramp_threshold &&
                        abs(deltas[i]) < ramp_threshold
            push!(stable_values, value)
        end
    end

    if length(stable_values) < 50
        return nothing  # Ανεπαρκή δεδομένα
    end

    # 5ο εκατοστημόριο της σταθερής λειτουργίας
    inferred_p_min = quantile(stable_values, percentile)

    # Εφαρμογή ορίων ανά τύπο καυσίμου (clamp σε MW)
    (min_frac, max_frac) = P_MIN_BOUNDS[get_p_min_bounds_category(fuel_type)]
    return clamp(inferred_p_min, min_frac * p_max, max_frac * p_max)
end
\end{minted}
\end{samepage}

\paragraph{Όρια ανά Τύπο Καυσίμου}
Για την αποφυγή μη ρεαλιστικών τιμών, εφαρμόζονται όρια (bounds) ανά τύπο καυσίμου. Τα όρια αυτά βασίζονται στα τεχνικά χαρακτηριστικά κάθε τεχνολογίας παραγωγής:

\begin{table}[htbp]
\centering
\caption{Όρια ελάχιστης ισχύος ($P^{\min}$) ανά τύπο καυσίμου}
\label{tab:p-min-bounds}
\begin{tabular}{@{}lcc@{}}
\toprule
\textbf{Τύπος Καυσίμου} & \textbf{Ελάχιστο (\%)} & \textbf{Μέγιστο (\%)} \\
\midrule
Λιγνίτης/Άνθρακας & 45 & 65 \\
Φυσικό αέριο (CCGT) & 35 & 55 \\
Φυσικό αέριο (OCGT) & 20 & 45 \\
Υδροηλεκτρικά/Μπαταρίες & 0 & 0 \\
Προεπιλογή & 30 & 60 \\
\bottomrule
\end{tabular}
\end{table}

Οι ευέλικτοι πόροι (υδροηλεκτρικά, μπαταρίες) λαμβάνουν $P^{\min} = 0$ χωρίς εξαγωγή, καθώς μπορούν να λειτουργήσουν σε οποιοδήποτε επίπεδο ισχύος. Η αντίστοιχη σταθερά \texttt{FLEXIBLE\_FUEL\_TYPES} περιλαμβάνει τους τύπους \texttt{Hydro Pumped Storage}, \texttt{Hydro Run-of-river}, \texttt{Hydro Water Reservoir}, \texttt{Battery} και \texttt{Other}. Αξίζει να σημειωθεί ότι οι προεπιλεγμένες τιμές \texttt{min\_load\_factor} στη δομή \texttt{FuelTypeParameters} ευθυγραμμίστηκαν με τα όρια εξαγωγής (π.χ. λιγνίτης: 0.45, ευέλικτοι πόροι: 0.0).

\paragraph{Εξαγωγή Ελάχιστων Χρόνων Λειτουργίας/Αδράνειας}
Οι ελάχιστοι χρόνοι υπολογίζονται αναλύοντας τους κύκλους λειτουργίας (on/off) της μονάδας:

\begin{samepage}
\begin{minted}[fontsize=\small,linenos,frame=lines]{julia}
function infer_uptime_downtime(historical_data::DataFrame, fuel_type::Symbol;
                               percentile::Float64=0.05)
    if nrow(historical_data) < MIN_DATA_POINTS_FOR_RAMP_INFERENCE
        return (nothing, nothing)
    end

    sort!(historical_data, :date_time_utc)
    gen_values = historical_data.actual_generation_output_mw

    # Υπολογισμός διάρκειας περιόδου σε ώρες
    resolution_minutes = parse_resolution_to_minutes(
        historical_data.resolution_code[1])
    period_hours = resolution_minutes / 60.0

    # Προσδιορισμός κατάστασης (on: > 1 MW, off: = 0)
    is_on = gen_values .> 1.0

    # Εντοπισμός κύκλων λειτουργίας
    on_durations, off_durations = Float64[], Float64[]
    current_state = is_on[1]
    current_duration = 1

    for i in 2:length(is_on)
        if is_on[i] == current_state
            current_duration += 1
        else
            duration_hours = current_duration * period_hours
            if current_state
                push!(on_durations, duration_hours)
            else
                push!(off_durations, duration_hours)
            end
            current_state = is_on[i]
            current_duration = 1
        end
    end

    # Φιλτράρισμα σύντομων glitches (< 2 περιόδων)
    min_dur = 2 * period_hours
    filter!(d -> d >= min_dur, on_durations)
    filter!(d -> d >= min_dur, off_durations)

    # Απαίτηση τουλάχιστον 5 κύκλων για αξιόπιστη εξαγωγή
    # Εφαρμογή ορίων ανά τύπο καυσίμου (βλ. Πίνακα~\ref{tab:uptime-downtime-bounds})
    bounds = UPTIME_DOWNTIME_BOUNDS[get_p_min_bounds_category(fuel_type)]
    min_uptime = length(on_durations) >= 5 ?
        round(Int, clamp(quantile(on_durations, percentile),
                         bounds[1]...)) : nothing
    min_downtime = length(off_durations) >= 5 ?
        round(Int, clamp(quantile(off_durations, percentile),
                         bounds[2]...)) : nothing

    return (min_uptime, min_downtime)
end
\end{minted}
\end{samepage}

Ομοίως με τα όρια $P^{\min}$, εφαρμόζονται όρια ανά τύπο καυσίμου για τους χρόνους λειτουργίας και αδράνειας, όπως παρουσιάζονται στον Πίνακα~\ref{tab:uptime-downtime-bounds}.

\begin{table}[htbp]
\centering
\caption{Όρια ελάχιστων χρόνων λειτουργίας/αδράνειας ανά τύπο καυσίμου (ώρες)}
\label{tab:uptime-downtime-bounds}
\begin{tabular}{@{}lcccc@{}}
\toprule
\textbf{Τύπος Καυσίμου} & \multicolumn{2}{c}{\textbf{Χρόνος Λειτουργίας}} & \multicolumn{2}{c}{\textbf{Χρόνος Αδράνειας}} \\
\cmidrule(lr){2-3} \cmidrule(lr){4-5}
& \textbf{Ελάχ.} & \textbf{Μέγ.} & \textbf{Ελάχ.} & \textbf{Μέγ.} \\
\midrule
Λιγνίτης/Άνθρακας & 8 & 48 & 4 & 24 \\
Φυσικό αέριο (CCGT) & 2 & 12 & 1 & 8 \\
Φυσικό αέριο (OCGT) & 1 & 4 & 1 & 4 \\
Προεπιλογή & 2 & 24 & 1 & 12 \\
\bottomrule
\end{tabular}
\end{table}

Η εφαρμογή ορίων διασφαλίζει ότι οι εξαχθείσες τιμές παραμένουν εντός τεχνικά ρεαλιστικού εύρους. Για μονάδες βάσης (π.χ. λιγνιτικές), η εξαγωγή χρόνων συχνά δεν είναι εφικτή λόγω σπάνιων κύκλων λειτουργίας/αδράνειας, οπότε χρησιμοποιούνται οι προεπιλεγμένες τιμές τύπου καυσίμου.

\paragraph{Ενοποίηση Εξαγωγής Παραμέτρων}
Η κεντρική συνάρτηση \texttt{infer\_parameters\_for\_generators} ενορχηστρώνει την εξαγωγή για ένα σύνολο μονάδων, εφαρμόζοντας διαφορετική λογική ανάλογα με τον τύπο καυσίμου:

\begin{samepage}
\begin{minted}[fontsize=\small,linenos,frame=lines]{julia}
function infer_parameters_for_generators(generators::Vector{Generator},
                                         day::Date)
    updated = Generator[]
    for gen in generators
        historical = get_historical_generation(gen.code, day)
        rates = infer_ramp_rates(historical, gen.p_max)

        if gen.fuel_type in FLEXIBLE_FUEL_TYPES
            # Ευέλικτοι πόροι: p_min = 0, χωρίς uptime/downtime
            push!(updated, Generator(gen.code, gen.name, gen.fuel_type,
                gen.location, gen.p_max, 0.0, gen.bidding_zone,
                gen.marginal_cost, rates.ramp_up, rates.ramp_down,
                nothing, nothing))
        else
            # Θερμικές μονάδες: εξαγωγή όλων των παραμέτρων
            p_min = infer_p_min(historical, gen.p_max, gen.fuel_type)
            (uptime, downtime) = infer_uptime_downtime(historical,
                                                        gen.fuel_type)
            push!(updated, Generator(gen.code, gen.name, gen.fuel_type,
                gen.location, gen.p_max,
                p_min !== nothing ? p_min : gen.p_min,
                gen.bidding_zone, gen.marginal_cost,
                rates.ramp_up, rates.ramp_down, uptime, downtime))
        end
    end
    return updated
end
\end{minted}
\end{samepage}

\paragraph{Αποθήκευση σε Κρυφή Μνήμη (Caching)}
Η εξαγωγή παραμέτρων από ιστορικά δεδομένα είναι υπολογιστικά ακριβή (περίπου 17 λεπτά για 39 μονάδες). Για τη βελτίωση της απόδοσης, οι εξαχθείσες παράμετροι αποθηκεύονται στον πίνακα \texttt{simulations.generator\_inferred\_parameters}:

\begin{samepage}
\begin{minted}[fontsize=\small,linenos,frame=lines]{julia}
# Ανάκτηση μονάδων με αυτόματη εξαγωγή και caching
generators = get_generators_with_inferred_params("GR", Date(2024, 6, 15);
    use_cache = true,       # Χρήση cached τιμών αν διαθέσιμες
    max_age_days = 7        # Ανανέωση cache κάθε 7 ημέρες
)
\end{minted}
\end{samepage}

Η χρήση της κρυφής μνήμης μειώνει τον χρόνο ανάκτησης από περίπου 17 λεπτά σε περίπου 2 δευτερόλεπτα---βελτίωση κατά 525 φορές.


\subsection{Αρχικές Συνθήκες Μονάδων Παραγωγής}
\label{subsec:initial-conditions}

Για τη ρεαλιστική μοντελοποίηση του προβλήματος Unit Commitment, η κατάσταση κάθε μονάδας κατά τη χρονική στιγμή $t=0$ (έναρξη του ορίζοντα βελτιστοποίησης) πρέπει να είναι γνωστή. Οι αρχικές συνθήκες καθορίζουν αν η μονάδα λειτουργούσε, σε ποιο επίπεδο ισχύος, και για πόσο χρόνο βρισκόταν στην τρέχουσα κατάστασή της. Αυτή η πληροφορία είναι απαραίτητη για τη σωστή επιβολή περιορισμών στις πρώτες περιόδους του ορίζοντα.

\paragraph{Δομή Αρχικών Συνθηκών}
Η δομή \texttt{InitialConditions} αναπαριστά την κατάσταση μιας μονάδας στο $t=0$:

\begin{samepage}
\begin{minted}[fontsize=\small,linenos,frame=lines]{julia}
struct InitialConditions
    is_on::Bool           # Λειτουργούσε η μονάδα στο t=0;
    output::Float64       # Ισχύς εξόδου σε MW (0 αν εκτός λειτουργίας)
    hours_on::Int         # Συνεχόμενες ώρες λειτουργίας
    hours_off::Int        # Συνεχόμενες ώρες αδράνειας
    thermal_state::Symbol # Θερμοκρασιακή κατάσταση (:hot, :warm, :cold)
end
\end{minted}
\end{samepage}

Η θερμοκρασιακή κατάσταση προσδιορίζεται από τον χρόνο αδράνειας: \texttt{:hot} αν $T_{\text{off}} \leq 8$ ώρες, \texttt{:warm} αν $8 < T_{\text{off}} \leq 48$ ώρες, και \texttt{:cold} αν $T_{\text{off}} > 48$ ώρες. Τα κατώφλια αυτά προσαρμόζονται ανά τύπο καυσίμου μέσω των αντίστοιχων πεδίων \texttt{warm\_threshold} και \texttt{cold\_threshold} της δομής \texttt{FuelTypeParameters}.

\paragraph{Εξαγωγή από Ιστορικά Δεδομένα}
Η κεντρική συνάρτηση \texttt{get\_initial\_conditions} εξάγει τις αρχικές συνθήκες για ένα σύνολο μονάδων χρησιμοποιώντας ένα μαζικό ερώτημα SQL (batch query) αντί $N$ ξεχωριστών ερωτημάτων. Η τεχνική αυτή επιτυγχάνει βελτίωση ταχύτητας κατά 36 φορές (από $\sim$9 λεπτά σε $\sim$15 δευτερόλεπτα για 40 μονάδες):

\begin{samepage}
\begin{minted}[fontsize=\small,linenos,frame=lines]{julia}
function get_initial_conditions(generators::Vector{Generator},
                                market_day::Date;
                                use_historical::Bool=true)
    # Μαζική ανάκτηση 72 ωρών ιστορικού για όλες τις μονάδες
    # σε ένα ερώτημα SQL (WHERE generation_unit_code = ANY($1))
    day_before = market_day - Day(1)
    start_utc = DateTime(day_before, Time(23, 0, 0))  # ~00:00 CET
    codes = [gen.code for gen in generators]
    all_historical = get_recent_generation_batch(codes, start_utc)

    conditions = Dict{String, InitialConditions}()
    for gen in generators
        gen_data = filter(r -> r.generation_unit_code == gen.code,
                          all_historical)
        ic = infer_initial_conditions_from_data(gen_data, gen.fuel_type)
        conditions[gen.code] = ic
    end
    return conditions
end
\end{minted}
\end{samepage}

Ο αλγόριθμος εξαγωγής αναλύει τα δεδομένα σε φθίνουσα χρονική σειρά: η πιο πρόσφατη μέτρηση $> 1$ MW δηλώνει ότι η μονάδα ήταν σε λειτουργία. Στη συνέχεια, μετρούνται οι συνεχόμενες περίοδοι στην ίδια κατάσταση για τον υπολογισμό του $T_{\text{on}}$ ή $T_{\text{off}}$. Σε περίπτωση απουσίας δεδομένων, χρησιμοποιούνται ευρετικές προεπιλεγμένες τιμές ανά τύπο καυσίμου: οι μονάδες βάσης (λιγνίτης, πυρηνικά) θεωρούνται σε λειτουργία, οι μονάδες μέσου φορτίου (CCGT) εκτός λειτουργίας αλλά θερμές, οι μονάδες αιχμής (OCGT) εκτός λειτουργίας και ψυχρές, και οι ευέλικτοι πόροι (υδροηλεκτρικά, μπαταρίες) εκτός λειτουργίας αλλά έτοιμοι.


\subsection{Τύποι Εντολών Αγοράς}
\label{subsec:market-orders}

Οι εντολές αγοράς υλοποιούνται ως ιεραρχία τύπων με αφηρημένο βασικό τύπο \texttt{MarketOrder}. Ο κύριος τύπος που χρησιμοποιείται είναι η απλή ωριαία εντολή \texttt{SimpleOrder}:

\begin{samepage}
\begin{minted}[fontsize=\small,linenos,frame=lines]{julia}
abstract type MarketOrder end

struct SimpleOrder <: MarketOrder
    type::Symbol           # :supply ή :demand
    price::Float64         # Τιμή (€/MWh)
    quantity::Float64      # Ποσότητα (MWh)
    zone::Symbol           # Ζώνη προσφοράς
    date_time::DateTime    # Χρόνος παράδοσης
    resolution_code::Int   # Ανάλυση σε λεπτά (60, 30, 15)
end
\end{minted}
\end{samepage}

Πέρα από τις απλές εντολές, υλοποιούνται και σύνθετοι τύποι όπως οι εντολές μπλοκ (\texttt{BlockOrder}), οι συνδεδεμένες εντολές μπλοκ (\texttt{LinkedBlockOrder}) και οι αποκλειστικές εντολές μπλοκ (\texttt{ExclusiveBlockOrder}), ακολουθώντας την προδιαγραφή του αλγορίθμου EUPHEMIA.


% =============================================================================
\section{Επιλυτής Unit Commitment}
\label{sec:uc-solver}

Ο επιλυτής Unit Commitment υλοποιείται στο module \texttt{UnitCommitment.jl} και αποτελεί το πρώτο βήμα της αλυσίδας υπολογισμών. Η κύρια συνάρτηση \texttt{solve\_unit\_commitment} δέχεται ως είσοδο τη ζώνη προσφοράς και την ημερομηνία, ανακτά τα απαραίτητα δεδομένα από τη βάση, διαμορφώνει το πρόβλημα βελτιστοποίησης και επιστρέφει τη βέλτιστη λύση.

\subsection{Δημιουργία Μοντέλου}
\label{subsec:uc-model-creation}

Η δημιουργία του μοντέλου JuMP ξεκινά με την επιλογή του επιλυτή και τον ορισμό των μεταβλητών απόφασης:

\begin{samepage}
\begin{minted}[fontsize=\small,linenos,frame=lines]{julia}
function solve_unit_commitment(bidding_zone::String, day::Date;
                               optimizer::String="auto",
                               mip_gap::Float64=0.01,
                               time_limit::Float64=600.0)
    # Φόρτωση δεδομένων
    generators = get_generators(bidding_zone, day)
    loads = get_loads(bidding_zone, day)
    renewables = get_generation_forecast_for_wind_and_solar(bidding_zone, day)
    initial_conditions = get_initial_conditions(generators, day)

    N = length(generators)  # Αριθμός μονάδων
    T = 24                  # Χρονικές περίοδοι
    periods_per_hour = 60 / resolution_minutes

    # Δημιουργία μοντέλου με ρύθμιση παραμέτρων επιλυτή
    model = Model(optimizer_func)
    set_silent(model)
    set_optimizer_attribute(model, MOI.RelativeGapTolerance(), mip_gap)
    set_optimizer_attribute(model, MOI.TimeLimitSec(), time_limit)

    # Εξαγωγή αρχικών συνθηκών σε πίνακες
    u0 = zeros(Int, N)      # Αρχική δέσμευση (0 ή 1)
    g0 = zeros(Float64, N)  # Αρχική παραγωγή (MW)
    T_on0 = zeros(Int, N)   # Ώρες ήδη σε λειτουργία
    T_off0 = zeros(Int, N)  # Ώρες ήδη εκτός λειτουργίας
    for (i, gen) in enumerate(generators)
        ic = initial_conditions[gen.code]
        u0[i] = ic.is_on ? 1 : 0
        g0[i] = ic.output
        T_on0[i] = ic.hours_on
        T_off0[i] = ic.hours_off
    end

    # Μεταβλητές παραγωγής (συνεχείς)
    @variable(model, 0 <= g[i=1:N, t=1:T] <= generators[i].p_max)

    # Μεταβλητές δέσμευσης (δυαδικές)
    @variable(model, u[i=1:N, t=1:T], Bin)  # Κατάσταση λειτουργίας
    @variable(model, v[i=1:N, t=1:T], Bin)  # Εκκίνηση
    @variable(model, z[i=1:N, t=1:T], Bin)  # Τερματισμός

    # Μεταβλητές για θερμοκρασιακές καταστάσεις εκκίνησης
    Θ = [:cold, :warm, :hot]
    @variable(model, v_θ[i=1:N, θ in Θ, t=1:T], Bin)

    # ... περαιτέρω μεταβλητές
end
\end{minted}
\end{samepage}

\subsection{Περιορισμοί Ισχύος}
\label{subsec:impl-uc-power-constraints}

Οι περιορισμοί ισχύος εξασφαλίζουν ότι η παραγωγή κάθε μονάδας βρίσκεται εντός των τεχνικών της ορίων:

\begin{samepage}
\begin{minted}[fontsize=\small,linenos,frame=lines]{julia}
# Ανώτατο όριο παραγωγής
@constraint(model, [i = 1:N, t = 1:T], 
    g[i, t] <= generators[i].p_max * u[i, t])

# Κατώτατο όριο παραγωγής με παραμέτρους ανά τύπο καυσίμου
for (i, gen) in enumerate(generators)
    params = fuel_params[i]
    # Το ελάχιστο είναι το μέγιστο μεταξύ p_min και του ελάχιστου 
    # φορτίου βάσει τύπου καυσίμου
    min_gen = max(gen.p_min, params.min_load_factor * gen.p_max)
    @constraint(model, [t = 1:T], g[i, t] >= min_gen * u[i, t])
end
\end{minted}
\end{samepage}

\subsection{Περιορισμοί Δέσμευσης}
\label{subsec:impl-uc-commitment-constraints}

Η σχέση μεταξύ κατάστασης λειτουργίας, εκκίνησης και τερματισμού υλοποιείται ως εξής, με ενσωμάτωση αρχικών συνθηκών στην περίοδο $t=1$:

\begin{samepage}
\begin{minted}[fontsize=\small,linenos,frame=lines]{julia}
# Σύνδεση t=1 με αρχική κατάσταση u0
if initial_conditions !== nothing
    @constraint(model, [i in 1:N],
        u[i, 1] == u0[i] + v[i, 1] - z[i, 1])
end

# Σύνδεση μεταβλητών δέσμευσης (t ≥ 2)
@constraint(model, [i in 1:N, t in 2:T],
    u[i, t] == u[i, t-1] + v[i, t] - z[i, t])

# Αποκλεισμός ταυτόχρονης εκκίνησης και τερματισμού
@constraint(model, [i = 1:N, t = 1:T], v[i, t] + z[i, t] <= 1)

# Ελάχιστος χρόνος λειτουργίας (ώρες → περίοδοι)
for (i, gen) in enumerate(generators)
    UT_hours = gen.min_uptime !== nothing ?
        gen.min_uptime : fuel_params[i].min_uptime
    UT = max(1, ceil(Int, UT_hours * periods_per_hour))
    if UT > 1
        @constraint(model, [t = UT:T],
            sum(v[i, τ] for τ in t-UT+1:t) <= u[i, t])
        # Αρχική συνθήκη: εναπομένων χρόνος λειτουργίας
        if initial_conditions !== nothing && u0[i] == 1
            remaining = max(0, UT - ceil(Int, T_on0[i] * periods_per_hour))
            for t in 1:min(remaining, T)
                @constraint(model, u[i, t] == 1)
            end
        end
    end
end

# Ελάχιστος χρόνος αδράνειας (ώρες → περίοδοι)
for (i, gen) in enumerate(generators)
    DT_hours = gen.min_downtime !== nothing ?
        gen.min_downtime : fuel_params[i].min_downtime
    DT = max(1, ceil(Int, DT_hours * periods_per_hour))
    if DT > 1
        @constraint(model, [t = DT:T],
            sum(z[i, τ] for τ in t-DT+1:t) <= 1 - u[i, t])
        # Αρχική συνθήκη: εναπομένων χρόνος αδράνειας
        if initial_conditions !== nothing && u0[i] == 0
            remaining = max(0, DT - ceil(Int, T_off0[i] * periods_per_hour))
            for t in 1:min(remaining, T)
                @constraint(model, u[i, t] == 0)
            end
        end
    end
end
\end{minted}
\end{samepage}

Η μετατροπή $\texttt{UT\_hours} \times \texttt{periods\_per\_hour}$ διασφαλίζει ότι οι περιορισμοί εφαρμόζονται σωστά ανεξάρτητα από τη χρονική ανάλυση (15, 30 ή 60 λεπτά). Οι αρχικές συνθήκες επιβάλλουν τη συνέχιση της τρέχουσας κατάστασης κάθε μονάδας: αν μια μονάδα λειτουργούσε αλλά δεν έχει συμπληρώσει τον ελάχιστο χρόνο λειτουργίας, εξαναγκάζεται να παραμείνει σε λειτουργία για τις εναπομένουσες περιόδους.

\subsection{Περιορισμοί Ρυθμού Μεταβολής}
\label{subsec:impl-uc-ramping}

Οι περιορισμοί ρυθμού μεταβολής (ramping) χρησιμοποιούν την τεχνική Big-M για την αποδέσμευση κατά τις μεταβατικές καταστάσεις:

\begin{samepage}
\begin{minted}[fontsize=\small,linenos,frame=lines]{julia}
# Υπολογισμός ρυθμών μεταβολής ανά μονάδα:
# εξαχθείσα τιμή (fraction/hour) ή προεπιλογή τύπου καυσίμου
# Μετατροπή σε MW/περίοδο: rate × p_max × period_hours
period_hours = resolution_minutes / 60.0
ramp_up = Float64[]
ramp_down = Float64[]
for (i, gen) in enumerate(generators)
    if gen.ramp_up !== nothing
        push!(ramp_up, gen.ramp_up * gen.p_max * period_hours)
    else
        push!(ramp_up, fuel_params[i].ramp_up_rate * gen.p_max * period_hours)
    end
    if gen.ramp_down !== nothing
        push!(ramp_down, gen.ramp_down * gen.p_max * period_hours)
    else
        push!(ramp_down, fuel_params[i].ramp_down_rate * gen.p_max * period_hours)
    end
end

# Big-M: κλιμακούμενο βάσει μεγέθους προβλήματος
max_p_max = maximum(gen.p_max for gen in generators)
M = 2.0 * max_p_max

# Περιορισμός αύξησης (t ≥ 2)
@constraint(model, [i in 1:N, t in 2:T],
    g[i, t] - g[i, t-1] <= ramp_up[i] + M * u_SU[i, t])

# Περιορισμός μείωσης (t ≥ 2)
@constraint(model, [i in 1:N, t in 2:T],
    g[i, t-1] - g[i, t] <= ramp_down[i] + M * u_SD[i, t])
\end{minted}
\end{samepage}

\subsection{Ισορροπία Ισχύος και Αντικειμενική Συνάρτηση}
\label{subsec:uc-balance-objective}

Ο περιορισμός ισορροπίας ισχύος εξασφαλίζει την κάλυψη της καθαρής ζήτησης:

\begin{samepage}
\begin{minted}[fontsize=\small,linenos,frame=lines]{julia}
# Περιορισμός ισορροπίας
@constraint(model, [t in 1:T],
    sum(g[i, t] for i in 1:N) == net_demand[t])
\end{minted}
\end{samepage}

Η αντικειμενική συνάρτηση αποτελείται από τρεις συνιστώσες κόστους: κόστος παραγωγής, κόστος εκκίνησης και κόστος χωρίς φορτίο. Πριν τη διαμόρφωση, υπολογίζονται οι παράμετροι κόστους:

\begin{samepage}
\begin{minted}[fontsize=\small,linenos,frame=lines]{julia}
# Κόστος εκκίνησης ανά μονάδα και θερμοκρασιακή κατάσταση (€)
# Βασικό κόστος = startup_cost_multiplier × marginal_cost × p_max
# Πολλαπλασιαστές: hot=1.0, warm=1.5, cold=2.5
C_SU = Dict{Tuple{Int,Symbol},Float64}()
for i in 1:N
    params = fuel_params[i]
    gen = generators[i]
    base = params.startup_cost_multiplier * gen.marginal_cost * gen.p_max
    C_SU[(i, :hot)]  = base * 1.0
    C_SU[(i, :warm)] = base * 1.5
    C_SU[(i, :cold)] = base * 2.5
end

# Κόστος χωρίς φορτίο ανά μονάδα (€/περίοδο)
# Πάγιο κόστος δέσμευσης ανεξάρτητο από παραγωγή
period_hours = resolution_minutes / 60.0
C_NL = Float64[]
for i in 1:N
    params = fuel_params[i]
    gen = generators[i]
    push!(C_NL, params.no_load_cost_fraction *
                gen.marginal_cost * gen.p_min * period_hours)
end

# Αντικειμενική: ελαχιστοποίηση συνολικού κόστους
@objective(model, Min,
    # 1. Κόστος παραγωγής (μεταβλητό)
    sum(generators[i].marginal_cost * g[i, t] for i in 1:N, t in 1:T)
    # 2. Κόστος εκκίνησης (εξαρτώμενο από θερμοκρασία)
    + sum(C_SU[(i, θ)] * v_θ[i, θ, t] for i in 1:N, θ in Θ, t in 1:T)
    # 3. Κόστος χωρίς φορτίο (πάγιο κόστος δέσμευσης)
    + sum(C_NL[i] * u[i, t] for i in 1:N, t in 1:T)
)
\end{minted}
\end{samepage}

Η τρισδιάστατη δομή κόστους αποτυπώνει ρεαλιστικά τα κίνητρα λειτουργίας: το κόστος εκκίνησης αποθαρρύνει τις συχνές εκκινήσεις/τερματισμούς (ιδιαίτερα τις ψυχρές εκκινήσεις, που κοστίζουν 2,5 φορές περισσότερο από τις θερμές), ενώ το κόστος χωρίς φορτίο αντανακλά τα πάγια λειτουργικά κόστη δέσμευσης μιας μονάδας ανεξάρτητα από το επίπεδο παραγωγής της.

\subsection{Επίλυση και Επιστροφή Αποτελεσμάτων}
\label{subsec:uc-solution}

Μετά την επίλυση, τα αποτελέσματα επιστρέφονται σε δομημένη μορφή:

\begin{samepage}
\begin{minted}[fontsize=\small,linenos,frame=lines]{julia}
optimize!(model)

if termination_status(model) == OPTIMAL
    return (
        status = OPTIMAL,
        generators = generators,
        generation = value.(g),       # Πίνακας παραγωγής N×T
        commitment = value.(u),       # Πίνακας δέσμευσης N×T
        startup = value.(v),          # Πίνακας εκκινήσεων N×T
        shutdown = value.(z),         # Πίνακας τερματισμών N×T
        objective_value = objective_value(model),
        solve_time = solve_time
    )
else
    return (status = termination_status(model),)
end
\end{minted}
\end{samepage}

\subsection{Αποθήκευση Αποτελεσμάτων σε Κρυφή Μνήμη}
\label{subsec:uc-caching}

Η επίλυση του προβλήματος Unit Commitment απαιτεί 10--15 λεπτά ανά ζώνη. Για την αποφυγή επαναληπτικών υπολογισμών, τα αποτελέσματα αποθηκεύονται σε τρεις πίνακες PostgreSQL (\texttt{uc\_results}, \texttt{uc\_generation}, \texttt{uc\_net\_demand}) με κλειδί \texttt{(bidding\_zone, market\_date, code\_version)}. Η συνάρτηση \texttt{solve\_unit\_commitment} ελέγχει αυτόματα τη διαθεσιμότητα αποθηκευμένων αποτελεσμάτων:

\begin{samepage}
\begin{minted}[fontsize=\small,linenos,frame=lines]{julia}
function solve_unit_commitment(bidding_zone::String, day::Date;
                               optimizer::String="auto",
                               use_cache::Bool=true,
                               force_rerun::Bool=false,
                               # ... λοιπές παράμετροι
                               )
    # Έλεγχος κρυφής μνήμης (εκτός αν απενεργοποιημένη)
    if use_cache && !force_rerun
        cached = load_uc_results(bidding_zone, day)
        if cached !== nothing
            return cached   # Επιστροφή σε ~2 δευτερόλεπτα
        end
    end

    # ... κανονική επίλυση (10-15 λεπτά) ...

    # Αποθήκευση αποτελεσμάτων (delete-before-insert)
    if use_cache
        save_uc_results(solution, bidding_zone, day)
    end
    return solution
end
\end{minted}
\end{samepage}

Η παράμετρος \texttt{force\_rerun} διαδίδεται μέσω ολόκληρης της αλυσίδας κλήσεων (\texttt{generate\_energy\_prices} $\to$ \texttt{create\_typed\_order\_book} $\to$ \texttt{generate\_market\_orders\_from\_uc} $\to$ \texttt{solve\_unit\_commitment}), επιτρέποντας τον επανυπολογισμό αποτελεσμάτων όταν αλλάξει ο κώδικας ή τα δεδομένα εισόδου. Η αποθήκευση ανέρχεται σε περίπου 195 KB ανά ζώνη/ημέρα---κόστος αμελητέο σε σύγκριση με την εξοικονόμηση χρόνου.


% =============================================================================
\section{Μηχανή Στρατηγικής Προσφορών}
\label{sec:bidding-strategy-engine}

Η μηχανή στρατηγικής προσφορών υλοποιείται στο module \texttt{BiddingStrategy.jl} και μετατρέπει τα αποτελέσματα του Unit Commitment σε εντολές αγοράς. Η κύρια συνάρτηση \texttt{generate\_market\_orders\_from\_uc} υποστηρίζει τρεις διακριτές στρατηγικές.

\subsection{Δομή Αποτελέσματος}
\label{subsec:bidding-result}

Τα αποτελέσματα της μετατροπής επιστρέφονται στη δομή \texttt{UCToBidsResult}:

\begin{samepage}
\begin{minted}[fontsize=\small,linenos,frame=lines]{julia}
struct UCToBidsResult
    supply_orders::Vector{SimpleOrder}   # Εντολές πώλησης
    demand_orders::Vector{SimpleOrder}   # Εντολές αγοράς
    bidding_zone::Symbol                 # Ζώνη προσφοράς
    day::Date                            # Ημερομηνία
    total_supply_quantity::Float64       # Συνολική προσφερόμενη ποσότητα
    total_demand_quantity::Float64       # Συνολική ζητούμενη ποσότητα
    success::Bool                        # Επιτυχία μετατροπής
    message::String                      # Μήνυμα κατάστασης
end
\end{minted}
\end{samepage}

\subsection{Στρατηγική Δεσμευμένων Μονάδων}
\label{subsec:committed-only-strategy}

Η συντηρητική στρατηγική \texttt{:committed\_only} δημιουργεί προσφορές μόνο για τις μονάδες που έχουν δεσμευτεί στη λύση του UC:

\begin{samepage}
\begin{minted}[fontsize=\small,linenos,frame=lines]{julia}
function apply_committed_only_strategy!(orders, uc_solution, generators, 
                                        day, markup_factor)
    for (i, gen) in enumerate(generators)
        for t in 1:24
            # Έλεγχος αν η μονάδα είναι δεσμευμένη
            if uc_solution.commitment[i, t] > COMMITMENT_THRESHOLD
                generation = uc_solution.generation[i, t]
                if generation > GENERATION_THRESHOLD
                    # Δημιουργία εντολής πώλησης
                    price = gen.marginal_cost * markup_factor
                    dt = DateTime(day, Time(t-1, 0))
                    push!(orders, SimpleOrder(
                        :supply, price, generation, 
                        Symbol(gen.bidding_zone), dt, 60
                    ))
                end
            end
        end
    end
end
\end{minted}
\end{samepage}

\subsection{Στρατηγική Μέγιστης Διαθεσιμότητας}
\label{subsec:all-available-max}

Η στρατηγική \texttt{:all\_available\_max} προσφέρει όλες τις μονάδες στη μέγιστη δυναμικότητά τους, ανεξάρτητα από τη λύση του UC. Αν και μη ρεαλιστική, χρησιμεύει ως βάση σύγκρισης:

\begin{samepage}
\begin{minted}[fontsize=\small,linenos,frame=lines]{julia}
function apply_all_available_max_strategy!(orders, generators, day, markup_factor)
    for gen in generators
        for t in 1:24
            price = gen.marginal_cost * markup_factor
            dt = DateTime(day, Time(t-1, 0))
            push!(orders, SimpleOrder(
                :supply, price, gen.p_max,
                Symbol(gen.bidding_zone), dt, 60
            ))
        end
    end
end
\end{minted}
\end{samepage}

\subsection{Στρατηγική Ρεαλιστικής Διαθεσιμότητας}
\label{subsec:all-available-realistic}

Η στρατηγική \texttt{:all\_available\_realistic} συνδυάζει πληροφορίες από τη λύση UC με τη διαθεσιμότητα όλων των μονάδων:

\begin{samepage}
\begin{minted}[fontsize=\small,linenos,frame=lines]{julia}
function apply_all_available_realistic_strategy!(orders, uc_solution, 
                                                  generators, day, markup_factor)
    for (i, gen) in enumerate(generators)
        for t in 1:24
            # Καθορισμός ποσότητας βάσει κατάστασης δέσμευσης
            if uc_solution.commitment[i, t] > COMMITMENT_THRESHOLD
                # Δεσμευμένη μονάδα: χρήση πραγματικής παραγωγής
                quantity = uc_solution.generation[i, t]
            elseif gen.p_min < 0.1 * gen.p_max
                # Μη δεσμευμένη με χαμηλό ελάχιστο: 20% της δυναμικότητας
                quantity = DEFAULT_UNCOMMITTED_UNIT_FRACTION * gen.p_max
            else
                # Μη δεσμευμένη: ελάχιστη ισχύς
                quantity = gen.p_min
            end
            
            if quantity > GENERATION_THRESHOLD
                price = gen.marginal_cost * markup_factor
                dt = DateTime(day, Time(t-1, 0))
                push!(orders, SimpleOrder(
                    :supply, price, quantity,
                    Symbol(gen.bidding_zone), dt, 60
                ))
            end
        end
    end
end
\end{minted}
\end{samepage}

\subsection{Δημιουργία Εντολών Ζήτησης}
\label{subsec:demand-orders-creation}

Οι εντολές ζήτησης δημιουργούνται με βάση τα δεδομένα φορτίου, με τιμή ίση με το ανώτατο όριο της αγοράς (price-taking demand):

\begin{samepage}
\begin{minted}[fontsize=\small,linenos,frame=lines]{julia}
function create_demand_orders!(orders, loads, day, bidding_zone, price_cap)
    for load in loads
        if load.value > DEMAND_THRESHOLD
            hour = parse(Int, load.timeslot[end-3:end-2])
            dt = DateTime(day, Time(hour, 0))
            push!(orders, SimpleOrder(
                :demand, price_cap, load.value,
                Symbol(bidding_zone), dt, 60
            ))
        end
    end
end
\end{minted}
\end{samepage}


% =============================================================================
\section{Εκκαθάριση Αγοράς με MPCC}
\label{sec:mpcc-market-clearing}

Η εκκαθάριση της αγοράς υλοποιείται στο module \texttt{MPCC.jl} χρησιμοποιώντας τη μέθοδο Mathematical Programming with Complementarity Constraints. Ο αλγόριθμος δέχεται ένα βιβλίο εντολών (\texttt{MPCCOrderBook}) και επιστρέφει τις τιμές εκκαθάρισης και τους βαθμούς αποδοχής των εντολών.

\subsection{Δομή Βιβλίου Εντολών}
\label{subsec:order-book-structure}

Το βιβλίο εντολών οργανώνει τις εντολές αγοράς μαζί με τις παραμέτρους του προβλήματος:

\begin{samepage}
\begin{minted}[fontsize=\small,linenos,frame=lines]{julia}
struct MPCCOrderBook
    orders::Vector{MarketOrder}              # Όλες οι εντολές
    nodes::Vector{String}                    # Ζώνες προσφοράς
    periods::Vector{String}                  # Χρονικές περίοδοι
    price_limits::Tuple{Float64,Float64}     # (min, max) τιμή
    network_topology::Union{Nothing,TransferCapacity}  # Περιορισμοί δικτύου
end
\end{minted}
\end{samepage}

\subsection{Μεταβλητές Απόφασης}
\label{subsec:mpcc-variables}

Οι κύριες μεταβλητές του μοντέλου MPCC είναι οι τιμές αγοράς, οι βαθμοί αποδοχής και οι ροές μεταφοράς:

\begin{samepage}
\begin{minted}[fontsize=\small,linenos,frame=lines]{julia}
function solve_mpcc_market_clearing(order_book::MPCCOrderBook; 
                                    big_m::Float64=BIG_M_PARAMETER)
    model = Model(HiGHS.Optimizer)
    set_silent(model)
    
    # Τιμές αγοράς ανά ζώνη και περίοδο
    @variable(model, market_price[node in order_book.nodes, 
                                  period in order_book.periods])
    
    # Βαθμός αποδοχής κάθε εντολής [0, 1]
    @variable(model, 0 <= stepwise_acceptance[order_id in order_ids] <= 1)
    
    # Dual μεταβλητή για περιορισμούς συμπληρωματικότητας
    @variable(model, stepwise_dual[order_id in order_ids] >= 0)
    
    # Μεταβλητή έκτακτης αποκοπής φορτίου
    @variable(model, load_shed[node in order_book.nodes, 
                               period in order_book.periods] >= 0)
    
    # Ροές μεταφοράς (αν υπάρχει δίκτυο)
    if order_book.network_topology !== nothing
        zone_pairs = get_zone_pairs(order_book.network_topology)
        @variable(model, flow[(s,d) in zone_pairs, t in order_book.periods])
    end
end
\end{minted}
\end{samepage}

\subsection{Περιορισμοί Συμπληρωματικότητας}
\label{subsec:complementarity-constraints}

Οι περιορισμοί συμπληρωματικότητας εξασφαλίζουν ότι μια εντολή γίνεται αποδεκτή μόνο όταν η τιμή αγοράς είναι ευνοϊκή. Υλοποιούνται με τη μέθοδο Big-M:

\begin{samepage}
\begin{minted}[fontsize=\small,linenos,frame=lines]{julia}
# Υπολογισμός του δεξιού μέλους της dual συνθήκης
for (i, order) in enumerate(simple_orders)
    order_id = order_ids[i]
    node_id = string(order.zone)
    time_period = extract_time_period(order.date_time, order_book.periods)
    
    # Ποσότητα με πρόσημο: αρνητική για πώληση, θετική για αγορά
    quantity = order.type == :supply ? -order.quantity : order.quantity
    
    stepwise_dual_rhs[order_id] = @expression(model,
        stepwise_dual[order_id] +
        quantity * market_price[node_id, time_period] -
        quantity * order.price
    )
end

# Περιορισμός dual
@constraint(model, [order_id in order_ids],
    0 <= stepwise_dual_rhs[order_id])

# Big-M μετασχηματισμός για συμπληρωματικότητα
@variable(model, complementarity_aux[order_id in order_ids], Bin)

@constraint(model, [order_id in order_ids],
    stepwise_acceptance[order_id] <= complementarity_aux[order_id] * big_m)

@constraint(model, [order_id in order_ids],
    stepwise_dual_rhs[order_id] <= (1 - complementarity_aux[order_id]) * big_m)
\end{minted}
\end{samepage}

\subsection{Περιορισμός Ισορροπίας Ισχύος}
\label{subsec:mpcc-power-balance}

Ο περιορισμός ισορροπίας ισχύος διαφοροποιείται ανάλογα με το αν υπάρχουν περιορισμοί δικτύου:

\begin{samepage}
\begin{minted}[fontsize=\small,linenos,frame=lines]{julia}
if has_network_constraints
    # Πολυζωνική ισορροπία με ροές μεταφοράς
    @constraint(model, [node_id in nodes, period in periods],
        # Συνεισφορά εντολών: πώληση (αρνητική) + αγορά (θετική)
        sum(stepwise_acceptance[order_ids[i]] * 
            (orders[i].type == :supply ? -orders[i].quantity : orders[i].quantity)
            for i in orders_at_node_period[(node_id, period)]) +
        # Εισροές από άλλες ζώνες
        sum(flow[(z, node_id), period] for z in connected_zones_to[node_id]) -
        # Εκροές προς άλλες ζώνες
        sum(flow[(node_id, z), period] for z in connected_zones_from[node_id]) +
        # Έκτακτη αποκοπή
        load_shed[node_id, period] == 0
    )
else
    # Μονοζωνική ισορροπία (χωρίς ροές)
    @constraint(model, [node_id in nodes, period in periods],
        sum(stepwise_acceptance[order_ids[i]] * 
            (orders[i].type == :supply ? -orders[i].quantity : orders[i].quantity)
            for i in orders_at_node_period[(node_id, period)]) +
        load_shed[node_id, period] == 0
    )
end
\end{minted}
\end{samepage}

\subsection{Αντικειμενική Συνάρτηση}
\label{subsec:mpcc-objective}

Η αντικειμενική συνάρτηση μεγιστοποιεί το οικονομικό πλεόνασμα, με ποινή για την αποκοπή φορτίου:

\begin{samepage}
\begin{minted}[fontsize=\small,linenos,frame=lines]{julia}
load_shed_penalty = 10000.0  # €/MWh

@objective(model, Max,
    # Πλεόνασμα από εντολές
    sum(stepwise_acceptance[order_ids[i]] *
        (order.type == :supply ? 
            -order.quantity * order.price :  # Κόστος παραγωγής
             order.quantity * order.price)   # Αξία κατανάλωσης
        for (i, order) in enumerate(simple_orders)) -
    # Ποινή αποκοπής φορτίου
    sum(load_shed_penalty * load_shed[n, p] 
        for n in order_book.nodes, p in order_book.periods)
)
\end{minted}
\end{samepage}

\subsection{Δομή Αποτελέσματος}
\label{subsec:mpcc-result}

Τα αποτελέσματα της εκκαθάρισης επιστρέφονται στη δομή \texttt{MPCCResult}:

\begin{samepage}
\begin{minted}[fontsize=\small,linenos,frame=lines]{julia}
struct MPCCResult
    status::Symbol                                # :optimal, :infeasible, κλπ
    objective_value::Float64                      # Τιμή αντικειμενικής
    market_prices::Dict{String,Dict{String,Float64}}  # Τιμές [zone][period]
    stepwise_acceptance::Dict{String,Float64}     # Αποδοχές εντολών
    block_acceptance::Dict{String,Float64}        # Αποδοχές block
    block_activation::Dict{String,Float64}        # Ενεργοποιήσεις block
    transmission_flows::Dict{String,Dict{String,Float64}}  # Ροές
    solve_time::Float64                           # Χρόνος επίλυσης
    total_time::Float64                           # Συνολικός χρόνος
    solver_name::String                           # Όνομα επιλυτή
    message::String                               # Μήνυμα
end
\end{minted}
\end{samepage}


% =============================================================================
\section{Μοντελοποίηση Δικτύου}
\label{sec:network-modeling}

Η μοντελοποίηση των περιορισμών δικτύου υλοποιείται στο module \texttt{Network.jl}. Το module παρέχει δύο προσεγγίσεις: τη βασισμένη σε γραμμές μεταφοράς (\texttt{NetworkTopology}) και τη βασισμένη σε ζώνες προσφοράς (\texttt{TransferCapacity}). Η δεύτερη προσέγγιση είναι η συνιστώμενη για χρήση με δεδομένα ENTSO-E.

\subsection{Δομή Διαθέσιμης Ικανότητας Μεταφοράς}
\label{subsec:transfer-capacity-struct}

Η δομή \texttt{TransferCapacity} αναπαριστά την ικανότητα μεταφοράς μεταξύ ζωνών προσφοράς:

\begin{samepage}
\begin{minted}[fontsize=\small,linenos,frame=lines]{julia}
struct TransferCapacity
    bidding_zones::Vector{String}    # Όλες οι ζώνες
    time_periods::Vector{String}     # Χρονικές περίοδοι
    # Ικανότητα στην κατεύθυνση source → sink
    capacity_forward::Dict{Tuple{String,String,String},Float64}
    # Ικανότητα στην κατεύθυνση sink → source
    capacity_backward::Dict{Tuple{String,String,String},Float64}
end
\end{minted}
\end{samepage}

\subsection{Ανάκτηση Δεδομένων ENTSO-E}
\label{subsec:entsoe-data-retrieval}

Η δημιουργία της δομής \texttt{TransferCapacity} από δεδομένα ENTSO-E γίνεται με την ακόλουθη συνάρτηση:

\begin{samepage}
\begin{minted}[fontsize=\small,linenos,frame=lines]{julia}
function create_transfer_capacity_from_entsoe(date::Date, 
                                              bidding_zones::Vector{String}=String[])
    zone_filter = isempty(bidding_zones) ? "" :
        "AND (out_map_code IN ('$(join(bidding_zones, "','"))') " *
        "OR in_map_code IN ('$(join(bidding_zones, "','"))'))"
    
    query = """
    SELECT
        out_map_code as source_zone,
        in_map_code as sink_zone,
        EXTRACT(HOUR FROM date_time_utc) + 1 as time_period,
        capacity_mw as capacity
    FROM entsoe.offered_transfer_capacities_implicit
    WHERE DATE(date_time_utc) = \$1
    $zone_filter
    ORDER BY out_map_code, in_map_code, date_time_utc
    """
    
    df = sql2df(query, [date])
    
    if nrow(df) == 0
        error("No transfer capacity data found for $date")
    end
    
    return build_transfer_capacity_from_dataframe(df)
end
\end{minted}
\end{samepage}

\subsection{Προσθήκη Περιορισμών στο Μοντέλο}
\label{subsec:adding-network-constraints}

Οι περιορισμοί ικανότητας μεταφοράς προστίθενται στο μοντέλο JuMP ως εξής:

\begin{samepage}
\begin{minted}[fontsize=\small,linenos,frame=lines]{julia}
function add_transfer_capacity_constraints!(model::Model, 
                                            transfer_capacity::TransferCapacity, 
                                            FLOW)
    zones = transfer_capacity.bidding_zones
    periods = transfer_capacity.time_periods
    cap_fwd = transfer_capacity.capacity_forward
    cap_bwd = transfer_capacity.capacity_backward
    
    # Περιορισμοί για κάθε ζεύγος ζωνών και περίοδο
    # Θετική ροή: source → sink, περιορίζεται από capacity_forward
    # Αρνητική ροή: sink → source, περιορίζεται από capacity_backward
    @constraint(model, 
        [source in zones, sink in zones, t in periods; source != sink],
        -get(cap_bwd, (source, sink, t), 0.0) <= 
            FLOW[source, sink, t] <= 
            get(cap_fwd, (source, sink, t), 0.0)
    )
    
    return model
end
\end{minted}
\end{samepage}

\subsection{Βοηθητικές Συναρτήσεις}
\label{subsec:network-utilities}

Το module παρέχει βοηθητικές συναρτήσεις για την εξαγωγή πληροφοριών συνδεσιμότητας:

\begin{samepage}
\begin{minted}[fontsize=\small,linenos,frame=lines]{julia}
# Επιστροφή ζευγών ζωνών με διαθέσιμη ικανότητα μεταφοράς
function get_zone_pairs(tc::TransferCapacity)
    pairs = Set{Tuple{String,String}}()
    for (source, sink, _) in keys(tc.capacity_forward)
        if get(tc.capacity_forward, (source, sink, "1"), 0.0) > 0
            push!(pairs, (source, sink))
        end
    end
    return collect(pairs)
end

# Επιστροφή συνδεδεμένων ζωνών για μια δεδομένη ζώνη
function get_connected_zones(tc::TransferCapacity, zone::String)
    connected = Set{String}()
    for (source, sink, _) in keys(tc.capacity_forward)
        if source == zone
            push!(connected, sink)
        elseif sink == zone
            push!(connected, source)
        end
    end
    return collect(connected)
end
\end{minted}
\end{samepage}


% =============================================================================
\section{Ολοκλήρωση και Κύρια Ροή Εργασιών}
\label{sec:integration}

Η ολοκλήρωση όλων των συνιστωσών επιτυγχάνεται μέσω δύο κύριων συναρτήσεων: η \texttt{generate\_energy\_prices} για μονοζωνική εκκαθάριση και η \texttt{run\_multi\_zone\_market\_clearing} για πολυζωνική εκκαθάριση με περιορισμούς διασυνοριακής μεταφοράς.

\subsection{Μονοζωνική Εκκαθάριση}
\label{subsec:single-zone-clearing}

Η συνάρτηση \texttt{generate\_energy\_prices} υλοποιεί την πλήρη ροή εργασιών για μια ζώνη προσφοράς:

\begin{samepage}
\begin{minted}[fontsize=\small,linenos,frame=lines]{julia}
function generate_energy_prices(bidding_zone::String, date::Date;
                                order_method::Symbol=:uc_based,
                                markup_factor::Float64=1.1,
                                save_to_db::Bool=true)
    
    # Βήμα 1: Δημιουργία βιβλίου εντολών
    if order_method == :uc_based
        # Επίλυση UC και μετατροπή σε εντολές
        uc_result = solve_unit_commitment(bidding_zone, date)
        bids_result = generate_market_orders_from_uc(
            bidding_zone, date;
            markup_factor=markup_factor,
            bidding_strategy=:committed_only
        )
        order_book = create_typed_order_book(bids_result, bidding_zone, date)
    elseif order_method == :alternative
        # Εναλλακτική μέθοδος (ταχύτερη)
        result = create_adjusted_order_book(bidding_zone, date)
        order_book = result.order_book
    end
    
    # Βήμα 2: Εκκαθάριση αγοράς
    mpcc_result = solve_mpcc_market_clearing(order_book)
    
    # Βήμα 3: Αποθήκευση αποτελεσμάτων
    if save_to_db && mpcc_result.status == :optimal
        save_energy_prices(mpcc_result, bidding_zone, date)
    end
    
    return mpcc_result
end
\end{minted}
\end{samepage}

\subsection{Πολυζωνική Εκκαθάριση}
\label{subsec:multi-zone-clearing}

Η συνάρτηση \texttt{run\_multi\_zone\_market\_clearing} εκτελεί ταυτόχρονη εκκαθάριση πολλαπλών ζωνών με περιορισμούς ATC:

\begin{samepage}
\begin{minted}[fontsize=\small,linenos,frame=lines]{julia}
function run_multi_zone_market_clearing(date::Date;
                                        zones::Vector{String}=String[],
                                        order_method::Symbol=:alternative,
                                        save_to_db::Bool=true)
    
    # Ανακάλυψη διαθέσιμων ζωνών αν δεν καθορίστηκαν
    if isempty(zones)
        zones = get_zones_with_transfer_capacity(date)
    end
    
    println("🌍 Running multi-zone clearing for $(length(zones)) zones")
    
    # Δημιουργία πολυζωνικού βιβλίου εντολών
    order_book = create_multi_zone_order_book(date, zones; 
                                              order_method=order_method)
    
    # Εκκαθάριση αγοράς
    result = solve_mpcc_market_clearing(order_book)
    
    # Αποθήκευση αποτελεσμάτων
    if save_to_db && result.status == :optimal
        for zone in zones
            save_energy_prices(result, zone, date)
        end
        save_transmission_flows(result, date)
    end
    
    # Εκτύπωση αποτελεσμάτων
    println("💰 Zonal Clearing Prices (average):")
    for zone in zones
        prices = values(result.market_prices[zone])
        avg = sum(prices) / length(prices)
        println("   $zone: €$(round(avg, digits=2))/MWh")
    end
    
    return result
end
\end{minted}
\end{samepage}

\paragraph{Παράλληλη Εκτέλεση UC}
Κατά την πολυζωνική εκκαθάριση με τη μέθοδο \texttt{:uc\_based}, η επίλυση Unit Commitment εκτελείται ανεξάρτητα για κάθε ζώνη. Η παράμετρος \texttt{parallel=true} ενεργοποιεί την ταυτόχρονη εκτέλεση αυτών των ανεξάρτητων υπολογισμών μέσω της συνάρτησης \texttt{pmap} του module \texttt{Distributed.jl}. Κάθε ζώνη ανατίθεται σε έναν ξεχωριστό εργάτη (worker process), ο οποίος εκτελεί πλήρως τον κύκλο ανάκτησης δεδομένων, επίλυσης UC και μετατροπής σε εντολές αγοράς. Η μέθοδος \texttt{pmap} εξασφαλίζει δυναμική κατανομή εργασιών: μόλις ένας εργάτης ολοκληρώσει μια ζώνη, αναλαμβάνει την επόμενη διαθέσιμη.

Σε περίπτωση που δεν υπάρχουν διαθέσιμοι εργάτες (\texttt{nworkers() == 0}), το σύστημα υποβαθμίζεται αυτόματα σε σειριακή εκτέλεση χωρίς σφάλμα. Η ασφάλεια νημάτων (thread safety) εξασφαλίζεται από το γεγονός ότι κάθε εργάτης λειτουργεί σε ξεχωριστό address space και τα κλειδιά κρυφής μνήμης είναι ειδικά ανά ζώνη, ενώ οι εγγραφές στη βάση δεδομένων χρησιμοποιούν συναλλαγές PostgreSQL. Η εκκαθάριση MPCC εκτελείται πάντα σειριακά μετά την ολοκλήρωση όλων των UC, καθώς απαιτεί το ενοποιημένο βιβλίο εντολών όλων των ζωνών.

\subsection{Αποθήκευση Αποτελεσμάτων}
\label{subsec:saving-results}

Τα αποτελέσματα αποθηκεύονται στη βάση δεδομένων μέσω των συναρτήσεων \texttt{save\_energy\_prices} και \texttt{save\_transmission\_flows}:

\begin{samepage}
\begin{minted}[fontsize=\small,linenos,frame=lines]{julia}
function save_energy_prices(result::MPCCResult, bidding_zone::String, date::Date)
    # Εξασφάλιση ύπαρξης πίνακα
    ensure_energy_prices_table()
    
    # Προετοιμασία δεδομένων
    prices = result.market_prices[bidding_zone]
    
    withdb() do conn
        for (period, price) in prices
            hour = parse(Int, period)
            execute(conn, """
                INSERT INTO simulations.energy_prices 
                    (bidding_zone, date, hour, price, created_at)
                VALUES (\$1, \$2, \$3, \$4, NOW())
                ON CONFLICT (bidding_zone, date, hour) 
                DO UPDATE SET price = EXCLUDED.price, created_at = NOW()
            """, [bidding_zone, date, hour, price])
        end
    end
end
\end{minted}
\end{samepage}

\subsection{Επεξεργασία Κατά Παρτίδες}
\label{subsec:batch-processing}

Για την επεξεργασία πολλαπλών ημερομηνιών, παρέχονται οι συναρτήσεις \texttt{generate\_energy\_prices\_for\_date\_range} και \texttt{run\_multi\_zone\_for\_date\_range}:

\begin{samepage}
\begin{minted}[fontsize=\small,linenos,frame=lines]{julia}
function run_multi_zone_for_date_range(start_date::Date, end_date::Date;
                                       zones::Vector{String}=String[],
                                       order_method::Symbol=:alternative,
                                       force_rerun::Bool=false,
                                       parallel::Bool=false)
    results = Dict{Date,MPCCResult}()

    for date in start_date:Day(1):end_date
        try
            result = run_multi_zone_market_clearing(date;
                zones=zones, order_method=order_method,
                force_rerun=force_rerun, parallel=parallel)
            results[date] = result
        catch e
            @warn "Failed to process $date: $e"
        end
    end
    return results
end
\end{minted}
\end{samepage}

\subsection{Παράδειγμα Χρήσης}
\label{subsec:usage-example}

Ένα πλήρες παράδειγμα χρήσης του συστήματος για τον υπολογισμό τιμών στην Ελλάδα:

\begin{samepage}
\begin{minted}[fontsize=\small,linenos,frame=lines]{julia}
using Euphemia
using Dates

# Μονοζωνική εκκαθάριση για την Ελλάδα
prices = generate_energy_prices("GR", Date(2024, 6, 15);
    order_method = :uc_based,
    markup_factor = 1.1,
    save_to_db = true
)

# Πολυζωνική εκκαθάριση Ελλάδας-Βουλγαρίας-Ιταλίας
result = run_multi_zone_market_clearing(Date(2024, 6, 15);
    zones = ["GR", "BG", "IT"],
    order_method = :alternative,
    save_to_db = true
)

# Εκτύπωση διασυνοριακών ροών
for (flow_id, flows) in result.transmission_flows
    avg_flow = sum(values(flows)) / length(flows)
    println("$flow_id: $(round(avg_flow, digits=1)) MW average")
end
\end{minted}
\end{samepage}


% =============================================================================
\section{Αυτοματοποιημένη Εκτέλεση}
\label{sec:automated-execution}

Για τη συστηματική παραγωγή τιμών σε ημερήσια βάση, αναπτύχθηκε υποδομή αυτοματοποίησης βασισμένη σε GitHub Actions. Δύο ροές εργασιών (workflows) εκτελούνται σε αυτο-φιλοξενούμενο μηχάνημα (self-hosted runner):

\begin{itemize}
\item \texttt{generate-energy-prices.yml}: Μονοζωνική εκκαθάριση, προγραμματισμένη στις 02:00 UTC ημερησίως.
\item \texttt{generate-multi-zone-prices.yml}: Πολυζωνική εκκαθάριση με διασυνοριακούς περιορισμούς, προγραμματισμένη στις 03:00 UTC ημερησίως.
\end{itemize}

Η πολυζωνική ροή εργασιών εκτελεί τα εξής βήματα: εγκατάσταση εξαρτήσεων Julia, ρύθμιση σύνδεσης βάσης δεδομένων, εκκίνηση παράλληλων εργατών (workers), εκτέλεση του σεναρίου \texttt{bin/multi\_zone\_main.jl} με τις κατάλληλες παραμέτρους, και αποστολή ειδοποιήσεων μέσω Slack σε περίπτωση επιτυχίας ή αποτυχίας. Κατά τη χρήση του εμπορικού επιλυτή Gurobi, ο αριθμός παράλληλων εργατών περιορίζεται αυτόματα σε 2, λόγω του ορίου ταυτόχρονων συνεδριών της άδειας WLS (Web License Service). Κάθε ροή εργασιών υποστηρίζει τόσο χειροκίνητη ενεργοποίηση (\texttt{workflow\_dispatch}) με παραμετροποιήσιμο εύρος ημερομηνιών, όσο και αυτόματη εκτέλεση μέσω χρονοπρογραμματισμού (\texttt{cron}).

% =============================================================================

\chapter{Υλοποίηση}
\label{ch:implementation}

Στο παρόν κεφάλαιο παρουσιάζεται η υλοποίηση του συστήματος προσομοίωσης της Προημερήσιας Αγοράς ηλεκτρικής ενέργειας. Η ανάλυση ακολουθεί τη ροή επεξεργασίας δεδομένων από τα μοντέλα οντοτήτων, μέσω του επιλυτή Unit Commitment και της μηχανής στρατηγικής προσφορών, έως την εκκαθάριση της αγοράς με τη μέθοδο MPCC. Για κάθε συνιστώσα παρέχονται αντιπροσωπευτικά αποσπάσματα κώδικα που επιδεικνύουν τις βασικές αλγοριθμικές επιλογές. Ο πλήρης κώδικας είναι διαθέσιμος στο αποθετήριο του έργου\footnote{\url{https://github.com/silentech-inc/energy-markets}}.


% =============================================================================
\section{Μοντελοποίηση Οντοτήτων}
\label{sec:entity-modeling}

Η υλοποίηση βασίζεται σε σαφώς ορισμένες δομές δεδομένων για την αναπαράσταση των οντοτήτων του συστήματος ηλεκτρικής ενέργειας. Οι δομές αυτές ορίζονται ως αμετάβλητα \texttt{struct} στη Julia, εξασφαλίζοντας ακεραιότητα δεδομένων και αποτελεσματικότητα.

\subsection{Μονάδα Παραγωγής (Generator)}
\label{subsec:generator-struct}

Η δομή \texttt{Generator} αναπαριστά μια θερμική μονάδα παραγωγής με όλα τα τεχνικά και οικονομικά χαρακτηριστικά που απαιτούνται για το πρόβλημα Unit Commitment:

\begin{samepage}
\begin{minted}[fontsize=\small,linenos,frame=lines]{julia}
struct Generator
    code::String           # Μοναδικός κωδικός (EIC code)
    name::String           # Όνομα μονάδας
    fuel_type::Symbol      # Τύπος καυσίμου (:lignite, :gas, :hydro, κλπ)
    location::String       # Τοποθεσία
    p_max::Float64         # Μέγιστη ισχύς (MW)
    p_min::Float64         # Ελάχιστη ισχύς (MW)
    bidding_zone::String   # Ζώνη προσφοράς (π.χ. "GR")
    marginal_cost::Float64 # Οριακό κόστος (€/MWh)
    # Προαιρετικές παράμετροι από εξαγωγή ιστορικών δεδομένων
    # nothing = χρήση προεπιλεγμένων τιμών ανά τύπο καυσίμου
    ramp_up::Union{Float64, Nothing}    # Ρυθμός αύξησης (κλάσμα p_max/ώρα)
    ramp_down::Union{Float64, Nothing}  # Ρυθμός μείωσης (κλάσμα p_max/ώρα)
    min_uptime::Union{Int, Nothing}     # Ελάχιστος χρόνος λειτουργίας (ώρες)
    min_downtime::Union{Int, Nothing}   # Ελάχιστος χρόνος αδράνειας (ώρες)

    # Κατασκευαστής με προεπιλεγμένες τιμές nothing
    function Generator(code, name, fuel_type, location, p_max, p_min,
                       bidding_zone, marginal_cost,
                       ramp_up=nothing, ramp_down=nothing,
                       min_uptime=nothing, min_downtime=nothing)
        new(code, name, fuel_type, location, p_max, p_min,
            bidding_zone, marginal_cost,
            ramp_up, ramp_down, min_uptime, min_downtime)
    end
end
\end{minted}
\end{samepage}

Ο τύπος καυσίμου (\texttt{fuel\_type}) αποθηκεύεται ως \texttt{Symbol} για αποτελεσματική σύγκριση και χρησιμοποιείται για την ανάκτηση παραμέτρων ειδικών ανά τύπο καυσίμου. Τα πεδία \texttt{ramp\_up}, \texttt{ramp\_down}, \texttt{min\_uptime} και \texttt{min\_downtime} ορίζονται ως \texttt{Union\{T, Nothing\}}, γεγονός που επιτρέπει τη χρήση της τιμής \texttt{nothing} ως σημαίας (sentinel) για την ενεργοποίηση των προεπιλεγμένων τιμών ανά τύπο καυσίμου. Ο εσωτερικός κατασκευαστής (inner constructor) ορίζει αυτά τα πεδία με προεπιλεγμένη τιμή \texttt{nothing}, διατηρώντας τη συμβατότητα με τον υπάρχοντα κώδικα. Οι τεχνικές παράμετροι εξάγονται αυτόματα από ιστορικά δεδομένα παραγωγής, όπως περιγράφεται στην Ενότητα~\ref{subsec:parameter-inference}.

Η ανάκτηση μονάδων παραγωγής από τη βάση δεδομένων πραγματοποιείται με τη συνάρτηση \texttt{get\_generators}, η οποία υποστηρίζει φιλτράρισμα μη διαθεσιμοτήτων:

\begin{samepage}
\begin{minted}[fontsize=\small,linenos,frame=lines]{julia}
function get_generators(bidding_zone::String, day::Date; 
                        exclude_unavailable::Bool=true)
    # Ερώτημα SQL με LEFT JOIN στον πίνακα unavailability
    # για αποκλεισμό ή μείωση ισχύος μονάδων με διακοπές
    
    query = """
    WITH active_outages AS (
        SELECT asset_code, MIN(available_capacity_mw) as available_capacity_mw
        FROM entsoe.unavailability_of_production_and_generation_units
        WHERE status = 'Active'
          AND start_outage_utc::timestamp <= \$2::date + interval '1 day'
          AND end_outage_utc::timestamp >= \$2::date
        GROUP BY asset_code
    )
    SELECT g.*, COALESCE(o.available_capacity_mw, g.generation_unit_installed_capacity_mw) as p_max
    FROM entsoe.production_and_generation_units g
    LEFT JOIN active_outages o ON g.generation_unit_code = o.asset_code
    WHERE g.area_map_code = \$1
      AND g.production_unit_status = 'COMMISSIONED'
      AND (o.asset_code IS NULL OR o.available_capacity_mw > 0)
    """
    # ... επεξεργασία αποτελεσμάτων
end
\end{minted}
\end{samepage}

\subsection{Παράμετροι Τύπου Καυσίμου}
\label{subsec:fuel-type-parameters}

Οι τεχνικοί περιορισμοί των μονάδων παραγωγής διαφέρουν σημαντικά ανάλογα με τον τύπο καυσίμου. Η δομή \texttt{FuelTypeParameters} συγκεντρώνει αυτές τις παραμέτρους:

\begin{samepage}
\begin{minted}[fontsize=\small,linenos,frame=lines]{julia}
struct FuelTypeParameters
    # Χαρακτηριστικά εκκίνησης (σε ΩΡΕΣ)
    cold_startup_time::Int          # Ώρες για ψυχρή εκκίνηση
    warm_startup_time::Int          # Ώρες για θερμή εκκίνηση
    hot_startup_time::Int           # Ώρες για καυτή εκκίνηση
    # Ελάχιστοι χρόνοι λειτουργίας (σε ΩΡΕΣ)
    min_uptime::Int                 # Ελάχιστες ώρες λειτουργίας
    min_downtime::Int               # Ελάχιστες ώρες αδράνειας
    # Ρυθμοί μεταβολής (κλάσμα δυναμικότητας ανά ΩΡΑ)
    ramp_up_rate::Float64           # Ρυθμός αύξησης (κλάσμα/ώρα)
    ramp_down_rate::Float64         # Ρυθμός μείωσης (κλάσμα/ώρα)
    # Λειτουργική ευελιξία
    min_load_factor::Float64        # Ελάχιστο φορτίο (κλάσμα)
    part_load_efficiency::Float64   # Απόδοση στο ελάχιστο φορτίο
    # Κατώφλια θερμοκρασιακής μετάβασης (ΩΡΕΣ εκτός)
    warm_threshold::Int             # Ώρες αδράνειας → θερμή κατάσταση
    cold_threshold::Int             # Ώρες αδράνειας → ψυχρή κατάσταση
    # Οικονομικές παράμετροι
    startup_cost_multiplier::Float64  # Πολλαπλασιαστής κόστους εκκίνησης
    no_load_cost_fraction::Float64    # Κόστος χωρίς φορτίο (κλάσμα)
end
\end{minted}
\end{samepage}

Σημαντικό σχεδιαστικό χαρακτηριστικό είναι ότι όλες οι χρονικές παράμετροι αποθηκεύονται σε \textbf{ώρες} (όχι σε περιόδους). Ο επιλυτής UC μετατρέπει αυτόματα σε περιόδους βάσει της πραγματικής χρονικής ανάλυσης των δεδομένων ($\text{periods\_per\_hour} = 60 / \text{resolution\_minutes}$). Τα πεδία \texttt{warm\_threshold} και \texttt{cold\_threshold} χρησιμοποιούνται τόσο στον προσδιορισμό θερμοκρασιακής κατάστασης εκκίνησης (Ενότητα~\ref{subsec:initial-conditions}) όσο και στους αντίστοιχους περιορισμούς του μοντέλου UC. Τα πεδία \texttt{startup\_cost\_multiplier} και \texttt{no\_load\_cost\_fraction} τροφοδοτούν τον υπολογισμό κόστους εκκίνησης και κόστους χωρίς φορτίο στην αντικειμενική συνάρτηση (Ενότητα~\ref{subsec:uc-balance-objective}).

Η συνάρτηση \texttt{get\_fuel\_type\_parameters} επιστρέφει τις κατάλληλες παραμέτρους για κάθε τύπο καυσίμου. Για παράδειγμα, οι λιγνιτικές μονάδες έχουν μεγαλύτερους χρόνους εκκίνησης και χαμηλότερους ρυθμούς μεταβολής σε σύγκριση με τις μονάδες συνδυασμένου κύκλου φυσικού αερίου.

\subsection{Φορτίο και Ανανεώσιμες Πηγές}
\label{subsec:load-renewables}

Η ζήτηση ηλεκτρικής ενέργειας αναπαρίσταται με τη δομή \texttt{Load}:

\begin{samepage}
\begin{minted}[fontsize=\small,linenos,frame=lines]{julia}
struct Load
    timeslot::String           # Χρονοθυρίδα (π.χ. "20240615-0000")
    resolution_code::String    # Χρονική ανάλυση ("PT15M", "PT60M")
    bidding_zone::String       # Ζώνη προσφοράς
    value::Float64             # Κατανάλωση (MW)
end
\end{minted}
\end{samepage}

Οι προβλέψεις παραγωγής από ΑΠΕ αναπαρίστανται με τη δομή \texttt{RenewablesGenerationForecast}:

\begin{samepage}
\begin{minted}[fontsize=\small,linenos,frame=lines]{julia}
struct RenewablesGenerationForecast
    date_time::String                    # Χρονοσφραγίδα
    resolution_code::String              # Χρονική ανάλυση
    bidding_zone::String                 # Ζώνη προσφοράς
    production_type::String              # "Solar" ή "Wind Onshore/Offshore"
    aggregated_generation_forecast::Float64  # Πρόβλεψη (MW)
end
\end{minted}
\end{samepage}

\subsection{Εξαγωγή Τεχνικών Παραμέτρων από Ιστορικά Δεδομένα}
\label{subsec:parameter-inference}

Οι τεχνικές παράμετροι των μονάδων παραγωγής (ρυθμοί μεταβολής, ελάχιστη ισχύς, ελάχιστοι χρόνοι λειτουργίας/αδράνειας) δεν είναι διαθέσιμες στα δημόσια δεδομένα του ENTSO-E. Για την αντιμετώπιση αυτού του περιορισμού, αναπτύχθηκε ένας μηχανισμός αυτόματης εξαγωγής (inference) αυτών των παραμέτρων από ιστορικά δεδομένα πραγματικής παραγωγής.

\paragraph{Πηγή Δεδομένων}
Τα ιστορικά δεδομένα παραγωγής ανακτώνται από τον πίνακα \texttt{actual\_generation\_output\_per\_generation\_unit} του ENTSO-E, ο οποίος περιέχει τη μετρούμενη παραγωγή κάθε μονάδας ανά 15 λεπτά ή ανά ώρα. Η συνάρτηση \texttt{get\_historical\_generation} ανακτά δεδομένα για παράθυρο 3 μηνών πριν από την επιλεγμένη ημερομηνία.

\paragraph{Εξαγωγή Ρυθμών Μεταβολής}
Οι ρυθμοί μεταβολής (ramp rates) υπολογίζονται από τη συνάρτηση \texttt{infer\_ramp\_rates}:

\begin{samepage}
\begin{minted}[fontsize=\small,linenos,frame=lines]{julia}
function infer_ramp_rates(historical_data::DataFrame, p_max::Float64;
                          percentile::Float64=0.95)
    if nrow(historical_data) < MIN_DATA_POINTS_FOR_RAMP_INFERENCE
        return (ramp_up=nothing, ramp_down=nothing)
    end

    # Ταξινόμηση κατά χρόνο και υπολογισμός μεταβολών
    sort!(historical_data, :date_time_utc)
    gen_values = historical_data.actual_generation_output_mw
    deltas = diff(gen_values)

    # Διαχωρισμός σε αυξήσεις και μειώσεις
    ramp_ups = filter(x -> x > 0, deltas)
    ramp_downs = map(abs, filter(x -> x < 0, deltas))

    # Κανονικοποίηση σε ωριαίο ρυθμό βάσει ανάλυσης δεδομένων
    resolution_code = historical_data.resolution_code[1]
    resolution_minutes = parse_resolution_to_minutes(resolution_code)
    hourly_factor = 60.0 / resolution_minutes

    # 95ο εκατοστημόριο, μετατροπή σε κλάσμα p_max/ώρα
    ramp_up = isempty(ramp_ups) ? nothing :
        (quantile(ramp_ups, percentile) * hourly_factor) / p_max
    ramp_down = isempty(ramp_downs) ? nothing :
        (quantile(ramp_downs, percentile) * hourly_factor) / p_max

    return (ramp_up=ramp_up, ramp_down=ramp_down)
end
\end{minted}
\end{samepage}

Οι ρυθμοί αποθηκεύονται ως κλάσμα της μέγιστης ισχύος ανά ώρα (π.χ. 0.20 σημαίνει ότι η μονάδα μπορεί να μεταβάλει την παραγωγή της κατά 20\% της $P^{\max}$ ανά ώρα), επιτρέποντας την εύκολη μετατροπή σε οποιαδήποτε χρονική ανάλυση.

\paragraph{Εξαγωγή Ελάχιστης Ισχύος}
Η ελάχιστη ισχύς ($P^{\min}$) εξάγεται αναλύοντας τις περιόδους σταθερής λειτουργίας, εξαιρώντας τις μεταβατικές καταστάσεις (εκκίνηση/τερματισμός):

\begin{samepage}
\begin{minted}[fontsize=\small,linenos,frame=lines]{julia}
function infer_p_min(historical_data::DataFrame, p_max::Float64,
                     fuel_type::Symbol; percentile::Float64=0.05)
    if nrow(historical_data) < MIN_DATA_POINTS_FOR_RAMP_INFERENCE
        return nothing
    end

    sort!(historical_data, :date_time_utc)
    gen_values = historical_data.actual_generation_output_mw
    deltas = diff(gen_values)

    # Ανίχνευση μεταβατικών καταστάσεων (|Δ| > 5% της p_max)
    ramp_threshold = 0.05 * p_max

    # Συλλογή σταθερών μη-μηδενικών τιμών
    stable_values = Float64[]
    for i in 2:(length(gen_values) - 1)
        value = gen_values[i]
        if value > 0 && abs(deltas[i-1]) < ramp_threshold &&
                        abs(deltas[i]) < ramp_threshold
            push!(stable_values, value)
        end
    end

    if length(stable_values) < 50
        return nothing  # Ανεπαρκή δεδομένα
    end

    # 5ο εκατοστημόριο της σταθερής λειτουργίας
    inferred_p_min = quantile(stable_values, percentile)

    # Εφαρμογή ορίων ανά τύπο καυσίμου (clamp σε MW)
    (min_frac, max_frac) = P_MIN_BOUNDS[get_p_min_bounds_category(fuel_type)]
    return clamp(inferred_p_min, min_frac * p_max, max_frac * p_max)
end
\end{minted}
\end{samepage}

\paragraph{Όρια ανά Τύπο Καυσίμου}
Για την αποφυγή μη ρεαλιστικών τιμών, εφαρμόζονται όρια (bounds) ανά τύπο καυσίμου. Τα όρια αυτά βασίζονται στα τεχνικά χαρακτηριστικά κάθε τεχνολογίας παραγωγής:

\begin{table}[htbp]
\centering
\caption{Όρια ελάχιστης ισχύος ($P^{\min}$) ανά τύπο καυσίμου}
\label{tab:p-min-bounds}
\begin{tabular}{@{}lcc@{}}
\toprule
\textbf{Τύπος Καυσίμου} & \textbf{Ελάχιστο (\%)} & \textbf{Μέγιστο (\%)} \\
\midrule
Λιγνίτης/Άνθρακας & 45 & 65 \\
Φυσικό αέριο (CCGT) & 35 & 55 \\
Φυσικό αέριο (OCGT) & 20 & 45 \\
Υδροηλεκτρικά/Μπαταρίες & 0 & 0 \\
Προεπιλογή & 30 & 60 \\
\bottomrule
\end{tabular}
\end{table}

Οι ευέλικτοι πόροι (υδροηλεκτρικά, μπαταρίες) λαμβάνουν $P^{\min} = 0$ χωρίς εξαγωγή, καθώς μπορούν να λειτουργήσουν σε οποιοδήποτε επίπεδο ισχύος. Η αντίστοιχη σταθερά \texttt{FLEXIBLE\_FUEL\_TYPES} περιλαμβάνει τους τύπους \texttt{Hydro Pumped Storage}, \texttt{Hydro Run-of-river}, \texttt{Hydro Water Reservoir}, \texttt{Battery} και \texttt{Other}. Αξίζει να σημειωθεί ότι οι προεπιλεγμένες τιμές \texttt{min\_load\_factor} στη δομή \texttt{FuelTypeParameters} ευθυγραμμίστηκαν με τα όρια εξαγωγής (π.χ. λιγνίτης: 0.45, ευέλικτοι πόροι: 0.0).

\paragraph{Εξαγωγή Ελάχιστων Χρόνων Λειτουργίας/Αδράνειας}
Οι ελάχιστοι χρόνοι υπολογίζονται αναλύοντας τους κύκλους λειτουργίας (on/off) της μονάδας:

\begin{samepage}
\begin{minted}[fontsize=\small,linenos,frame=lines]{julia}
function infer_uptime_downtime(historical_data::DataFrame, fuel_type::Symbol;
                               percentile::Float64=0.05)
    if nrow(historical_data) < MIN_DATA_POINTS_FOR_RAMP_INFERENCE
        return (nothing, nothing)
    end

    sort!(historical_data, :date_time_utc)
    gen_values = historical_data.actual_generation_output_mw

    # Υπολογισμός διάρκειας περιόδου σε ώρες
    resolution_minutes = parse_resolution_to_minutes(
        historical_data.resolution_code[1])
    period_hours = resolution_minutes / 60.0

    # Προσδιορισμός κατάστασης (on: > 1 MW, off: = 0)
    is_on = gen_values .> 1.0

    # Εντοπισμός κύκλων λειτουργίας
    on_durations, off_durations = Float64[], Float64[]
    current_state = is_on[1]
    current_duration = 1

    for i in 2:length(is_on)
        if is_on[i] == current_state
            current_duration += 1
        else
            duration_hours = current_duration * period_hours
            if current_state
                push!(on_durations, duration_hours)
            else
                push!(off_durations, duration_hours)
            end
            current_state = is_on[i]
            current_duration = 1
        end
    end

    # Φιλτράρισμα σύντομων glitches (< 2 περιόδων)
    min_dur = 2 * period_hours
    filter!(d -> d >= min_dur, on_durations)
    filter!(d -> d >= min_dur, off_durations)

    # Απαίτηση τουλάχιστον 5 κύκλων για αξιόπιστη εξαγωγή
    # Εφαρμογή ορίων ανά τύπο καυσίμου (βλ. Πίνακα~\ref{tab:uptime-downtime-bounds})
    bounds = UPTIME_DOWNTIME_BOUNDS[get_p_min_bounds_category(fuel_type)]
    min_uptime = length(on_durations) >= 5 ?
        round(Int, clamp(quantile(on_durations, percentile),
                         bounds[1]...)) : nothing
    min_downtime = length(off_durations) >= 5 ?
        round(Int, clamp(quantile(off_durations, percentile),
                         bounds[2]...)) : nothing

    return (min_uptime, min_downtime)
end
\end{minted}
\end{samepage}

Ομοίως με τα όρια $P^{\min}$, εφαρμόζονται όρια ανά τύπο καυσίμου για τους χρόνους λειτουργίας και αδράνειας, όπως παρουσιάζονται στον Πίνακα~\ref{tab:uptime-downtime-bounds}.

\begin{table}[htbp]
\centering
\caption{Όρια ελάχιστων χρόνων λειτουργίας/αδράνειας ανά τύπο καυσίμου (ώρες)}
\label{tab:uptime-downtime-bounds}
\begin{tabular}{@{}lcccc@{}}
\toprule
\textbf{Τύπος Καυσίμου} & \multicolumn{2}{c}{\textbf{Χρόνος Λειτουργίας}} & \multicolumn{2}{c}{\textbf{Χρόνος Αδράνειας}} \\
\cmidrule(lr){2-3} \cmidrule(lr){4-5}
& \textbf{Ελάχ.} & \textbf{Μέγ.} & \textbf{Ελάχ.} & \textbf{Μέγ.} \\
\midrule
Λιγνίτης/Άνθρακας & 8 & 48 & 4 & 24 \\
Φυσικό αέριο (CCGT) & 2 & 12 & 1 & 8 \\
Φυσικό αέριο (OCGT) & 1 & 4 & 1 & 4 \\
Προεπιλογή & 2 & 24 & 1 & 12 \\
\bottomrule
\end{tabular}
\end{table}

Η εφαρμογή ορίων διασφαλίζει ότι οι εξαχθείσες τιμές παραμένουν εντός τεχνικά ρεαλιστικού εύρους. Για μονάδες βάσης (π.χ. λιγνιτικές), η εξαγωγή χρόνων συχνά δεν είναι εφικτή λόγω σπάνιων κύκλων λειτουργίας/αδράνειας, οπότε χρησιμοποιούνται οι προεπιλεγμένες τιμές τύπου καυσίμου.

\paragraph{Ενοποίηση Εξαγωγής Παραμέτρων}
Η κεντρική συνάρτηση \texttt{infer\_parameters\_for\_generators} ενορχηστρώνει την εξαγωγή για ένα σύνολο μονάδων, εφαρμόζοντας διαφορετική λογική ανάλογα με τον τύπο καυσίμου:

\begin{samepage}
\begin{minted}[fontsize=\small,linenos,frame=lines]{julia}
function infer_parameters_for_generators(generators::Vector{Generator},
                                         day::Date)
    updated = Generator[]
    for gen in generators
        historical = get_historical_generation(gen.code, day)
        rates = infer_ramp_rates(historical, gen.p_max)

        if gen.fuel_type in FLEXIBLE_FUEL_TYPES
            # Ευέλικτοι πόροι: p_min = 0, χωρίς uptime/downtime
            push!(updated, Generator(gen.code, gen.name, gen.fuel_type,
                gen.location, gen.p_max, 0.0, gen.bidding_zone,
                gen.marginal_cost, rates.ramp_up, rates.ramp_down,
                nothing, nothing))
        else
            # Θερμικές μονάδες: εξαγωγή όλων των παραμέτρων
            p_min = infer_p_min(historical, gen.p_max, gen.fuel_type)
            (uptime, downtime) = infer_uptime_downtime(historical,
                                                        gen.fuel_type)
            push!(updated, Generator(gen.code, gen.name, gen.fuel_type,
                gen.location, gen.p_max,
                p_min !== nothing ? p_min : gen.p_min,
                gen.bidding_zone, gen.marginal_cost,
                rates.ramp_up, rates.ramp_down, uptime, downtime))
        end
    end
    return updated
end
\end{minted}
\end{samepage}

\paragraph{Αποθήκευση σε Κρυφή Μνήμη (Caching)}
Η εξαγωγή παραμέτρων από ιστορικά δεδομένα είναι υπολογιστικά ακριβή (περίπου 17 λεπτά για 39 μονάδες). Για τη βελτίωση της απόδοσης, οι εξαχθείσες παράμετροι αποθηκεύονται στον πίνακα \texttt{simulations.generator\_inferred\_parameters}:

\begin{samepage}
\begin{minted}[fontsize=\small,linenos,frame=lines]{julia}
# Ανάκτηση μονάδων με αυτόματη εξαγωγή και caching
generators = get_generators_with_inferred_params("GR", Date(2024, 6, 15);
    use_cache = true,       # Χρήση cached τιμών αν διαθέσιμες
    max_age_days = 7        # Ανανέωση cache κάθε 7 ημέρες
)
\end{minted}
\end{samepage}

Η χρήση της κρυφής μνήμης μειώνει τον χρόνο ανάκτησης από περίπου 17 λεπτά σε περίπου 2 δευτερόλεπτα---βελτίωση κατά 525 φορές.


\subsection{Αρχικές Συνθήκες Μονάδων Παραγωγής}
\label{subsec:initial-conditions}

Για τη ρεαλιστική μοντελοποίηση του προβλήματος Unit Commitment, η κατάσταση κάθε μονάδας κατά τη χρονική στιγμή $t=0$ (έναρξη του ορίζοντα βελτιστοποίησης) πρέπει να είναι γνωστή. Οι αρχικές συνθήκες καθορίζουν αν η μονάδα λειτουργούσε, σε ποιο επίπεδο ισχύος, και για πόσο χρόνο βρισκόταν στην τρέχουσα κατάστασή της. Αυτή η πληροφορία είναι απαραίτητη για τη σωστή επιβολή περιορισμών στις πρώτες περιόδους του ορίζοντα.

\paragraph{Δομή Αρχικών Συνθηκών}
Η δομή \texttt{InitialConditions} αναπαριστά την κατάσταση μιας μονάδας στο $t=0$:

\begin{samepage}
\begin{minted}[fontsize=\small,linenos,frame=lines]{julia}
struct InitialConditions
    is_on::Bool           # Λειτουργούσε η μονάδα στο t=0;
    output::Float64       # Ισχύς εξόδου σε MW (0 αν εκτός λειτουργίας)
    hours_on::Int         # Συνεχόμενες ώρες λειτουργίας
    hours_off::Int        # Συνεχόμενες ώρες αδράνειας
    thermal_state::Symbol # Θερμοκρασιακή κατάσταση (:hot, :warm, :cold)
end
\end{minted}
\end{samepage}

Η θερμοκρασιακή κατάσταση προσδιορίζεται από τον χρόνο αδράνειας: \texttt{:hot} αν $T_{\text{off}} \leq 8$ ώρες, \texttt{:warm} αν $8 < T_{\text{off}} \leq 48$ ώρες, και \texttt{:cold} αν $T_{\text{off}} > 48$ ώρες. Τα κατώφλια αυτά προσαρμόζονται ανά τύπο καυσίμου μέσω των αντίστοιχων πεδίων \texttt{warm\_threshold} και \texttt{cold\_threshold} της δομής \texttt{FuelTypeParameters}.

\paragraph{Εξαγωγή από Ιστορικά Δεδομένα}
Η κεντρική συνάρτηση \texttt{get\_initial\_conditions} εξάγει τις αρχικές συνθήκες για ένα σύνολο μονάδων χρησιμοποιώντας ένα μαζικό ερώτημα SQL (batch query) αντί $N$ ξεχωριστών ερωτημάτων. Η τεχνική αυτή επιτυγχάνει βελτίωση ταχύτητας κατά 36 φορές (από $\sim$9 λεπτά σε $\sim$15 δευτερόλεπτα για 40 μονάδες):

\begin{samepage}
\begin{minted}[fontsize=\small,linenos,frame=lines]{julia}
function get_initial_conditions(generators::Vector{Generator},
                                market_day::Date;
                                use_historical::Bool=true)
    # Μαζική ανάκτηση 72 ωρών ιστορικού για όλες τις μονάδες
    # σε ένα ερώτημα SQL (WHERE generation_unit_code = ANY($1))
    day_before = market_day - Day(1)
    start_utc = DateTime(day_before, Time(23, 0, 0))  # ~00:00 CET
    codes = [gen.code for gen in generators]
    all_historical = get_recent_generation_batch(codes, start_utc)

    conditions = Dict{String, InitialConditions}()
    for gen in generators
        gen_data = filter(r -> r.generation_unit_code == gen.code,
                          all_historical)
        ic = infer_initial_conditions_from_data(gen_data, gen.fuel_type)
        conditions[gen.code] = ic
    end
    return conditions
end
\end{minted}
\end{samepage}

Ο αλγόριθμος εξαγωγής αναλύει τα δεδομένα σε φθίνουσα χρονική σειρά: η πιο πρόσφατη μέτρηση $> 1$ MW δηλώνει ότι η μονάδα ήταν σε λειτουργία. Στη συνέχεια, μετρούνται οι συνεχόμενες περίοδοι στην ίδια κατάσταση για τον υπολογισμό του $T_{\text{on}}$ ή $T_{\text{off}}$. Σε περίπτωση απουσίας δεδομένων, χρησιμοποιούνται ευρετικές προεπιλεγμένες τιμές ανά τύπο καυσίμου: οι μονάδες βάσης (λιγνίτης, πυρηνικά) θεωρούνται σε λειτουργία, οι μονάδες μέσου φορτίου (CCGT) εκτός λειτουργίας αλλά θερμές, οι μονάδες αιχμής (OCGT) εκτός λειτουργίας και ψυχρές, και οι ευέλικτοι πόροι (υδροηλεκτρικά, μπαταρίες) εκτός λειτουργίας αλλά έτοιμοι.


\subsection{Τύποι Εντολών Αγοράς}
\label{subsec:market-orders}

Οι εντολές αγοράς υλοποιούνται ως ιεραρχία τύπων με αφηρημένο βασικό τύπο \texttt{MarketOrder}. Ο κύριος τύπος που χρησιμοποιείται είναι η απλή ωριαία εντολή \texttt{SimpleOrder}:

\begin{samepage}
\begin{minted}[fontsize=\small,linenos,frame=lines]{julia}
abstract type MarketOrder end

struct SimpleOrder <: MarketOrder
    type::Symbol           # :supply ή :demand
    price::Float64         # Τιμή (€/MWh)
    quantity::Float64      # Ποσότητα (MWh)
    zone::Symbol           # Ζώνη προσφοράς
    date_time::DateTime    # Χρόνος παράδοσης
    resolution_code::Int   # Ανάλυση σε λεπτά (60, 30, 15)
end
\end{minted}
\end{samepage}

Πέρα από τις απλές εντολές, υλοποιούνται και σύνθετοι τύποι όπως οι εντολές μπλοκ (\texttt{BlockOrder}), οι συνδεδεμένες εντολές μπλοκ (\texttt{LinkedBlockOrder}) και οι αποκλειστικές εντολές μπλοκ (\texttt{ExclusiveBlockOrder}), ακολουθώντας την προδιαγραφή του αλγορίθμου EUPHEMIA.


% =============================================================================
\section{Επιλυτής Unit Commitment}
\label{sec:uc-solver}

Ο επιλυτής Unit Commitment υλοποιείται στο module \texttt{UnitCommitment.jl} και αποτελεί το πρώτο βήμα της αλυσίδας υπολογισμών. Η κύρια συνάρτηση \texttt{solve\_unit\_commitment} δέχεται ως είσοδο τη ζώνη προσφοράς και την ημερομηνία, ανακτά τα απαραίτητα δεδομένα από τη βάση, διαμορφώνει το πρόβλημα βελτιστοποίησης και επιστρέφει τη βέλτιστη λύση.

\subsection{Δημιουργία Μοντέλου}
\label{subsec:uc-model-creation}

Η δημιουργία του μοντέλου JuMP ξεκινά με την επιλογή του επιλυτή και τον ορισμό των μεταβλητών απόφασης:

\begin{samepage}
\begin{minted}[fontsize=\small,linenos,frame=lines]{julia}
function solve_unit_commitment(bidding_zone::String, day::Date;
                               optimizer::String="auto",
                               mip_gap::Float64=0.01,
                               time_limit::Float64=600.0)
    # Φόρτωση δεδομένων
    generators = get_generators(bidding_zone, day)
    loads = get_loads(bidding_zone, day)
    renewables = get_generation_forecast_for_wind_and_solar(bidding_zone, day)
    initial_conditions = get_initial_conditions(generators, day)

    N = length(generators)  # Αριθμός μονάδων
    T = 24                  # Χρονικές περίοδοι
    periods_per_hour = 60 / resolution_minutes

    # Δημιουργία μοντέλου με ρύθμιση παραμέτρων επιλυτή
    model = Model(optimizer_func)
    set_silent(model)
    set_optimizer_attribute(model, MOI.RelativeGapTolerance(), mip_gap)
    set_optimizer_attribute(model, MOI.TimeLimitSec(), time_limit)

    # Εξαγωγή αρχικών συνθηκών σε πίνακες
    u0 = zeros(Int, N)      # Αρχική δέσμευση (0 ή 1)
    g0 = zeros(Float64, N)  # Αρχική παραγωγή (MW)
    T_on0 = zeros(Int, N)   # Ώρες ήδη σε λειτουργία
    T_off0 = zeros(Int, N)  # Ώρες ήδη εκτός λειτουργίας
    for (i, gen) in enumerate(generators)
        ic = initial_conditions[gen.code]
        u0[i] = ic.is_on ? 1 : 0
        g0[i] = ic.output
        T_on0[i] = ic.hours_on
        T_off0[i] = ic.hours_off
    end

    # Μεταβλητές παραγωγής (συνεχείς)
    @variable(model, 0 <= g[i=1:N, t=1:T] <= generators[i].p_max)

    # Μεταβλητές δέσμευσης (δυαδικές)
    @variable(model, u[i=1:N, t=1:T], Bin)  # Κατάσταση λειτουργίας
    @variable(model, v[i=1:N, t=1:T], Bin)  # Εκκίνηση
    @variable(model, z[i=1:N, t=1:T], Bin)  # Τερματισμός

    # Μεταβλητές για θερμοκρασιακές καταστάσεις εκκίνησης
    Θ = [:cold, :warm, :hot]
    @variable(model, v_θ[i=1:N, θ in Θ, t=1:T], Bin)

    # ... περαιτέρω μεταβλητές
end
\end{minted}
\end{samepage}

\subsection{Περιορισμοί Ισχύος}
\label{subsec:impl-uc-power-constraints}

Οι περιορισμοί ισχύος εξασφαλίζουν ότι η παραγωγή κάθε μονάδας βρίσκεται εντός των τεχνικών της ορίων:

\begin{samepage}
\begin{minted}[fontsize=\small,linenos,frame=lines]{julia}
# Ανώτατο όριο παραγωγής
@constraint(model, [i = 1:N, t = 1:T], 
    g[i, t] <= generators[i].p_max * u[i, t])

# Κατώτατο όριο παραγωγής με παραμέτρους ανά τύπο καυσίμου
for (i, gen) in enumerate(generators)
    params = fuel_params[i]
    # Το ελάχιστο είναι το μέγιστο μεταξύ p_min και του ελάχιστου 
    # φορτίου βάσει τύπου καυσίμου
    min_gen = max(gen.p_min, params.min_load_factor * gen.p_max)
    @constraint(model, [t = 1:T], g[i, t] >= min_gen * u[i, t])
end
\end{minted}
\end{samepage}

\subsection{Περιορισμοί Δέσμευσης}
\label{subsec:impl-uc-commitment-constraints}

Η σχέση μεταξύ κατάστασης λειτουργίας, εκκίνησης και τερματισμού υλοποιείται ως εξής, με ενσωμάτωση αρχικών συνθηκών στην περίοδο $t=1$:

\begin{samepage}
\begin{minted}[fontsize=\small,linenos,frame=lines]{julia}
# Σύνδεση t=1 με αρχική κατάσταση u0
if initial_conditions !== nothing
    @constraint(model, [i in 1:N],
        u[i, 1] == u0[i] + v[i, 1] - z[i, 1])
end

# Σύνδεση μεταβλητών δέσμευσης (t ≥ 2)
@constraint(model, [i in 1:N, t in 2:T],
    u[i, t] == u[i, t-1] + v[i, t] - z[i, t])

# Αποκλεισμός ταυτόχρονης εκκίνησης και τερματισμού
@constraint(model, [i = 1:N, t = 1:T], v[i, t] + z[i, t] <= 1)

# Ελάχιστος χρόνος λειτουργίας (ώρες → περίοδοι)
for (i, gen) in enumerate(generators)
    UT_hours = gen.min_uptime !== nothing ?
        gen.min_uptime : fuel_params[i].min_uptime
    UT = max(1, ceil(Int, UT_hours * periods_per_hour))
    if UT > 1
        @constraint(model, [t = UT:T],
            sum(v[i, τ] for τ in t-UT+1:t) <= u[i, t])
        # Αρχική συνθήκη: εναπομένων χρόνος λειτουργίας
        if initial_conditions !== nothing && u0[i] == 1
            remaining = max(0, UT - ceil(Int, T_on0[i] * periods_per_hour))
            for t in 1:min(remaining, T)
                @constraint(model, u[i, t] == 1)
            end
        end
    end
end

# Ελάχιστος χρόνος αδράνειας (ώρες → περίοδοι)
for (i, gen) in enumerate(generators)
    DT_hours = gen.min_downtime !== nothing ?
        gen.min_downtime : fuel_params[i].min_downtime
    DT = max(1, ceil(Int, DT_hours * periods_per_hour))
    if DT > 1
        @constraint(model, [t = DT:T],
            sum(z[i, τ] for τ in t-DT+1:t) <= 1 - u[i, t])
        # Αρχική συνθήκη: εναπομένων χρόνος αδράνειας
        if initial_conditions !== nothing && u0[i] == 0
            remaining = max(0, DT - ceil(Int, T_off0[i] * periods_per_hour))
            for t in 1:min(remaining, T)
                @constraint(model, u[i, t] == 0)
            end
        end
    end
end
\end{minted}
\end{samepage}

Η μετατροπή $\texttt{UT\_hours} \times \texttt{periods\_per\_hour}$ διασφαλίζει ότι οι περιορισμοί εφαρμόζονται σωστά ανεξάρτητα από τη χρονική ανάλυση (15, 30 ή 60 λεπτά). Οι αρχικές συνθήκες επιβάλλουν τη συνέχιση της τρέχουσας κατάστασης κάθε μονάδας: αν μια μονάδα λειτουργούσε αλλά δεν έχει συμπληρώσει τον ελάχιστο χρόνο λειτουργίας, εξαναγκάζεται να παραμείνει σε λειτουργία για τις εναπομένουσες περιόδους.

\subsection{Περιορισμοί Ρυθμού Μεταβολής}
\label{subsec:impl-uc-ramping}

Οι περιορισμοί ρυθμού μεταβολής (ramping) χρησιμοποιούν την τεχνική Big-M για την αποδέσμευση κατά τις μεταβατικές καταστάσεις:

\begin{samepage}
\begin{minted}[fontsize=\small,linenos,frame=lines]{julia}
# Υπολογισμός ρυθμών μεταβολής ανά μονάδα:
# εξαχθείσα τιμή (fraction/hour) ή προεπιλογή τύπου καυσίμου
# Μετατροπή σε MW/περίοδο: rate × p_max × period_hours
period_hours = resolution_minutes / 60.0
ramp_up = Float64[]
ramp_down = Float64[]
for (i, gen) in enumerate(generators)
    if gen.ramp_up !== nothing
        push!(ramp_up, gen.ramp_up * gen.p_max * period_hours)
    else
        push!(ramp_up, fuel_params[i].ramp_up_rate * gen.p_max * period_hours)
    end
    if gen.ramp_down !== nothing
        push!(ramp_down, gen.ramp_down * gen.p_max * period_hours)
    else
        push!(ramp_down, fuel_params[i].ramp_down_rate * gen.p_max * period_hours)
    end
end

# Big-M: κλιμακούμενο βάσει μεγέθους προβλήματος
max_p_max = maximum(gen.p_max for gen in generators)
M = 2.0 * max_p_max

# Περιορισμός αύξησης (t ≥ 2)
@constraint(model, [i in 1:N, t in 2:T],
    g[i, t] - g[i, t-1] <= ramp_up[i] + M * u_SU[i, t])

# Περιορισμός μείωσης (t ≥ 2)
@constraint(model, [i in 1:N, t in 2:T],
    g[i, t-1] - g[i, t] <= ramp_down[i] + M * u_SD[i, t])
\end{minted}
\end{samepage}

\subsection{Ισορροπία Ισχύος και Αντικειμενική Συνάρτηση}
\label{subsec:uc-balance-objective}

Ο περιορισμός ισορροπίας ισχύος εξασφαλίζει την κάλυψη της καθαρής ζήτησης:

\begin{samepage}
\begin{minted}[fontsize=\small,linenos,frame=lines]{julia}
# Περιορισμός ισορροπίας
@constraint(model, [t in 1:T],
    sum(g[i, t] for i in 1:N) == net_demand[t])
\end{minted}
\end{samepage}

Η αντικειμενική συνάρτηση αποτελείται από τρεις συνιστώσες κόστους: κόστος παραγωγής, κόστος εκκίνησης και κόστος χωρίς φορτίο. Πριν τη διαμόρφωση, υπολογίζονται οι παράμετροι κόστους:

\begin{samepage}
\begin{minted}[fontsize=\small,linenos,frame=lines]{julia}
# Κόστος εκκίνησης ανά μονάδα και θερμοκρασιακή κατάσταση (€)
# Βασικό κόστος = startup_cost_multiplier × marginal_cost × p_max
# Πολλαπλασιαστές: hot=1.0, warm=1.5, cold=2.5
C_SU = Dict{Tuple{Int,Symbol},Float64}()
for i in 1:N
    params = fuel_params[i]
    gen = generators[i]
    base = params.startup_cost_multiplier * gen.marginal_cost * gen.p_max
    C_SU[(i, :hot)]  = base * 1.0
    C_SU[(i, :warm)] = base * 1.5
    C_SU[(i, :cold)] = base * 2.5
end

# Κόστος χωρίς φορτίο ανά μονάδα (€/περίοδο)
# Πάγιο κόστος δέσμευσης ανεξάρτητο από παραγωγή
period_hours = resolution_minutes / 60.0
C_NL = Float64[]
for i in 1:N
    params = fuel_params[i]
    gen = generators[i]
    push!(C_NL, params.no_load_cost_fraction *
                gen.marginal_cost * gen.p_min * period_hours)
end

# Αντικειμενική: ελαχιστοποίηση συνολικού κόστους
@objective(model, Min,
    # 1. Κόστος παραγωγής (μεταβλητό)
    sum(generators[i].marginal_cost * g[i, t] for i in 1:N, t in 1:T)
    # 2. Κόστος εκκίνησης (εξαρτώμενο από θερμοκρασία)
    + sum(C_SU[(i, θ)] * v_θ[i, θ, t] for i in 1:N, θ in Θ, t in 1:T)
    # 3. Κόστος χωρίς φορτίο (πάγιο κόστος δέσμευσης)
    + sum(C_NL[i] * u[i, t] for i in 1:N, t in 1:T)
)
\end{minted}
\end{samepage}

Η τρισδιάστατη δομή κόστους αποτυπώνει ρεαλιστικά τα κίνητρα λειτουργίας: το κόστος εκκίνησης αποθαρρύνει τις συχνές εκκινήσεις/τερματισμούς (ιδιαίτερα τις ψυχρές εκκινήσεις, που κοστίζουν 2,5 φορές περισσότερο από τις θερμές), ενώ το κόστος χωρίς φορτίο αντανακλά τα πάγια λειτουργικά κόστη δέσμευσης μιας μονάδας ανεξάρτητα από το επίπεδο παραγωγής της.

\subsection{Επίλυση και Επιστροφή Αποτελεσμάτων}
\label{subsec:uc-solution}

Μετά την επίλυση, τα αποτελέσματα επιστρέφονται σε δομημένη μορφή:

\begin{samepage}
\begin{minted}[fontsize=\small,linenos,frame=lines]{julia}
optimize!(model)

if termination_status(model) == OPTIMAL
    return (
        status = OPTIMAL,
        generators = generators,
        generation = value.(g),       # Πίνακας παραγωγής N×T
        commitment = value.(u),       # Πίνακας δέσμευσης N×T
        startup = value.(v),          # Πίνακας εκκινήσεων N×T
        shutdown = value.(z),         # Πίνακας τερματισμών N×T
        objective_value = objective_value(model),
        solve_time = solve_time
    )
else
    return (status = termination_status(model),)
end
\end{minted}
\end{samepage}

\subsection{Αποθήκευση Αποτελεσμάτων σε Κρυφή Μνήμη}
\label{subsec:uc-caching}

Η επίλυση του προβλήματος Unit Commitment απαιτεί 10--15 λεπτά ανά ζώνη. Για την αποφυγή επαναληπτικών υπολογισμών, τα αποτελέσματα αποθηκεύονται σε τρεις πίνακες PostgreSQL (\texttt{uc\_results}, \texttt{uc\_generation}, \texttt{uc\_net\_demand}) με κλειδί \texttt{(bidding\_zone, market\_date, code\_version)}. Η συνάρτηση \texttt{solve\_unit\_commitment} ελέγχει αυτόματα τη διαθεσιμότητα αποθηκευμένων αποτελεσμάτων:

\begin{samepage}
\begin{minted}[fontsize=\small,linenos,frame=lines]{julia}
function solve_unit_commitment(bidding_zone::String, day::Date;
                               optimizer::String="auto",
                               use_cache::Bool=true,
                               force_rerun::Bool=false,
                               # ... λοιπές παράμετροι
                               )
    # Έλεγχος κρυφής μνήμης (εκτός αν απενεργοποιημένη)
    if use_cache && !force_rerun
        cached = load_uc_results(bidding_zone, day)
        if cached !== nothing
            return cached   # Επιστροφή σε ~2 δευτερόλεπτα
        end
    end

    # ... κανονική επίλυση (10-15 λεπτά) ...

    # Αποθήκευση αποτελεσμάτων (delete-before-insert)
    if use_cache
        save_uc_results(solution, bidding_zone, day)
    end
    return solution
end
\end{minted}
\end{samepage}

Η παράμετρος \texttt{force\_rerun} διαδίδεται μέσω ολόκληρης της αλυσίδας κλήσεων (\texttt{generate\_energy\_prices} $\to$ \texttt{create\_typed\_order\_book} $\to$ \texttt{generate\_market\_orders\_from\_uc} $\to$ \texttt{solve\_unit\_commitment}), επιτρέποντας τον επανυπολογισμό αποτελεσμάτων όταν αλλάξει ο κώδικας ή τα δεδομένα εισόδου. Η αποθήκευση ανέρχεται σε περίπου 195 KB ανά ζώνη/ημέρα---κόστος αμελητέο σε σύγκριση με την εξοικονόμηση χρόνου.


% =============================================================================
\section{Μηχανή Στρατηγικής Προσφορών}
\label{sec:bidding-strategy-engine}

Η μηχανή στρατηγικής προσφορών υλοποιείται στο module \texttt{BiddingStrategy.jl} και μετατρέπει τα αποτελέσματα του Unit Commitment σε εντολές αγοράς. Η κύρια συνάρτηση \texttt{generate\_market\_orders\_from\_uc} υποστηρίζει τρεις διακριτές στρατηγικές.

\subsection{Δομή Αποτελέσματος}
\label{subsec:bidding-result}

Τα αποτελέσματα της μετατροπής επιστρέφονται στη δομή \texttt{UCToBidsResult}:

\begin{samepage}
\begin{minted}[fontsize=\small,linenos,frame=lines]{julia}
struct UCToBidsResult
    supply_orders::Vector{SimpleOrder}   # Εντολές πώλησης
    demand_orders::Vector{SimpleOrder}   # Εντολές αγοράς
    bidding_zone::Symbol                 # Ζώνη προσφοράς
    day::Date                            # Ημερομηνία
    total_supply_quantity::Float64       # Συνολική προσφερόμενη ποσότητα
    total_demand_quantity::Float64       # Συνολική ζητούμενη ποσότητα
    success::Bool                        # Επιτυχία μετατροπής
    message::String                      # Μήνυμα κατάστασης
end
\end{minted}
\end{samepage}

\subsection{Στρατηγική Δεσμευμένων Μονάδων}
\label{subsec:committed-only-strategy}

Η συντηρητική στρατηγική \texttt{:committed\_only} δημιουργεί προσφορές μόνο για τις μονάδες που έχουν δεσμευτεί στη λύση του UC:

\begin{samepage}
\begin{minted}[fontsize=\small,linenos,frame=lines]{julia}
function apply_committed_only_strategy!(orders, uc_solution, generators, 
                                        day, markup_factor)
    for (i, gen) in enumerate(generators)
        for t in 1:24
            # Έλεγχος αν η μονάδα είναι δεσμευμένη
            if uc_solution.commitment[i, t] > COMMITMENT_THRESHOLD
                generation = uc_solution.generation[i, t]
                if generation > GENERATION_THRESHOLD
                    # Δημιουργία εντολής πώλησης
                    price = gen.marginal_cost * markup_factor
                    dt = DateTime(day, Time(t-1, 0))
                    push!(orders, SimpleOrder(
                        :supply, price, generation, 
                        Symbol(gen.bidding_zone), dt, 60
                    ))
                end
            end
        end
    end
end
\end{minted}
\end{samepage}

\subsection{Στρατηγική Μέγιστης Διαθεσιμότητας}
\label{subsec:all-available-max}

Η στρατηγική \texttt{:all\_available\_max} προσφέρει όλες τις μονάδες στη μέγιστη δυναμικότητά τους, ανεξάρτητα από τη λύση του UC. Αν και μη ρεαλιστική, χρησιμεύει ως βάση σύγκρισης:

\begin{samepage}
\begin{minted}[fontsize=\small,linenos,frame=lines]{julia}
function apply_all_available_max_strategy!(orders, generators, day, markup_factor)
    for gen in generators
        for t in 1:24
            price = gen.marginal_cost * markup_factor
            dt = DateTime(day, Time(t-1, 0))
            push!(orders, SimpleOrder(
                :supply, price, gen.p_max,
                Symbol(gen.bidding_zone), dt, 60
            ))
        end
    end
end
\end{minted}
\end{samepage}

\subsection{Στρατηγική Ρεαλιστικής Διαθεσιμότητας}
\label{subsec:all-available-realistic}

Η στρατηγική \texttt{:all\_available\_realistic} συνδυάζει πληροφορίες από τη λύση UC με τη διαθεσιμότητα όλων των μονάδων:

\begin{samepage}
\begin{minted}[fontsize=\small,linenos,frame=lines]{julia}
function apply_all_available_realistic_strategy!(orders, uc_solution, 
                                                  generators, day, markup_factor)
    for (i, gen) in enumerate(generators)
        for t in 1:24
            # Καθορισμός ποσότητας βάσει κατάστασης δέσμευσης
            if uc_solution.commitment[i, t] > COMMITMENT_THRESHOLD
                # Δεσμευμένη μονάδα: χρήση πραγματικής παραγωγής
                quantity = uc_solution.generation[i, t]
            elseif gen.p_min < 0.1 * gen.p_max
                # Μη δεσμευμένη με χαμηλό ελάχιστο: 20% της δυναμικότητας
                quantity = DEFAULT_UNCOMMITTED_UNIT_FRACTION * gen.p_max
            else
                # Μη δεσμευμένη: ελάχιστη ισχύς
                quantity = gen.p_min
            end
            
            if quantity > GENERATION_THRESHOLD
                price = gen.marginal_cost * markup_factor
                dt = DateTime(day, Time(t-1, 0))
                push!(orders, SimpleOrder(
                    :supply, price, quantity,
                    Symbol(gen.bidding_zone), dt, 60
                ))
            end
        end
    end
end
\end{minted}
\end{samepage}

\subsection{Δημιουργία Εντολών Ζήτησης}
\label{subsec:demand-orders-creation}

Οι εντολές ζήτησης δημιουργούνται με βάση τα δεδομένα φορτίου, με τιμή ίση με το ανώτατο όριο της αγοράς (price-taking demand):

\begin{samepage}
\begin{minted}[fontsize=\small,linenos,frame=lines]{julia}
function create_demand_orders!(orders, loads, day, bidding_zone, price_cap)
    for load in loads
        if load.value > DEMAND_THRESHOLD
            hour = parse(Int, load.timeslot[end-3:end-2])
            dt = DateTime(day, Time(hour, 0))
            push!(orders, SimpleOrder(
                :demand, price_cap, load.value,
                Symbol(bidding_zone), dt, 60
            ))
        end
    end
end
\end{minted}
\end{samepage}


% =============================================================================
\section{Εκκαθάριση Αγοράς με MPCC}
\label{sec:mpcc-market-clearing}

Η εκκαθάριση της αγοράς υλοποιείται στο module \texttt{MPCC.jl} χρησιμοποιώντας τη μέθοδο Mathematical Programming with Complementarity Constraints. Ο αλγόριθμος δέχεται ένα βιβλίο εντολών (\texttt{MPCCOrderBook}) και επιστρέφει τις τιμές εκκαθάρισης και τους βαθμούς αποδοχής των εντολών.

\subsection{Δομή Βιβλίου Εντολών}
\label{subsec:order-book-structure}

Το βιβλίο εντολών οργανώνει τις εντολές αγοράς μαζί με τις παραμέτρους του προβλήματος:

\begin{samepage}
\begin{minted}[fontsize=\small,linenos,frame=lines]{julia}
struct MPCCOrderBook
    orders::Vector{MarketOrder}              # Όλες οι εντολές
    nodes::Vector{String}                    # Ζώνες προσφοράς
    periods::Vector{String}                  # Χρονικές περίοδοι
    price_limits::Tuple{Float64,Float64}     # (min, max) τιμή
    network_topology::Union{Nothing,TransferCapacity}  # Περιορισμοί δικτύου
end
\end{minted}
\end{samepage}

\subsection{Μεταβλητές Απόφασης}
\label{subsec:mpcc-variables}

Οι κύριες μεταβλητές του μοντέλου MPCC είναι οι τιμές αγοράς, οι βαθμοί αποδοχής και οι ροές μεταφοράς:

\begin{samepage}
\begin{minted}[fontsize=\small,linenos,frame=lines]{julia}
function solve_mpcc_market_clearing(order_book::MPCCOrderBook; 
                                    big_m::Float64=BIG_M_PARAMETER)
    model = Model(HiGHS.Optimizer)
    set_silent(model)
    
    # Τιμές αγοράς ανά ζώνη και περίοδο
    @variable(model, market_price[node in order_book.nodes, 
                                  period in order_book.periods])
    
    # Βαθμός αποδοχής κάθε εντολής [0, 1]
    @variable(model, 0 <= stepwise_acceptance[order_id in order_ids] <= 1)
    
    # Dual μεταβλητή για περιορισμούς συμπληρωματικότητας
    @variable(model, stepwise_dual[order_id in order_ids] >= 0)
    
    # Μεταβλητή έκτακτης αποκοπής φορτίου
    @variable(model, load_shed[node in order_book.nodes, 
                               period in order_book.periods] >= 0)
    
    # Ροές μεταφοράς (αν υπάρχει δίκτυο)
    if order_book.network_topology !== nothing
        zone_pairs = get_zone_pairs(order_book.network_topology)
        @variable(model, flow[(s,d) in zone_pairs, t in order_book.periods])
    end
end
\end{minted}
\end{samepage}

\subsection{Περιορισμοί Συμπληρωματικότητας}
\label{subsec:complementarity-constraints}

Οι περιορισμοί συμπληρωματικότητας εξασφαλίζουν ότι μια εντολή γίνεται αποδεκτή μόνο όταν η τιμή αγοράς είναι ευνοϊκή. Υλοποιούνται με τη μέθοδο Big-M:

\begin{samepage}
\begin{minted}[fontsize=\small,linenos,frame=lines]{julia}
# Υπολογισμός του δεξιού μέλους της dual συνθήκης
for (i, order) in enumerate(simple_orders)
    order_id = order_ids[i]
    node_id = string(order.zone)
    time_period = extract_time_period(order.date_time, order_book.periods)
    
    # Ποσότητα με πρόσημο: αρνητική για πώληση, θετική για αγορά
    quantity = order.type == :supply ? -order.quantity : order.quantity
    
    stepwise_dual_rhs[order_id] = @expression(model,
        stepwise_dual[order_id] +
        quantity * market_price[node_id, time_period] -
        quantity * order.price
    )
end

# Περιορισμός dual
@constraint(model, [order_id in order_ids],
    0 <= stepwise_dual_rhs[order_id])

# Big-M μετασχηματισμός για συμπληρωματικότητα
@variable(model, complementarity_aux[order_id in order_ids], Bin)

@constraint(model, [order_id in order_ids],
    stepwise_acceptance[order_id] <= complementarity_aux[order_id] * big_m)

@constraint(model, [order_id in order_ids],
    stepwise_dual_rhs[order_id] <= (1 - complementarity_aux[order_id]) * big_m)
\end{minted}
\end{samepage}

\subsection{Περιορισμός Ισορροπίας Ισχύος}
\label{subsec:mpcc-power-balance}

Ο περιορισμός ισορροπίας ισχύος διαφοροποιείται ανάλογα με το αν υπάρχουν περιορισμοί δικτύου:

\begin{samepage}
\begin{minted}[fontsize=\small,linenos,frame=lines]{julia}
if has_network_constraints
    # Πολυζωνική ισορροπία με ροές μεταφοράς
    @constraint(model, [node_id in nodes, period in periods],
        # Συνεισφορά εντολών: πώληση (αρνητική) + αγορά (θετική)
        sum(stepwise_acceptance[order_ids[i]] * 
            (orders[i].type == :supply ? -orders[i].quantity : orders[i].quantity)
            for i in orders_at_node_period[(node_id, period)]) +
        # Εισροές από άλλες ζώνες
        sum(flow[(z, node_id), period] for z in connected_zones_to[node_id]) -
        # Εκροές προς άλλες ζώνες
        sum(flow[(node_id, z), period] for z in connected_zones_from[node_id]) +
        # Έκτακτη αποκοπή
        load_shed[node_id, period] == 0
    )
else
    # Μονοζωνική ισορροπία (χωρίς ροές)
    @constraint(model, [node_id in nodes, period in periods],
        sum(stepwise_acceptance[order_ids[i]] * 
            (orders[i].type == :supply ? -orders[i].quantity : orders[i].quantity)
            for i in orders_at_node_period[(node_id, period)]) +
        load_shed[node_id, period] == 0
    )
end
\end{minted}
\end{samepage}

\subsection{Αντικειμενική Συνάρτηση}
\label{subsec:mpcc-objective}

Η αντικειμενική συνάρτηση μεγιστοποιεί το οικονομικό πλεόνασμα, με ποινή για την αποκοπή φορτίου:

\begin{samepage}
\begin{minted}[fontsize=\small,linenos,frame=lines]{julia}
load_shed_penalty = 10000.0  # €/MWh

@objective(model, Max,
    # Πλεόνασμα από εντολές
    sum(stepwise_acceptance[order_ids[i]] *
        (order.type == :supply ? 
            -order.quantity * order.price :  # Κόστος παραγωγής
             order.quantity * order.price)   # Αξία κατανάλωσης
        for (i, order) in enumerate(simple_orders)) -
    # Ποινή αποκοπής φορτίου
    sum(load_shed_penalty * load_shed[n, p] 
        for n in order_book.nodes, p in order_book.periods)
)
\end{minted}
\end{samepage}

\subsection{Δομή Αποτελέσματος}
\label{subsec:mpcc-result}

Τα αποτελέσματα της εκκαθάρισης επιστρέφονται στη δομή \texttt{MPCCResult}:

\begin{samepage}
\begin{minted}[fontsize=\small,linenos,frame=lines]{julia}
struct MPCCResult
    status::Symbol                                # :optimal, :infeasible, κλπ
    objective_value::Float64                      # Τιμή αντικειμενικής
    market_prices::Dict{String,Dict{String,Float64}}  # Τιμές [zone][period]
    stepwise_acceptance::Dict{String,Float64}     # Αποδοχές εντολών
    block_acceptance::Dict{String,Float64}        # Αποδοχές block
    block_activation::Dict{String,Float64}        # Ενεργοποιήσεις block
    transmission_flows::Dict{String,Dict{String,Float64}}  # Ροές
    solve_time::Float64                           # Χρόνος επίλυσης
    total_time::Float64                           # Συνολικός χρόνος
    solver_name::String                           # Όνομα επιλυτή
    message::String                               # Μήνυμα
end
\end{minted}
\end{samepage}


% =============================================================================
\section{Μοντελοποίηση Δικτύου}
\label{sec:network-modeling}

Η μοντελοποίηση των περιορισμών δικτύου υλοποιείται στο module \texttt{Network.jl}. Το module παρέχει δύο προσεγγίσεις: τη βασισμένη σε γραμμές μεταφοράς (\texttt{NetworkTopology}) και τη βασισμένη σε ζώνες προσφοράς (\texttt{TransferCapacity}). Η δεύτερη προσέγγιση είναι η συνιστώμενη για χρήση με δεδομένα ENTSO-E.

\subsection{Δομή Διαθέσιμης Ικανότητας Μεταφοράς}
\label{subsec:transfer-capacity-struct}

Η δομή \texttt{TransferCapacity} αναπαριστά την ικανότητα μεταφοράς μεταξύ ζωνών προσφοράς:

\begin{samepage}
\begin{minted}[fontsize=\small,linenos,frame=lines]{julia}
struct TransferCapacity
    bidding_zones::Vector{String}    # Όλες οι ζώνες
    time_periods::Vector{String}     # Χρονικές περίοδοι
    # Ικανότητα στην κατεύθυνση source → sink
    capacity_forward::Dict{Tuple{String,String,String},Float64}
    # Ικανότητα στην κατεύθυνση sink → source
    capacity_backward::Dict{Tuple{String,String,String},Float64}
end
\end{minted}
\end{samepage}

\subsection{Ανάκτηση Δεδομένων ENTSO-E}
\label{subsec:entsoe-data-retrieval}

Η δημιουργία της δομής \texttt{TransferCapacity} από δεδομένα ENTSO-E γίνεται με την ακόλουθη συνάρτηση:

\begin{samepage}
\begin{minted}[fontsize=\small,linenos,frame=lines]{julia}
function create_transfer_capacity_from_entsoe(date::Date, 
                                              bidding_zones::Vector{String}=String[])
    zone_filter = isempty(bidding_zones) ? "" :
        "AND (out_map_code IN ('$(join(bidding_zones, "','"))') " *
        "OR in_map_code IN ('$(join(bidding_zones, "','"))'))"
    
    query = """
    SELECT
        out_map_code as source_zone,
        in_map_code as sink_zone,
        EXTRACT(HOUR FROM date_time_utc) + 1 as time_period,
        capacity_mw as capacity
    FROM entsoe.offered_transfer_capacities_implicit
    WHERE DATE(date_time_utc) = \$1
    $zone_filter
    ORDER BY out_map_code, in_map_code, date_time_utc
    """
    
    df = sql2df(query, [date])
    
    if nrow(df) == 0
        error("No transfer capacity data found for $date")
    end
    
    return build_transfer_capacity_from_dataframe(df)
end
\end{minted}
\end{samepage}

\subsection{Προσθήκη Περιορισμών στο Μοντέλο}
\label{subsec:adding-network-constraints}

Οι περιορισμοί ικανότητας μεταφοράς προστίθενται στο μοντέλο JuMP ως εξής:

\begin{samepage}
\begin{minted}[fontsize=\small,linenos,frame=lines]{julia}
function add_transfer_capacity_constraints!(model::Model, 
                                            transfer_capacity::TransferCapacity, 
                                            FLOW)
    zones = transfer_capacity.bidding_zones
    periods = transfer_capacity.time_periods
    cap_fwd = transfer_capacity.capacity_forward
    cap_bwd = transfer_capacity.capacity_backward
    
    # Περιορισμοί για κάθε ζεύγος ζωνών και περίοδο
    # Θετική ροή: source → sink, περιορίζεται από capacity_forward
    # Αρνητική ροή: sink → source, περιορίζεται από capacity_backward
    @constraint(model, 
        [source in zones, sink in zones, t in periods; source != sink],
        -get(cap_bwd, (source, sink, t), 0.0) <= 
            FLOW[source, sink, t] <= 
            get(cap_fwd, (source, sink, t), 0.0)
    )
    
    return model
end
\end{minted}
\end{samepage}

\subsection{Βοηθητικές Συναρτήσεις}
\label{subsec:network-utilities}

Το module παρέχει βοηθητικές συναρτήσεις για την εξαγωγή πληροφοριών συνδεσιμότητας:

\begin{samepage}
\begin{minted}[fontsize=\small,linenos,frame=lines]{julia}
# Επιστροφή ζευγών ζωνών με διαθέσιμη ικανότητα μεταφοράς
function get_zone_pairs(tc::TransferCapacity)
    pairs = Set{Tuple{String,String}}()
    for (source, sink, _) in keys(tc.capacity_forward)
        if get(tc.capacity_forward, (source, sink, "1"), 0.0) > 0
            push!(pairs, (source, sink))
        end
    end
    return collect(pairs)
end

# Επιστροφή συνδεδεμένων ζωνών για μια δεδομένη ζώνη
function get_connected_zones(tc::TransferCapacity, zone::String)
    connected = Set{String}()
    for (source, sink, _) in keys(tc.capacity_forward)
        if source == zone
            push!(connected, sink)
        elseif sink == zone
            push!(connected, source)
        end
    end
    return collect(connected)
end
\end{minted}
\end{samepage}


% =============================================================================
\section{Ολοκλήρωση και Κύρια Ροή Εργασιών}
\label{sec:integration}

Η ολοκλήρωση όλων των συνιστωσών επιτυγχάνεται μέσω δύο κύριων συναρτήσεων: η \texttt{generate\_energy\_prices} για μονοζωνική εκκαθάριση και η \texttt{run\_multi\_zone\_market\_clearing} για πολυζωνική εκκαθάριση με περιορισμούς διασυνοριακής μεταφοράς.

\subsection{Μονοζωνική Εκκαθάριση}
\label{subsec:single-zone-clearing}

Η συνάρτηση \texttt{generate\_energy\_prices} υλοποιεί την πλήρη ροή εργασιών για μια ζώνη προσφοράς:

\begin{samepage}
\begin{minted}[fontsize=\small,linenos,frame=lines]{julia}
function generate_energy_prices(bidding_zone::String, date::Date;
                                order_method::Symbol=:uc_based,
                                markup_factor::Float64=1.1,
                                save_to_db::Bool=true)
    
    # Βήμα 1: Δημιουργία βιβλίου εντολών
    if order_method == :uc_based
        # Επίλυση UC και μετατροπή σε εντολές
        uc_result = solve_unit_commitment(bidding_zone, date)
        bids_result = generate_market_orders_from_uc(
            bidding_zone, date;
            markup_factor=markup_factor,
            bidding_strategy=:committed_only
        )
        order_book = create_typed_order_book(bids_result, bidding_zone, date)
    elseif order_method == :alternative
        # Εναλλακτική μέθοδος (ταχύτερη)
        result = create_adjusted_order_book(bidding_zone, date)
        order_book = result.order_book
    end
    
    # Βήμα 2: Εκκαθάριση αγοράς
    mpcc_result = solve_mpcc_market_clearing(order_book)
    
    # Βήμα 3: Αποθήκευση αποτελεσμάτων
    if save_to_db && mpcc_result.status == :optimal
        save_energy_prices(mpcc_result, bidding_zone, date)
    end
    
    return mpcc_result
end
\end{minted}
\end{samepage}

\subsection{Πολυζωνική Εκκαθάριση}
\label{subsec:multi-zone-clearing}

Η συνάρτηση \texttt{run\_multi\_zone\_market\_clearing} εκτελεί ταυτόχρονη εκκαθάριση πολλαπλών ζωνών με περιορισμούς ATC:

\begin{samepage}
\begin{minted}[fontsize=\small,linenos,frame=lines]{julia}
function run_multi_zone_market_clearing(date::Date;
                                        zones::Vector{String}=String[],
                                        order_method::Symbol=:alternative,
                                        save_to_db::Bool=true)
    
    # Ανακάλυψη διαθέσιμων ζωνών αν δεν καθορίστηκαν
    if isempty(zones)
        zones = get_zones_with_transfer_capacity(date)
    end
    
    println("🌍 Running multi-zone clearing for $(length(zones)) zones")
    
    # Δημιουργία πολυζωνικού βιβλίου εντολών
    order_book = create_multi_zone_order_book(date, zones; 
                                              order_method=order_method)
    
    # Εκκαθάριση αγοράς
    result = solve_mpcc_market_clearing(order_book)
    
    # Αποθήκευση αποτελεσμάτων
    if save_to_db && result.status == :optimal
        for zone in zones
            save_energy_prices(result, zone, date)
        end
        save_transmission_flows(result, date)
    end
    
    # Εκτύπωση αποτελεσμάτων
    println("💰 Zonal Clearing Prices (average):")
    for zone in zones
        prices = values(result.market_prices[zone])
        avg = sum(prices) / length(prices)
        println("   $zone: €$(round(avg, digits=2))/MWh")
    end
    
    return result
end
\end{minted}
\end{samepage}

\paragraph{Παράλληλη Εκτέλεση UC}
Κατά την πολυζωνική εκκαθάριση με τη μέθοδο \texttt{:uc\_based}, η επίλυση Unit Commitment εκτελείται ανεξάρτητα για κάθε ζώνη. Η παράμετρος \texttt{parallel=true} ενεργοποιεί την ταυτόχρονη εκτέλεση αυτών των ανεξάρτητων υπολογισμών μέσω της συνάρτησης \texttt{pmap} του module \texttt{Distributed.jl}. Κάθε ζώνη ανατίθεται σε έναν ξεχωριστό εργάτη (worker process), ο οποίος εκτελεί πλήρως τον κύκλο ανάκτησης δεδομένων, επίλυσης UC και μετατροπής σε εντολές αγοράς. Η μέθοδος \texttt{pmap} εξασφαλίζει δυναμική κατανομή εργασιών: μόλις ένας εργάτης ολοκληρώσει μια ζώνη, αναλαμβάνει την επόμενη διαθέσιμη.

Σε περίπτωση που δεν υπάρχουν διαθέσιμοι εργάτες (\texttt{nworkers() == 0}), το σύστημα υποβαθμίζεται αυτόματα σε σειριακή εκτέλεση χωρίς σφάλμα. Η ασφάλεια νημάτων (thread safety) εξασφαλίζεται από το γεγονός ότι κάθε εργάτης λειτουργεί σε ξεχωριστό address space και τα κλειδιά κρυφής μνήμης είναι ειδικά ανά ζώνη, ενώ οι εγγραφές στη βάση δεδομένων χρησιμοποιούν συναλλαγές PostgreSQL. Η εκκαθάριση MPCC εκτελείται πάντα σειριακά μετά την ολοκλήρωση όλων των UC, καθώς απαιτεί το ενοποιημένο βιβλίο εντολών όλων των ζωνών.

\subsection{Αποθήκευση Αποτελεσμάτων}
\label{subsec:saving-results}

Τα αποτελέσματα αποθηκεύονται στη βάση δεδομένων μέσω των συναρτήσεων \texttt{save\_energy\_prices} και \texttt{save\_transmission\_flows}:

\begin{samepage}
\begin{minted}[fontsize=\small,linenos,frame=lines]{julia}
function save_energy_prices(result::MPCCResult, bidding_zone::String, date::Date)
    # Εξασφάλιση ύπαρξης πίνακα
    ensure_energy_prices_table()
    
    # Προετοιμασία δεδομένων
    prices = result.market_prices[bidding_zone]
    
    withdb() do conn
        for (period, price) in prices
            hour = parse(Int, period)
            execute(conn, """
                INSERT INTO simulations.energy_prices 
                    (bidding_zone, date, hour, price, created_at)
                VALUES (\$1, \$2, \$3, \$4, NOW())
                ON CONFLICT (bidding_zone, date, hour) 
                DO UPDATE SET price = EXCLUDED.price, created_at = NOW()
            """, [bidding_zone, date, hour, price])
        end
    end
end
\end{minted}
\end{samepage}

\subsection{Επεξεργασία Κατά Παρτίδες}
\label{subsec:batch-processing}

Για την επεξεργασία πολλαπλών ημερομηνιών, παρέχονται οι συναρτήσεις \texttt{generate\_energy\_prices\_for\_date\_range} και \texttt{run\_multi\_zone\_for\_date\_range}:

\begin{samepage}
\begin{minted}[fontsize=\small,linenos,frame=lines]{julia}
function run_multi_zone_for_date_range(start_date::Date, end_date::Date;
                                       zones::Vector{String}=String[],
                                       order_method::Symbol=:alternative,
                                       force_rerun::Bool=false,
                                       parallel::Bool=false)
    results = Dict{Date,MPCCResult}()

    for date in start_date:Day(1):end_date
        try
            result = run_multi_zone_market_clearing(date;
                zones=zones, order_method=order_method,
                force_rerun=force_rerun, parallel=parallel)
            results[date] = result
        catch e
            @warn "Failed to process $date: $e"
        end
    end
    return results
end
\end{minted}
\end{samepage}

\subsection{Παράδειγμα Χρήσης}
\label{subsec:usage-example}

Ένα πλήρες παράδειγμα χρήσης του συστήματος για τον υπολογισμό τιμών στην Ελλάδα:

\begin{samepage}
\begin{minted}[fontsize=\small,linenos,frame=lines]{julia}
using Euphemia
using Dates

# Μονοζωνική εκκαθάριση για την Ελλάδα
prices = generate_energy_prices("GR", Date(2024, 6, 15);
    order_method = :uc_based,
    markup_factor = 1.1,
    save_to_db = true
)

# Πολυζωνική εκκαθάριση Ελλάδας-Βουλγαρίας-Ιταλίας
result = run_multi_zone_market_clearing(Date(2024, 6, 15);
    zones = ["GR", "BG", "IT"],
    order_method = :alternative,
    save_to_db = true
)

# Εκτύπωση διασυνοριακών ροών
for (flow_id, flows) in result.transmission_flows
    avg_flow = sum(values(flows)) / length(flows)
    println("$flow_id: $(round(avg_flow, digits=1)) MW average")
end
\end{minted}
\end{samepage}


% =============================================================================
\section{Αυτοματοποιημένη Εκτέλεση}
\label{sec:automated-execution}

Για τη συστηματική παραγωγή τιμών σε ημερήσια βάση, αναπτύχθηκε υποδομή αυτοματοποίησης βασισμένη σε GitHub Actions. Δύο ροές εργασιών (workflows) εκτελούνται σε αυτο-φιλοξενούμενο μηχάνημα (self-hosted runner):

\begin{itemize}
\item \texttt{generate-energy-prices.yml}: Μονοζωνική εκκαθάριση, προγραμματισμένη στις 02:00 UTC ημερησίως.
\item \texttt{generate-multi-zone-prices.yml}: Πολυζωνική εκκαθάριση με διασυνοριακούς περιορισμούς, προγραμματισμένη στις 03:00 UTC ημερησίως.
\end{itemize}

Η πολυζωνική ροή εργασιών εκτελεί τα εξής βήματα: εγκατάσταση εξαρτήσεων Julia, ρύθμιση σύνδεσης βάσης δεδομένων, εκκίνηση παράλληλων εργατών (workers), εκτέλεση του σεναρίου \texttt{bin/multi\_zone\_main.jl} με τις κατάλληλες παραμέτρους, και αποστολή ειδοποιήσεων μέσω Slack σε περίπτωση επιτυχίας ή αποτυχίας. Κατά τη χρήση του εμπορικού επιλυτή Gurobi, ο αριθμός παράλληλων εργατών περιορίζεται αυτόματα σε 2, λόγω του ορίου ταυτόχρονων συνεδριών της άδειας WLS (Web License Service). Κάθε ροή εργασιών υποστηρίζει τόσο χειροκίνητη ενεργοποίηση (\texttt{workflow\_dispatch}) με παραμετροποιήσιμο εύρος ημερομηνιών, όσο και αυτόματη εκτέλεση μέσω χρονοπρογραμματισμού (\texttt{cron}).

% =============================================================================

\chapter{Υλοποίηση}
\label{ch:implementation}

Στο παρόν κεφάλαιο παρουσιάζεται η υλοποίηση του συστήματος προσομοίωσης της Προημερήσιας Αγοράς ηλεκτρικής ενέργειας. Η ανάλυση ακολουθεί τη ροή επεξεργασίας δεδομένων από τα μοντέλα οντοτήτων, μέσω του επιλυτή Unit Commitment και της μηχανής στρατηγικής προσφορών, έως την εκκαθάριση της αγοράς με τη μέθοδο MPCC. Για κάθε συνιστώσα παρέχονται αντιπροσωπευτικά αποσπάσματα κώδικα που επιδεικνύουν τις βασικές αλγοριθμικές επιλογές. Ο πλήρης κώδικας είναι διαθέσιμος στο αποθετήριο του έργου\footnote{\url{https://github.com/silentech-inc/energy-markets}}.


% =============================================================================
\section{Μοντελοποίηση Οντοτήτων}
\label{sec:entity-modeling}

Η υλοποίηση βασίζεται σε σαφώς ορισμένες δομές δεδομένων για την αναπαράσταση των οντοτήτων του συστήματος ηλεκτρικής ενέργειας. Οι δομές αυτές ορίζονται ως αμετάβλητα \texttt{struct} στη Julia, εξασφαλίζοντας ακεραιότητα δεδομένων και αποτελεσματικότητα.

\subsection{Μονάδα Παραγωγής (Generator)}
\label{subsec:generator-struct}

Η δομή \texttt{Generator} αναπαριστά μια θερμική μονάδα παραγωγής με όλα τα τεχνικά και οικονομικά χαρακτηριστικά που απαιτούνται για το πρόβλημα Unit Commitment:

\begin{samepage}
\begin{minted}[fontsize=\small,linenos,frame=lines]{julia}
struct Generator
    code::String           # Μοναδικός κωδικός (EIC code)
    name::String           # Όνομα μονάδας
    fuel_type::Symbol      # Τύπος καυσίμου (:lignite, :gas, :hydro, κλπ)
    location::String       # Τοποθεσία
    p_max::Float64         # Μέγιστη ισχύς (MW)
    p_min::Float64         # Ελάχιστη ισχύς (MW)
    bidding_zone::String   # Ζώνη προσφοράς (π.χ. "GR")
    marginal_cost::Float64 # Οριακό κόστος (€/MWh)
    # Προαιρετικές παράμετροι από εξαγωγή ιστορικών δεδομένων
    # nothing = χρήση προεπιλεγμένων τιμών ανά τύπο καυσίμου
    ramp_up::Union{Float64, Nothing}    # Ρυθμός αύξησης (κλάσμα p_max/ώρα)
    ramp_down::Union{Float64, Nothing}  # Ρυθμός μείωσης (κλάσμα p_max/ώρα)
    min_uptime::Union{Int, Nothing}     # Ελάχιστος χρόνος λειτουργίας (ώρες)
    min_downtime::Union{Int, Nothing}   # Ελάχιστος χρόνος αδράνειας (ώρες)

    # Κατασκευαστής με προεπιλεγμένες τιμές nothing
    function Generator(code, name, fuel_type, location, p_max, p_min,
                       bidding_zone, marginal_cost,
                       ramp_up=nothing, ramp_down=nothing,
                       min_uptime=nothing, min_downtime=nothing)
        new(code, name, fuel_type, location, p_max, p_min,
            bidding_zone, marginal_cost,
            ramp_up, ramp_down, min_uptime, min_downtime)
    end
end
\end{minted}
\end{samepage}

Ο τύπος καυσίμου (\texttt{fuel\_type}) αποθηκεύεται ως \texttt{Symbol} για αποτελεσματική σύγκριση και χρησιμοποιείται για την ανάκτηση παραμέτρων ειδικών ανά τύπο καυσίμου. Τα πεδία \texttt{ramp\_up}, \texttt{ramp\_down}, \texttt{min\_uptime} και \texttt{min\_downtime} ορίζονται ως \texttt{Union\{T, Nothing\}}, γεγονός που επιτρέπει τη χρήση της τιμής \texttt{nothing} ως σημαίας (sentinel) για την ενεργοποίηση των προεπιλεγμένων τιμών ανά τύπο καυσίμου. Ο εσωτερικός κατασκευαστής (inner constructor) ορίζει αυτά τα πεδία με προεπιλεγμένη τιμή \texttt{nothing}, διατηρώντας τη συμβατότητα με τον υπάρχοντα κώδικα. Οι τεχνικές παράμετροι εξάγονται αυτόματα από ιστορικά δεδομένα παραγωγής, όπως περιγράφεται στην Ενότητα~\ref{subsec:parameter-inference}.

Η ανάκτηση μονάδων παραγωγής από τη βάση δεδομένων πραγματοποιείται με τη συνάρτηση \texttt{get\_generators}, η οποία υποστηρίζει φιλτράρισμα μη διαθεσιμοτήτων:

\begin{samepage}
\begin{minted}[fontsize=\small,linenos,frame=lines]{julia}
function get_generators(bidding_zone::String, day::Date; 
                        exclude_unavailable::Bool=true)
    # Ερώτημα SQL με LEFT JOIN στον πίνακα unavailability
    # για αποκλεισμό ή μείωση ισχύος μονάδων με διακοπές
    
    query = """
    WITH active_outages AS (
        SELECT asset_code, MIN(available_capacity_mw) as available_capacity_mw
        FROM entsoe.unavailability_of_production_and_generation_units
        WHERE status = 'Active'
          AND start_outage_utc::timestamp <= \$2::date + interval '1 day'
          AND end_outage_utc::timestamp >= \$2::date
        GROUP BY asset_code
    )
    SELECT g.*, COALESCE(o.available_capacity_mw, g.generation_unit_installed_capacity_mw) as p_max
    FROM entsoe.production_and_generation_units g
    LEFT JOIN active_outages o ON g.generation_unit_code = o.asset_code
    WHERE g.area_map_code = \$1
      AND g.production_unit_status = 'COMMISSIONED'
      AND (o.asset_code IS NULL OR o.available_capacity_mw > 0)
    """
    # ... επεξεργασία αποτελεσμάτων
end
\end{minted}
\end{samepage}

\subsection{Παράμετροι Τύπου Καυσίμου}
\label{subsec:fuel-type-parameters}

Οι τεχνικοί περιορισμοί των μονάδων παραγωγής διαφέρουν σημαντικά ανάλογα με τον τύπο καυσίμου. Η δομή \texttt{FuelTypeParameters} συγκεντρώνει αυτές τις παραμέτρους:

\begin{samepage}
\begin{minted}[fontsize=\small,linenos,frame=lines]{julia}
struct FuelTypeParameters
    # Χαρακτηριστικά εκκίνησης (σε ΩΡΕΣ)
    cold_startup_time::Int          # Ώρες για ψυχρή εκκίνηση
    warm_startup_time::Int          # Ώρες για θερμή εκκίνηση
    hot_startup_time::Int           # Ώρες για καυτή εκκίνηση
    # Ελάχιστοι χρόνοι λειτουργίας (σε ΩΡΕΣ)
    min_uptime::Int                 # Ελάχιστες ώρες λειτουργίας
    min_downtime::Int               # Ελάχιστες ώρες αδράνειας
    # Ρυθμοί μεταβολής (κλάσμα δυναμικότητας ανά ΩΡΑ)
    ramp_up_rate::Float64           # Ρυθμός αύξησης (κλάσμα/ώρα)
    ramp_down_rate::Float64         # Ρυθμός μείωσης (κλάσμα/ώρα)
    # Λειτουργική ευελιξία
    min_load_factor::Float64        # Ελάχιστο φορτίο (κλάσμα)
    part_load_efficiency::Float64   # Απόδοση στο ελάχιστο φορτίο
    # Κατώφλια θερμοκρασιακής μετάβασης (ΩΡΕΣ εκτός)
    warm_threshold::Int             # Ώρες αδράνειας → θερμή κατάσταση
    cold_threshold::Int             # Ώρες αδράνειας → ψυχρή κατάσταση
    # Οικονομικές παράμετροι
    startup_cost_multiplier::Float64  # Πολλαπλασιαστής κόστους εκκίνησης
    no_load_cost_fraction::Float64    # Κόστος χωρίς φορτίο (κλάσμα)
end
\end{minted}
\end{samepage}

Σημαντικό σχεδιαστικό χαρακτηριστικό είναι ότι όλες οι χρονικές παράμετροι αποθηκεύονται σε \textbf{ώρες} (όχι σε περιόδους). Ο επιλυτής UC μετατρέπει αυτόματα σε περιόδους βάσει της πραγματικής χρονικής ανάλυσης των δεδομένων ($\text{periods\_per\_hour} = 60 / \text{resolution\_minutes}$). Τα πεδία \texttt{warm\_threshold} και \texttt{cold\_threshold} χρησιμοποιούνται τόσο στον προσδιορισμό θερμοκρασιακής κατάστασης εκκίνησης (Ενότητα~\ref{subsec:initial-conditions}) όσο και στους αντίστοιχους περιορισμούς του μοντέλου UC. Τα πεδία \texttt{startup\_cost\_multiplier} και \texttt{no\_load\_cost\_fraction} τροφοδοτούν τον υπολογισμό κόστους εκκίνησης και κόστους χωρίς φορτίο στην αντικειμενική συνάρτηση (Ενότητα~\ref{subsec:uc-balance-objective}).

Η συνάρτηση \texttt{get\_fuel\_type\_parameters} επιστρέφει τις κατάλληλες παραμέτρους για κάθε τύπο καυσίμου. Για παράδειγμα, οι λιγνιτικές μονάδες έχουν μεγαλύτερους χρόνους εκκίνησης και χαμηλότερους ρυθμούς μεταβολής σε σύγκριση με τις μονάδες συνδυασμένου κύκλου φυσικού αερίου.

\subsection{Φορτίο και Ανανεώσιμες Πηγές}
\label{subsec:load-renewables}

Η ζήτηση ηλεκτρικής ενέργειας αναπαρίσταται με τη δομή \texttt{Load}:

\begin{samepage}
\begin{minted}[fontsize=\small,linenos,frame=lines]{julia}
struct Load
    timeslot::String           # Χρονοθυρίδα (π.χ. "20240615-0000")
    resolution_code::String    # Χρονική ανάλυση ("PT15M", "PT60M")
    bidding_zone::String       # Ζώνη προσφοράς
    value::Float64             # Κατανάλωση (MW)
end
\end{minted}
\end{samepage}

Οι προβλέψεις παραγωγής από ΑΠΕ αναπαρίστανται με τη δομή \texttt{RenewablesGenerationForecast}:

\begin{samepage}
\begin{minted}[fontsize=\small,linenos,frame=lines]{julia}
struct RenewablesGenerationForecast
    date_time::String                    # Χρονοσφραγίδα
    resolution_code::String              # Χρονική ανάλυση
    bidding_zone::String                 # Ζώνη προσφοράς
    production_type::String              # "Solar" ή "Wind Onshore/Offshore"
    aggregated_generation_forecast::Float64  # Πρόβλεψη (MW)
end
\end{minted}
\end{samepage}

\subsection{Εξαγωγή Τεχνικών Παραμέτρων από Ιστορικά Δεδομένα}
\label{subsec:parameter-inference}

Οι τεχνικές παράμετροι των μονάδων παραγωγής (ρυθμοί μεταβολής, ελάχιστη ισχύς, ελάχιστοι χρόνοι λειτουργίας/αδράνειας) δεν είναι διαθέσιμες στα δημόσια δεδομένα του ENTSO-E. Για την αντιμετώπιση αυτού του περιορισμού, αναπτύχθηκε ένας μηχανισμός αυτόματης εξαγωγής (inference) αυτών των παραμέτρων από ιστορικά δεδομένα πραγματικής παραγωγής.

\paragraph{Πηγή Δεδομένων}
Τα ιστορικά δεδομένα παραγωγής ανακτώνται από τον πίνακα \texttt{actual\_generation\_output\_per\_generation\_unit} του ENTSO-E, ο οποίος περιέχει τη μετρούμενη παραγωγή κάθε μονάδας ανά 15 λεπτά ή ανά ώρα. Η συνάρτηση \texttt{get\_historical\_generation} ανακτά δεδομένα για παράθυρο 3 μηνών πριν από την επιλεγμένη ημερομηνία.

\paragraph{Εξαγωγή Ρυθμών Μεταβολής}
Οι ρυθμοί μεταβολής (ramp rates) υπολογίζονται από τη συνάρτηση \texttt{infer\_ramp\_rates}:

\begin{samepage}
\begin{minted}[fontsize=\small,linenos,frame=lines]{julia}
function infer_ramp_rates(historical_data::DataFrame, p_max::Float64;
                          percentile::Float64=0.95)
    if nrow(historical_data) < MIN_DATA_POINTS_FOR_RAMP_INFERENCE
        return (ramp_up=nothing, ramp_down=nothing)
    end

    # Ταξινόμηση κατά χρόνο και υπολογισμός μεταβολών
    sort!(historical_data, :date_time_utc)
    gen_values = historical_data.actual_generation_output_mw
    deltas = diff(gen_values)

    # Διαχωρισμός σε αυξήσεις και μειώσεις
    ramp_ups = filter(x -> x > 0, deltas)
    ramp_downs = map(abs, filter(x -> x < 0, deltas))

    # Κανονικοποίηση σε ωριαίο ρυθμό βάσει ανάλυσης δεδομένων
    resolution_code = historical_data.resolution_code[1]
    resolution_minutes = parse_resolution_to_minutes(resolution_code)
    hourly_factor = 60.0 / resolution_minutes

    # 95ο εκατοστημόριο, μετατροπή σε κλάσμα p_max/ώρα
    ramp_up = isempty(ramp_ups) ? nothing :
        (quantile(ramp_ups, percentile) * hourly_factor) / p_max
    ramp_down = isempty(ramp_downs) ? nothing :
        (quantile(ramp_downs, percentile) * hourly_factor) / p_max

    return (ramp_up=ramp_up, ramp_down=ramp_down)
end
\end{minted}
\end{samepage}

Οι ρυθμοί αποθηκεύονται ως κλάσμα της μέγιστης ισχύος ανά ώρα (π.χ. 0.20 σημαίνει ότι η μονάδα μπορεί να μεταβάλει την παραγωγή της κατά 20\% της $P^{\max}$ ανά ώρα), επιτρέποντας την εύκολη μετατροπή σε οποιαδήποτε χρονική ανάλυση.

\paragraph{Εξαγωγή Ελάχιστης Ισχύος}
Η ελάχιστη ισχύς ($P^{\min}$) εξάγεται αναλύοντας τις περιόδους σταθερής λειτουργίας, εξαιρώντας τις μεταβατικές καταστάσεις (εκκίνηση/τερματισμός):

\begin{samepage}
\begin{minted}[fontsize=\small,linenos,frame=lines]{julia}
function infer_p_min(historical_data::DataFrame, p_max::Float64,
                     fuel_type::Symbol; percentile::Float64=0.05)
    if nrow(historical_data) < MIN_DATA_POINTS_FOR_RAMP_INFERENCE
        return nothing
    end

    sort!(historical_data, :date_time_utc)
    gen_values = historical_data.actual_generation_output_mw
    deltas = diff(gen_values)

    # Ανίχνευση μεταβατικών καταστάσεων (|Δ| > 5% της p_max)
    ramp_threshold = 0.05 * p_max

    # Συλλογή σταθερών μη-μηδενικών τιμών
    stable_values = Float64[]
    for i in 2:(length(gen_values) - 1)
        value = gen_values[i]
        if value > 0 && abs(deltas[i-1]) < ramp_threshold &&
                        abs(deltas[i]) < ramp_threshold
            push!(stable_values, value)
        end
    end

    if length(stable_values) < 50
        return nothing  # Ανεπαρκή δεδομένα
    end

    # 5ο εκατοστημόριο της σταθερής λειτουργίας
    inferred_p_min = quantile(stable_values, percentile)

    # Εφαρμογή ορίων ανά τύπο καυσίμου (clamp σε MW)
    (min_frac, max_frac) = P_MIN_BOUNDS[get_p_min_bounds_category(fuel_type)]
    return clamp(inferred_p_min, min_frac * p_max, max_frac * p_max)
end
\end{minted}
\end{samepage}

\paragraph{Όρια ανά Τύπο Καυσίμου}
Για την αποφυγή μη ρεαλιστικών τιμών, εφαρμόζονται όρια (bounds) ανά τύπο καυσίμου. Τα όρια αυτά βασίζονται στα τεχνικά χαρακτηριστικά κάθε τεχνολογίας παραγωγής:

\begin{table}[htbp]
\centering
\caption{Όρια ελάχιστης ισχύος ($P^{\min}$) ανά τύπο καυσίμου}
\label{tab:p-min-bounds}
\begin{tabular}{@{}lcc@{}}
\toprule
\textbf{Τύπος Καυσίμου} & \textbf{Ελάχιστο (\%)} & \textbf{Μέγιστο (\%)} \\
\midrule
Λιγνίτης/Άνθρακας & 45 & 65 \\
Φυσικό αέριο (CCGT) & 35 & 55 \\
Φυσικό αέριο (OCGT) & 20 & 45 \\
Υδροηλεκτρικά/Μπαταρίες & 0 & 0 \\
Προεπιλογή & 30 & 60 \\
\bottomrule
\end{tabular}
\end{table}

Οι ευέλικτοι πόροι (υδροηλεκτρικά, μπαταρίες) λαμβάνουν $P^{\min} = 0$ χωρίς εξαγωγή, καθώς μπορούν να λειτουργήσουν σε οποιοδήποτε επίπεδο ισχύος. Η αντίστοιχη σταθερά \texttt{FLEXIBLE\_FUEL\_TYPES} περιλαμβάνει τους τύπους \texttt{Hydro Pumped Storage}, \texttt{Hydro Run-of-river}, \texttt{Hydro Water Reservoir}, \texttt{Battery} και \texttt{Other}. Αξίζει να σημειωθεί ότι οι προεπιλεγμένες τιμές \texttt{min\_load\_factor} στη δομή \texttt{FuelTypeParameters} ευθυγραμμίστηκαν με τα όρια εξαγωγής (π.χ. λιγνίτης: 0.45, ευέλικτοι πόροι: 0.0).

\paragraph{Εξαγωγή Ελάχιστων Χρόνων Λειτουργίας/Αδράνειας}
Οι ελάχιστοι χρόνοι υπολογίζονται αναλύοντας τους κύκλους λειτουργίας (on/off) της μονάδας:

\begin{samepage}
\begin{minted}[fontsize=\small,linenos,frame=lines]{julia}
function infer_uptime_downtime(historical_data::DataFrame, fuel_type::Symbol;
                               percentile::Float64=0.05)
    if nrow(historical_data) < MIN_DATA_POINTS_FOR_RAMP_INFERENCE
        return (nothing, nothing)
    end

    sort!(historical_data, :date_time_utc)
    gen_values = historical_data.actual_generation_output_mw

    # Υπολογισμός διάρκειας περιόδου σε ώρες
    resolution_minutes = parse_resolution_to_minutes(
        historical_data.resolution_code[1])
    period_hours = resolution_minutes / 60.0

    # Προσδιορισμός κατάστασης (on: > 1 MW, off: = 0)
    is_on = gen_values .> 1.0

    # Εντοπισμός κύκλων λειτουργίας
    on_durations, off_durations = Float64[], Float64[]
    current_state = is_on[1]
    current_duration = 1

    for i in 2:length(is_on)
        if is_on[i] == current_state
            current_duration += 1
        else
            duration_hours = current_duration * period_hours
            if current_state
                push!(on_durations, duration_hours)
            else
                push!(off_durations, duration_hours)
            end
            current_state = is_on[i]
            current_duration = 1
        end
    end

    # Φιλτράρισμα σύντομων glitches (< 2 περιόδων)
    min_dur = 2 * period_hours
    filter!(d -> d >= min_dur, on_durations)
    filter!(d -> d >= min_dur, off_durations)

    # Απαίτηση τουλάχιστον 5 κύκλων για αξιόπιστη εξαγωγή
    # Εφαρμογή ορίων ανά τύπο καυσίμου (βλ. Πίνακα~\ref{tab:uptime-downtime-bounds})
    bounds = UPTIME_DOWNTIME_BOUNDS[get_p_min_bounds_category(fuel_type)]
    min_uptime = length(on_durations) >= 5 ?
        round(Int, clamp(quantile(on_durations, percentile),
                         bounds[1]...)) : nothing
    min_downtime = length(off_durations) >= 5 ?
        round(Int, clamp(quantile(off_durations, percentile),
                         bounds[2]...)) : nothing

    return (min_uptime, min_downtime)
end
\end{minted}
\end{samepage}

Ομοίως με τα όρια $P^{\min}$, εφαρμόζονται όρια ανά τύπο καυσίμου για τους χρόνους λειτουργίας και αδράνειας, όπως παρουσιάζονται στον Πίνακα~\ref{tab:uptime-downtime-bounds}.

\begin{table}[htbp]
\centering
\caption{Όρια ελάχιστων χρόνων λειτουργίας/αδράνειας ανά τύπο καυσίμου (ώρες)}
\label{tab:uptime-downtime-bounds}
\begin{tabular}{@{}lcccc@{}}
\toprule
\textbf{Τύπος Καυσίμου} & \multicolumn{2}{c}{\textbf{Χρόνος Λειτουργίας}} & \multicolumn{2}{c}{\textbf{Χρόνος Αδράνειας}} \\
\cmidrule(lr){2-3} \cmidrule(lr){4-5}
& \textbf{Ελάχ.} & \textbf{Μέγ.} & \textbf{Ελάχ.} & \textbf{Μέγ.} \\
\midrule
Λιγνίτης/Άνθρακας & 8 & 48 & 4 & 24 \\
Φυσικό αέριο (CCGT) & 2 & 12 & 1 & 8 \\
Φυσικό αέριο (OCGT) & 1 & 4 & 1 & 4 \\
Προεπιλογή & 2 & 24 & 1 & 12 \\
\bottomrule
\end{tabular}
\end{table}

Η εφαρμογή ορίων διασφαλίζει ότι οι εξαχθείσες τιμές παραμένουν εντός τεχνικά ρεαλιστικού εύρους. Για μονάδες βάσης (π.χ. λιγνιτικές), η εξαγωγή χρόνων συχνά δεν είναι εφικτή λόγω σπάνιων κύκλων λειτουργίας/αδράνειας, οπότε χρησιμοποιούνται οι προεπιλεγμένες τιμές τύπου καυσίμου.

\paragraph{Ενοποίηση Εξαγωγής Παραμέτρων}
Η κεντρική συνάρτηση \texttt{infer\_parameters\_for\_generators} ενορχηστρώνει την εξαγωγή για ένα σύνολο μονάδων, εφαρμόζοντας διαφορετική λογική ανάλογα με τον τύπο καυσίμου:

\begin{samepage}
\begin{minted}[fontsize=\small,linenos,frame=lines]{julia}
function infer_parameters_for_generators(generators::Vector{Generator},
                                         day::Date)
    updated = Generator[]
    for gen in generators
        historical = get_historical_generation(gen.code, day)
        rates = infer_ramp_rates(historical, gen.p_max)

        if gen.fuel_type in FLEXIBLE_FUEL_TYPES
            # Ευέλικτοι πόροι: p_min = 0, χωρίς uptime/downtime
            push!(updated, Generator(gen.code, gen.name, gen.fuel_type,
                gen.location, gen.p_max, 0.0, gen.bidding_zone,
                gen.marginal_cost, rates.ramp_up, rates.ramp_down,
                nothing, nothing))
        else
            # Θερμικές μονάδες: εξαγωγή όλων των παραμέτρων
            p_min = infer_p_min(historical, gen.p_max, gen.fuel_type)
            (uptime, downtime) = infer_uptime_downtime(historical,
                                                        gen.fuel_type)
            push!(updated, Generator(gen.code, gen.name, gen.fuel_type,
                gen.location, gen.p_max,
                p_min !== nothing ? p_min : gen.p_min,
                gen.bidding_zone, gen.marginal_cost,
                rates.ramp_up, rates.ramp_down, uptime, downtime))
        end
    end
    return updated
end
\end{minted}
\end{samepage}

\paragraph{Αποθήκευση σε Κρυφή Μνήμη (Caching)}
Η εξαγωγή παραμέτρων από ιστορικά δεδομένα είναι υπολογιστικά ακριβή (περίπου 17 λεπτά για 39 μονάδες). Για τη βελτίωση της απόδοσης, οι εξαχθείσες παράμετροι αποθηκεύονται στον πίνακα \texttt{simulations.generator\_inferred\_parameters}:

\begin{samepage}
\begin{minted}[fontsize=\small,linenos,frame=lines]{julia}
# Ανάκτηση μονάδων με αυτόματη εξαγωγή και caching
generators = get_generators_with_inferred_params("GR", Date(2024, 6, 15);
    use_cache = true,       # Χρήση cached τιμών αν διαθέσιμες
    max_age_days = 7        # Ανανέωση cache κάθε 7 ημέρες
)
\end{minted}
\end{samepage}

Η χρήση της κρυφής μνήμης μειώνει τον χρόνο ανάκτησης από περίπου 17 λεπτά σε περίπου 2 δευτερόλεπτα---βελτίωση κατά 525 φορές.


\subsection{Αρχικές Συνθήκες Μονάδων Παραγωγής}
\label{subsec:initial-conditions}

Για τη ρεαλιστική μοντελοποίηση του προβλήματος Unit Commitment, η κατάσταση κάθε μονάδας κατά τη χρονική στιγμή $t=0$ (έναρξη του ορίζοντα βελτιστοποίησης) πρέπει να είναι γνωστή. Οι αρχικές συνθήκες καθορίζουν αν η μονάδα λειτουργούσε, σε ποιο επίπεδο ισχύος, και για πόσο χρόνο βρισκόταν στην τρέχουσα κατάστασή της. Αυτή η πληροφορία είναι απαραίτητη για τη σωστή επιβολή περιορισμών στις πρώτες περιόδους του ορίζοντα.

\paragraph{Δομή Αρχικών Συνθηκών}
Η δομή \texttt{InitialConditions} αναπαριστά την κατάσταση μιας μονάδας στο $t=0$:

\begin{samepage}
\begin{minted}[fontsize=\small,linenos,frame=lines]{julia}
struct InitialConditions
    is_on::Bool           # Λειτουργούσε η μονάδα στο t=0;
    output::Float64       # Ισχύς εξόδου σε MW (0 αν εκτός λειτουργίας)
    hours_on::Int         # Συνεχόμενες ώρες λειτουργίας
    hours_off::Int        # Συνεχόμενες ώρες αδράνειας
    thermal_state::Symbol # Θερμοκρασιακή κατάσταση (:hot, :warm, :cold)
end
\end{minted}
\end{samepage}

Η θερμοκρασιακή κατάσταση προσδιορίζεται από τον χρόνο αδράνειας: \texttt{:hot} αν $T_{\text{off}} \leq 8$ ώρες, \texttt{:warm} αν $8 < T_{\text{off}} \leq 48$ ώρες, και \texttt{:cold} αν $T_{\text{off}} > 48$ ώρες. Τα κατώφλια αυτά προσαρμόζονται ανά τύπο καυσίμου μέσω των αντίστοιχων πεδίων \texttt{warm\_threshold} και \texttt{cold\_threshold} της δομής \texttt{FuelTypeParameters}.

\paragraph{Εξαγωγή από Ιστορικά Δεδομένα}
Η κεντρική συνάρτηση \texttt{get\_initial\_conditions} εξάγει τις αρχικές συνθήκες για ένα σύνολο μονάδων χρησιμοποιώντας ένα μαζικό ερώτημα SQL (batch query) αντί $N$ ξεχωριστών ερωτημάτων. Η τεχνική αυτή επιτυγχάνει βελτίωση ταχύτητας κατά 36 φορές (από $\sim$9 λεπτά σε $\sim$15 δευτερόλεπτα για 40 μονάδες):

\begin{samepage}
\begin{minted}[fontsize=\small,linenos,frame=lines]{julia}
function get_initial_conditions(generators::Vector{Generator},
                                market_day::Date;
                                use_historical::Bool=true)
    # Μαζική ανάκτηση 72 ωρών ιστορικού για όλες τις μονάδες
    # σε ένα ερώτημα SQL (WHERE generation_unit_code = ANY($1))
    day_before = market_day - Day(1)
    start_utc = DateTime(day_before, Time(23, 0, 0))  # ~00:00 CET
    codes = [gen.code for gen in generators]
    all_historical = get_recent_generation_batch(codes, start_utc)

    conditions = Dict{String, InitialConditions}()
    for gen in generators
        gen_data = filter(r -> r.generation_unit_code == gen.code,
                          all_historical)
        ic = infer_initial_conditions_from_data(gen_data, gen.fuel_type)
        conditions[gen.code] = ic
    end
    return conditions
end
\end{minted}
\end{samepage}

Ο αλγόριθμος εξαγωγής αναλύει τα δεδομένα σε φθίνουσα χρονική σειρά: η πιο πρόσφατη μέτρηση $> 1$ MW δηλώνει ότι η μονάδα ήταν σε λειτουργία. Στη συνέχεια, μετρούνται οι συνεχόμενες περίοδοι στην ίδια κατάσταση για τον υπολογισμό του $T_{\text{on}}$ ή $T_{\text{off}}$. Σε περίπτωση απουσίας δεδομένων, χρησιμοποιούνται ευρετικές προεπιλεγμένες τιμές ανά τύπο καυσίμου: οι μονάδες βάσης (λιγνίτης, πυρηνικά) θεωρούνται σε λειτουργία, οι μονάδες μέσου φορτίου (CCGT) εκτός λειτουργίας αλλά θερμές, οι μονάδες αιχμής (OCGT) εκτός λειτουργίας και ψυχρές, και οι ευέλικτοι πόροι (υδροηλεκτρικά, μπαταρίες) εκτός λειτουργίας αλλά έτοιμοι.


\subsection{Τύποι Εντολών Αγοράς}
\label{subsec:market-orders}

Οι εντολές αγοράς υλοποιούνται ως ιεραρχία τύπων με αφηρημένο βασικό τύπο \texttt{MarketOrder}. Ο κύριος τύπος που χρησιμοποιείται είναι η απλή ωριαία εντολή \texttt{SimpleOrder}:

\begin{samepage}
\begin{minted}[fontsize=\small,linenos,frame=lines]{julia}
abstract type MarketOrder end

struct SimpleOrder <: MarketOrder
    type::Symbol           # :supply ή :demand
    price::Float64         # Τιμή (€/MWh)
    quantity::Float64      # Ποσότητα (MWh)
    zone::Symbol           # Ζώνη προσφοράς
    date_time::DateTime    # Χρόνος παράδοσης
    resolution_code::Int   # Ανάλυση σε λεπτά (60, 30, 15)
end
\end{minted}
\end{samepage}

Πέρα από τις απλές εντολές, υλοποιούνται και σύνθετοι τύποι όπως οι εντολές μπλοκ (\texttt{BlockOrder}), οι συνδεδεμένες εντολές μπλοκ (\texttt{LinkedBlockOrder}) και οι αποκλειστικές εντολές μπλοκ (\texttt{ExclusiveBlockOrder}), ακολουθώντας την προδιαγραφή του αλγορίθμου EUPHEMIA.


% =============================================================================
\section{Επιλυτής Unit Commitment}
\label{sec:uc-solver}

Ο επιλυτής Unit Commitment υλοποιείται στο module \texttt{UnitCommitment.jl} και αποτελεί το πρώτο βήμα της αλυσίδας υπολογισμών. Η κύρια συνάρτηση \texttt{solve\_unit\_commitment} δέχεται ως είσοδο τη ζώνη προσφοράς και την ημερομηνία, ανακτά τα απαραίτητα δεδομένα από τη βάση, διαμορφώνει το πρόβλημα βελτιστοποίησης και επιστρέφει τη βέλτιστη λύση.

\subsection{Δημιουργία Μοντέλου}
\label{subsec:uc-model-creation}

Η δημιουργία του μοντέλου JuMP ξεκινά με την επιλογή του επιλυτή και τον ορισμό των μεταβλητών απόφασης:

\begin{samepage}
\begin{minted}[fontsize=\small,linenos,frame=lines]{julia}
function solve_unit_commitment(bidding_zone::String, day::Date;
                               optimizer::String="auto",
                               mip_gap::Float64=0.01,
                               time_limit::Float64=600.0)
    # Φόρτωση δεδομένων
    generators = get_generators(bidding_zone, day)
    loads = get_loads(bidding_zone, day)
    renewables = get_generation_forecast_for_wind_and_solar(bidding_zone, day)
    initial_conditions = get_initial_conditions(generators, day)

    N = length(generators)  # Αριθμός μονάδων
    T = 24                  # Χρονικές περίοδοι
    periods_per_hour = 60 / resolution_minutes

    # Δημιουργία μοντέλου με ρύθμιση παραμέτρων επιλυτή
    model = Model(optimizer_func)
    set_silent(model)
    set_optimizer_attribute(model, MOI.RelativeGapTolerance(), mip_gap)
    set_optimizer_attribute(model, MOI.TimeLimitSec(), time_limit)

    # Εξαγωγή αρχικών συνθηκών σε πίνακες
    u0 = zeros(Int, N)      # Αρχική δέσμευση (0 ή 1)
    g0 = zeros(Float64, N)  # Αρχική παραγωγή (MW)
    T_on0 = zeros(Int, N)   # Ώρες ήδη σε λειτουργία
    T_off0 = zeros(Int, N)  # Ώρες ήδη εκτός λειτουργίας
    for (i, gen) in enumerate(generators)
        ic = initial_conditions[gen.code]
        u0[i] = ic.is_on ? 1 : 0
        g0[i] = ic.output
        T_on0[i] = ic.hours_on
        T_off0[i] = ic.hours_off
    end

    # Μεταβλητές παραγωγής (συνεχείς)
    @variable(model, 0 <= g[i=1:N, t=1:T] <= generators[i].p_max)

    # Μεταβλητές δέσμευσης (δυαδικές)
    @variable(model, u[i=1:N, t=1:T], Bin)  # Κατάσταση λειτουργίας
    @variable(model, v[i=1:N, t=1:T], Bin)  # Εκκίνηση
    @variable(model, z[i=1:N, t=1:T], Bin)  # Τερματισμός

    # Μεταβλητές για θερμοκρασιακές καταστάσεις εκκίνησης
    Θ = [:cold, :warm, :hot]
    @variable(model, v_θ[i=1:N, θ in Θ, t=1:T], Bin)

    # ... περαιτέρω μεταβλητές
end
\end{minted}
\end{samepage}

\subsection{Περιορισμοί Ισχύος}
\label{subsec:impl-uc-power-constraints}

Οι περιορισμοί ισχύος εξασφαλίζουν ότι η παραγωγή κάθε μονάδας βρίσκεται εντός των τεχνικών της ορίων:

\begin{samepage}
\begin{minted}[fontsize=\small,linenos,frame=lines]{julia}
# Ανώτατο όριο παραγωγής
@constraint(model, [i = 1:N, t = 1:T], 
    g[i, t] <= generators[i].p_max * u[i, t])

# Κατώτατο όριο παραγωγής με παραμέτρους ανά τύπο καυσίμου
for (i, gen) in enumerate(generators)
    params = fuel_params[i]
    # Το ελάχιστο είναι το μέγιστο μεταξύ p_min και του ελάχιστου 
    # φορτίου βάσει τύπου καυσίμου
    min_gen = max(gen.p_min, params.min_load_factor * gen.p_max)
    @constraint(model, [t = 1:T], g[i, t] >= min_gen * u[i, t])
end
\end{minted}
\end{samepage}

\subsection{Περιορισμοί Δέσμευσης}
\label{subsec:impl-uc-commitment-constraints}

Η σχέση μεταξύ κατάστασης λειτουργίας, εκκίνησης και τερματισμού υλοποιείται ως εξής, με ενσωμάτωση αρχικών συνθηκών στην περίοδο $t=1$:

\begin{samepage}
\begin{minted}[fontsize=\small,linenos,frame=lines]{julia}
# Σύνδεση t=1 με αρχική κατάσταση u0
if initial_conditions !== nothing
    @constraint(model, [i in 1:N],
        u[i, 1] == u0[i] + v[i, 1] - z[i, 1])
end

# Σύνδεση μεταβλητών δέσμευσης (t ≥ 2)
@constraint(model, [i in 1:N, t in 2:T],
    u[i, t] == u[i, t-1] + v[i, t] - z[i, t])

# Αποκλεισμός ταυτόχρονης εκκίνησης και τερματισμού
@constraint(model, [i = 1:N, t = 1:T], v[i, t] + z[i, t] <= 1)

# Ελάχιστος χρόνος λειτουργίας (ώρες → περίοδοι)
for (i, gen) in enumerate(generators)
    UT_hours = gen.min_uptime !== nothing ?
        gen.min_uptime : fuel_params[i].min_uptime
    UT = max(1, ceil(Int, UT_hours * periods_per_hour))
    if UT > 1
        @constraint(model, [t = UT:T],
            sum(v[i, τ] for τ in t-UT+1:t) <= u[i, t])
        # Αρχική συνθήκη: εναπομένων χρόνος λειτουργίας
        if initial_conditions !== nothing && u0[i] == 1
            remaining = max(0, UT - ceil(Int, T_on0[i] * periods_per_hour))
            for t in 1:min(remaining, T)
                @constraint(model, u[i, t] == 1)
            end
        end
    end
end

# Ελάχιστος χρόνος αδράνειας (ώρες → περίοδοι)
for (i, gen) in enumerate(generators)
    DT_hours = gen.min_downtime !== nothing ?
        gen.min_downtime : fuel_params[i].min_downtime
    DT = max(1, ceil(Int, DT_hours * periods_per_hour))
    if DT > 1
        @constraint(model, [t = DT:T],
            sum(z[i, τ] for τ in t-DT+1:t) <= 1 - u[i, t])
        # Αρχική συνθήκη: εναπομένων χρόνος αδράνειας
        if initial_conditions !== nothing && u0[i] == 0
            remaining = max(0, DT - ceil(Int, T_off0[i] * periods_per_hour))
            for t in 1:min(remaining, T)
                @constraint(model, u[i, t] == 0)
            end
        end
    end
end
\end{minted}
\end{samepage}

Η μετατροπή $\texttt{UT\_hours} \times \texttt{periods\_per\_hour}$ διασφαλίζει ότι οι περιορισμοί εφαρμόζονται σωστά ανεξάρτητα από τη χρονική ανάλυση (15, 30 ή 60 λεπτά). Οι αρχικές συνθήκες επιβάλλουν τη συνέχιση της τρέχουσας κατάστασης κάθε μονάδας: αν μια μονάδα λειτουργούσε αλλά δεν έχει συμπληρώσει τον ελάχιστο χρόνο λειτουργίας, εξαναγκάζεται να παραμείνει σε λειτουργία για τις εναπομένουσες περιόδους.

\subsection{Περιορισμοί Ρυθμού Μεταβολής}
\label{subsec:impl-uc-ramping}

Οι περιορισμοί ρυθμού μεταβολής (ramping) χρησιμοποιούν την τεχνική Big-M για την αποδέσμευση κατά τις μεταβατικές καταστάσεις:

\begin{samepage}
\begin{minted}[fontsize=\small,linenos,frame=lines]{julia}
# Υπολογισμός ρυθμών μεταβολής ανά μονάδα:
# εξαχθείσα τιμή (fraction/hour) ή προεπιλογή τύπου καυσίμου
# Μετατροπή σε MW/περίοδο: rate × p_max × period_hours
period_hours = resolution_minutes / 60.0
ramp_up = Float64[]
ramp_down = Float64[]
for (i, gen) in enumerate(generators)
    if gen.ramp_up !== nothing
        push!(ramp_up, gen.ramp_up * gen.p_max * period_hours)
    else
        push!(ramp_up, fuel_params[i].ramp_up_rate * gen.p_max * period_hours)
    end
    if gen.ramp_down !== nothing
        push!(ramp_down, gen.ramp_down * gen.p_max * period_hours)
    else
        push!(ramp_down, fuel_params[i].ramp_down_rate * gen.p_max * period_hours)
    end
end

# Big-M: κλιμακούμενο βάσει μεγέθους προβλήματος
max_p_max = maximum(gen.p_max for gen in generators)
M = 2.0 * max_p_max

# Περιορισμός αύξησης (t ≥ 2)
@constraint(model, [i in 1:N, t in 2:T],
    g[i, t] - g[i, t-1] <= ramp_up[i] + M * u_SU[i, t])

# Περιορισμός μείωσης (t ≥ 2)
@constraint(model, [i in 1:N, t in 2:T],
    g[i, t-1] - g[i, t] <= ramp_down[i] + M * u_SD[i, t])
\end{minted}
\end{samepage}

\subsection{Ισορροπία Ισχύος και Αντικειμενική Συνάρτηση}
\label{subsec:uc-balance-objective}

Ο περιορισμός ισορροπίας ισχύος εξασφαλίζει την κάλυψη της καθαρής ζήτησης:

\begin{samepage}
\begin{minted}[fontsize=\small,linenos,frame=lines]{julia}
# Περιορισμός ισορροπίας
@constraint(model, [t in 1:T],
    sum(g[i, t] for i in 1:N) == net_demand[t])
\end{minted}
\end{samepage}

Η αντικειμενική συνάρτηση αποτελείται από τρεις συνιστώσες κόστους: κόστος παραγωγής, κόστος εκκίνησης και κόστος χωρίς φορτίο. Πριν τη διαμόρφωση, υπολογίζονται οι παράμετροι κόστους:

\begin{samepage}
\begin{minted}[fontsize=\small,linenos,frame=lines]{julia}
# Κόστος εκκίνησης ανά μονάδα και θερμοκρασιακή κατάσταση (€)
# Βασικό κόστος = startup_cost_multiplier × marginal_cost × p_max
# Πολλαπλασιαστές: hot=1.0, warm=1.5, cold=2.5
C_SU = Dict{Tuple{Int,Symbol},Float64}()
for i in 1:N
    params = fuel_params[i]
    gen = generators[i]
    base = params.startup_cost_multiplier * gen.marginal_cost * gen.p_max
    C_SU[(i, :hot)]  = base * 1.0
    C_SU[(i, :warm)] = base * 1.5
    C_SU[(i, :cold)] = base * 2.5
end

# Κόστος χωρίς φορτίο ανά μονάδα (€/περίοδο)
# Πάγιο κόστος δέσμευσης ανεξάρτητο από παραγωγή
period_hours = resolution_minutes / 60.0
C_NL = Float64[]
for i in 1:N
    params = fuel_params[i]
    gen = generators[i]
    push!(C_NL, params.no_load_cost_fraction *
                gen.marginal_cost * gen.p_min * period_hours)
end

# Αντικειμενική: ελαχιστοποίηση συνολικού κόστους
@objective(model, Min,
    # 1. Κόστος παραγωγής (μεταβλητό)
    sum(generators[i].marginal_cost * g[i, t] for i in 1:N, t in 1:T)
    # 2. Κόστος εκκίνησης (εξαρτώμενο από θερμοκρασία)
    + sum(C_SU[(i, θ)] * v_θ[i, θ, t] for i in 1:N, θ in Θ, t in 1:T)
    # 3. Κόστος χωρίς φορτίο (πάγιο κόστος δέσμευσης)
    + sum(C_NL[i] * u[i, t] for i in 1:N, t in 1:T)
)
\end{minted}
\end{samepage}

Η τρισδιάστατη δομή κόστους αποτυπώνει ρεαλιστικά τα κίνητρα λειτουργίας: το κόστος εκκίνησης αποθαρρύνει τις συχνές εκκινήσεις/τερματισμούς (ιδιαίτερα τις ψυχρές εκκινήσεις, που κοστίζουν 2,5 φορές περισσότερο από τις θερμές), ενώ το κόστος χωρίς φορτίο αντανακλά τα πάγια λειτουργικά κόστη δέσμευσης μιας μονάδας ανεξάρτητα από το επίπεδο παραγωγής της.

\subsection{Επίλυση και Επιστροφή Αποτελεσμάτων}
\label{subsec:uc-solution}

Μετά την επίλυση, τα αποτελέσματα επιστρέφονται σε δομημένη μορφή:

\begin{samepage}
\begin{minted}[fontsize=\small,linenos,frame=lines]{julia}
optimize!(model)

if termination_status(model) == OPTIMAL
    return (
        status = OPTIMAL,
        generators = generators,
        generation = value.(g),       # Πίνακας παραγωγής N×T
        commitment = value.(u),       # Πίνακας δέσμευσης N×T
        startup = value.(v),          # Πίνακας εκκινήσεων N×T
        shutdown = value.(z),         # Πίνακας τερματισμών N×T
        objective_value = objective_value(model),
        solve_time = solve_time
    )
else
    return (status = termination_status(model),)
end
\end{minted}
\end{samepage}

\subsection{Αποθήκευση Αποτελεσμάτων σε Κρυφή Μνήμη}
\label{subsec:uc-caching}

Η επίλυση του προβλήματος Unit Commitment απαιτεί 10--15 λεπτά ανά ζώνη. Για την αποφυγή επαναληπτικών υπολογισμών, τα αποτελέσματα αποθηκεύονται σε τρεις πίνακες PostgreSQL (\texttt{uc\_results}, \texttt{uc\_generation}, \texttt{uc\_net\_demand}) με κλειδί \texttt{(bidding\_zone, market\_date, code\_version)}. Η συνάρτηση \texttt{solve\_unit\_commitment} ελέγχει αυτόματα τη διαθεσιμότητα αποθηκευμένων αποτελεσμάτων:

\begin{samepage}
\begin{minted}[fontsize=\small,linenos,frame=lines]{julia}
function solve_unit_commitment(bidding_zone::String, day::Date;
                               optimizer::String="auto",
                               use_cache::Bool=true,
                               force_rerun::Bool=false,
                               # ... λοιπές παράμετροι
                               )
    # Έλεγχος κρυφής μνήμης (εκτός αν απενεργοποιημένη)
    if use_cache && !force_rerun
        cached = load_uc_results(bidding_zone, day)
        if cached !== nothing
            return cached   # Επιστροφή σε ~2 δευτερόλεπτα
        end
    end

    # ... κανονική επίλυση (10-15 λεπτά) ...

    # Αποθήκευση αποτελεσμάτων (delete-before-insert)
    if use_cache
        save_uc_results(solution, bidding_zone, day)
    end
    return solution
end
\end{minted}
\end{samepage}

Η παράμετρος \texttt{force\_rerun} διαδίδεται μέσω ολόκληρης της αλυσίδας κλήσεων (\texttt{generate\_energy\_prices} $\to$ \texttt{create\_typed\_order\_book} $\to$ \texttt{generate\_market\_orders\_from\_uc} $\to$ \texttt{solve\_unit\_commitment}), επιτρέποντας τον επανυπολογισμό αποτελεσμάτων όταν αλλάξει ο κώδικας ή τα δεδομένα εισόδου. Η αποθήκευση ανέρχεται σε περίπου 195 KB ανά ζώνη/ημέρα---κόστος αμελητέο σε σύγκριση με την εξοικονόμηση χρόνου.


% =============================================================================
\section{Μηχανή Στρατηγικής Προσφορών}
\label{sec:bidding-strategy-engine}

Η μηχανή στρατηγικής προσφορών υλοποιείται στο module \texttt{BiddingStrategy.jl} και μετατρέπει τα αποτελέσματα του Unit Commitment σε εντολές αγοράς. Η κύρια συνάρτηση \texttt{generate\_market\_orders\_from\_uc} υποστηρίζει τρεις διακριτές στρατηγικές.

\subsection{Δομή Αποτελέσματος}
\label{subsec:bidding-result}

Τα αποτελέσματα της μετατροπής επιστρέφονται στη δομή \texttt{UCToBidsResult}:

\begin{samepage}
\begin{minted}[fontsize=\small,linenos,frame=lines]{julia}
struct UCToBidsResult
    supply_orders::Vector{SimpleOrder}   # Εντολές πώλησης
    demand_orders::Vector{SimpleOrder}   # Εντολές αγοράς
    bidding_zone::Symbol                 # Ζώνη προσφοράς
    day::Date                            # Ημερομηνία
    total_supply_quantity::Float64       # Συνολική προσφερόμενη ποσότητα
    total_demand_quantity::Float64       # Συνολική ζητούμενη ποσότητα
    success::Bool                        # Επιτυχία μετατροπής
    message::String                      # Μήνυμα κατάστασης
end
\end{minted}
\end{samepage}

\subsection{Στρατηγική Δεσμευμένων Μονάδων}
\label{subsec:committed-only-strategy}

Η συντηρητική στρατηγική \texttt{:committed\_only} δημιουργεί προσφορές μόνο για τις μονάδες που έχουν δεσμευτεί στη λύση του UC:

\begin{samepage}
\begin{minted}[fontsize=\small,linenos,frame=lines]{julia}
function apply_committed_only_strategy!(orders, uc_solution, generators, 
                                        day, markup_factor)
    for (i, gen) in enumerate(generators)
        for t in 1:24
            # Έλεγχος αν η μονάδα είναι δεσμευμένη
            if uc_solution.commitment[i, t] > COMMITMENT_THRESHOLD
                generation = uc_solution.generation[i, t]
                if generation > GENERATION_THRESHOLD
                    # Δημιουργία εντολής πώλησης
                    price = gen.marginal_cost * markup_factor
                    dt = DateTime(day, Time(t-1, 0))
                    push!(orders, SimpleOrder(
                        :supply, price, generation, 
                        Symbol(gen.bidding_zone), dt, 60
                    ))
                end
            end
        end
    end
end
\end{minted}
\end{samepage}

\subsection{Στρατηγική Μέγιστης Διαθεσιμότητας}
\label{subsec:all-available-max}

Η στρατηγική \texttt{:all\_available\_max} προσφέρει όλες τις μονάδες στη μέγιστη δυναμικότητά τους, ανεξάρτητα από τη λύση του UC. Αν και μη ρεαλιστική, χρησιμεύει ως βάση σύγκρισης:

\begin{samepage}
\begin{minted}[fontsize=\small,linenos,frame=lines]{julia}
function apply_all_available_max_strategy!(orders, generators, day, markup_factor)
    for gen in generators
        for t in 1:24
            price = gen.marginal_cost * markup_factor
            dt = DateTime(day, Time(t-1, 0))
            push!(orders, SimpleOrder(
                :supply, price, gen.p_max,
                Symbol(gen.bidding_zone), dt, 60
            ))
        end
    end
end
\end{minted}
\end{samepage}

\subsection{Στρατηγική Ρεαλιστικής Διαθεσιμότητας}
\label{subsec:all-available-realistic}

Η στρατηγική \texttt{:all\_available\_realistic} συνδυάζει πληροφορίες από τη λύση UC με τη διαθεσιμότητα όλων των μονάδων:

\begin{samepage}
\begin{minted}[fontsize=\small,linenos,frame=lines]{julia}
function apply_all_available_realistic_strategy!(orders, uc_solution, 
                                                  generators, day, markup_factor)
    for (i, gen) in enumerate(generators)
        for t in 1:24
            # Καθορισμός ποσότητας βάσει κατάστασης δέσμευσης
            if uc_solution.commitment[i, t] > COMMITMENT_THRESHOLD
                # Δεσμευμένη μονάδα: χρήση πραγματικής παραγωγής
                quantity = uc_solution.generation[i, t]
            elseif gen.p_min < 0.1 * gen.p_max
                # Μη δεσμευμένη με χαμηλό ελάχιστο: 20% της δυναμικότητας
                quantity = DEFAULT_UNCOMMITTED_UNIT_FRACTION * gen.p_max
            else
                # Μη δεσμευμένη: ελάχιστη ισχύς
                quantity = gen.p_min
            end
            
            if quantity > GENERATION_THRESHOLD
                price = gen.marginal_cost * markup_factor
                dt = DateTime(day, Time(t-1, 0))
                push!(orders, SimpleOrder(
                    :supply, price, quantity,
                    Symbol(gen.bidding_zone), dt, 60
                ))
            end
        end
    end
end
\end{minted}
\end{samepage}

\subsection{Δημιουργία Εντολών Ζήτησης}
\label{subsec:demand-orders-creation}

Οι εντολές ζήτησης δημιουργούνται με βάση τα δεδομένα φορτίου, με τιμή ίση με το ανώτατο όριο της αγοράς (price-taking demand):

\begin{samepage}
\begin{minted}[fontsize=\small,linenos,frame=lines]{julia}
function create_demand_orders!(orders, loads, day, bidding_zone, price_cap)
    for load in loads
        if load.value > DEMAND_THRESHOLD
            hour = parse(Int, load.timeslot[end-3:end-2])
            dt = DateTime(day, Time(hour, 0))
            push!(orders, SimpleOrder(
                :demand, price_cap, load.value,
                Symbol(bidding_zone), dt, 60
            ))
        end
    end
end
\end{minted}
\end{samepage}


% =============================================================================
\section{Εκκαθάριση Αγοράς με MPCC}
\label{sec:mpcc-market-clearing}

Η εκκαθάριση της αγοράς υλοποιείται στο module \texttt{MPCC.jl} χρησιμοποιώντας τη μέθοδο Mathematical Programming with Complementarity Constraints. Ο αλγόριθμος δέχεται ένα βιβλίο εντολών (\texttt{MPCCOrderBook}) και επιστρέφει τις τιμές εκκαθάρισης και τους βαθμούς αποδοχής των εντολών.

\subsection{Δομή Βιβλίου Εντολών}
\label{subsec:order-book-structure}

Το βιβλίο εντολών οργανώνει τις εντολές αγοράς μαζί με τις παραμέτρους του προβλήματος:

\begin{samepage}
\begin{minted}[fontsize=\small,linenos,frame=lines]{julia}
struct MPCCOrderBook
    orders::Vector{MarketOrder}              # Όλες οι εντολές
    nodes::Vector{String}                    # Ζώνες προσφοράς
    periods::Vector{String}                  # Χρονικές περίοδοι
    price_limits::Tuple{Float64,Float64}     # (min, max) τιμή
    network_topology::Union{Nothing,TransferCapacity}  # Περιορισμοί δικτύου
end
\end{minted}
\end{samepage}

\subsection{Μεταβλητές Απόφασης}
\label{subsec:mpcc-variables}

Οι κύριες μεταβλητές του μοντέλου MPCC είναι οι τιμές αγοράς, οι βαθμοί αποδοχής και οι ροές μεταφοράς:

\begin{samepage}
\begin{minted}[fontsize=\small,linenos,frame=lines]{julia}
function solve_mpcc_market_clearing(order_book::MPCCOrderBook; 
                                    big_m::Float64=BIG_M_PARAMETER)
    model = Model(HiGHS.Optimizer)
    set_silent(model)
    
    # Τιμές αγοράς ανά ζώνη και περίοδο
    @variable(model, market_price[node in order_book.nodes, 
                                  period in order_book.periods])
    
    # Βαθμός αποδοχής κάθε εντολής [0, 1]
    @variable(model, 0 <= stepwise_acceptance[order_id in order_ids] <= 1)
    
    # Dual μεταβλητή για περιορισμούς συμπληρωματικότητας
    @variable(model, stepwise_dual[order_id in order_ids] >= 0)
    
    # Μεταβλητή έκτακτης αποκοπής φορτίου
    @variable(model, load_shed[node in order_book.nodes, 
                               period in order_book.periods] >= 0)
    
    # Ροές μεταφοράς (αν υπάρχει δίκτυο)
    if order_book.network_topology !== nothing
        zone_pairs = get_zone_pairs(order_book.network_topology)
        @variable(model, flow[(s,d) in zone_pairs, t in order_book.periods])
    end
end
\end{minted}
\end{samepage}

\subsection{Περιορισμοί Συμπληρωματικότητας}
\label{subsec:complementarity-constraints}

Οι περιορισμοί συμπληρωματικότητας εξασφαλίζουν ότι μια εντολή γίνεται αποδεκτή μόνο όταν η τιμή αγοράς είναι ευνοϊκή. Υλοποιούνται με τη μέθοδο Big-M:

\begin{samepage}
\begin{minted}[fontsize=\small,linenos,frame=lines]{julia}
# Υπολογισμός του δεξιού μέλους της dual συνθήκης
for (i, order) in enumerate(simple_orders)
    order_id = order_ids[i]
    node_id = string(order.zone)
    time_period = extract_time_period(order.date_time, order_book.periods)
    
    # Ποσότητα με πρόσημο: αρνητική για πώληση, θετική για αγορά
    quantity = order.type == :supply ? -order.quantity : order.quantity
    
    stepwise_dual_rhs[order_id] = @expression(model,
        stepwise_dual[order_id] +
        quantity * market_price[node_id, time_period] -
        quantity * order.price
    )
end

# Περιορισμός dual
@constraint(model, [order_id in order_ids],
    0 <= stepwise_dual_rhs[order_id])

# Big-M μετασχηματισμός για συμπληρωματικότητα
@variable(model, complementarity_aux[order_id in order_ids], Bin)

@constraint(model, [order_id in order_ids],
    stepwise_acceptance[order_id] <= complementarity_aux[order_id] * big_m)

@constraint(model, [order_id in order_ids],
    stepwise_dual_rhs[order_id] <= (1 - complementarity_aux[order_id]) * big_m)
\end{minted}
\end{samepage}

\subsection{Περιορισμός Ισορροπίας Ισχύος}
\label{subsec:mpcc-power-balance}

Ο περιορισμός ισορροπίας ισχύος διαφοροποιείται ανάλογα με το αν υπάρχουν περιορισμοί δικτύου:

\begin{samepage}
\begin{minted}[fontsize=\small,linenos,frame=lines]{julia}
if has_network_constraints
    # Πολυζωνική ισορροπία με ροές μεταφοράς
    @constraint(model, [node_id in nodes, period in periods],
        # Συνεισφορά εντολών: πώληση (αρνητική) + αγορά (θετική)
        sum(stepwise_acceptance[order_ids[i]] * 
            (orders[i].type == :supply ? -orders[i].quantity : orders[i].quantity)
            for i in orders_at_node_period[(node_id, period)]) +
        # Εισροές από άλλες ζώνες
        sum(flow[(z, node_id), period] for z in connected_zones_to[node_id]) -
        # Εκροές προς άλλες ζώνες
        sum(flow[(node_id, z), period] for z in connected_zones_from[node_id]) +
        # Έκτακτη αποκοπή
        load_shed[node_id, period] == 0
    )
else
    # Μονοζωνική ισορροπία (χωρίς ροές)
    @constraint(model, [node_id in nodes, period in periods],
        sum(stepwise_acceptance[order_ids[i]] * 
            (orders[i].type == :supply ? -orders[i].quantity : orders[i].quantity)
            for i in orders_at_node_period[(node_id, period)]) +
        load_shed[node_id, period] == 0
    )
end
\end{minted}
\end{samepage}

\subsection{Αντικειμενική Συνάρτηση}
\label{subsec:mpcc-objective}

Η αντικειμενική συνάρτηση μεγιστοποιεί το οικονομικό πλεόνασμα, με ποινή για την αποκοπή φορτίου:

\begin{samepage}
\begin{minted}[fontsize=\small,linenos,frame=lines]{julia}
load_shed_penalty = 10000.0  # €/MWh

@objective(model, Max,
    # Πλεόνασμα από εντολές
    sum(stepwise_acceptance[order_ids[i]] *
        (order.type == :supply ? 
            -order.quantity * order.price :  # Κόστος παραγωγής
             order.quantity * order.price)   # Αξία κατανάλωσης
        for (i, order) in enumerate(simple_orders)) -
    # Ποινή αποκοπής φορτίου
    sum(load_shed_penalty * load_shed[n, p] 
        for n in order_book.nodes, p in order_book.periods)
)
\end{minted}
\end{samepage}

\subsection{Δομή Αποτελέσματος}
\label{subsec:mpcc-result}

Τα αποτελέσματα της εκκαθάρισης επιστρέφονται στη δομή \texttt{MPCCResult}:

\begin{samepage}
\begin{minted}[fontsize=\small,linenos,frame=lines]{julia}
struct MPCCResult
    status::Symbol                                # :optimal, :infeasible, κλπ
    objective_value::Float64                      # Τιμή αντικειμενικής
    market_prices::Dict{String,Dict{String,Float64}}  # Τιμές [zone][period]
    stepwise_acceptance::Dict{String,Float64}     # Αποδοχές εντολών
    block_acceptance::Dict{String,Float64}        # Αποδοχές block
    block_activation::Dict{String,Float64}        # Ενεργοποιήσεις block
    transmission_flows::Dict{String,Dict{String,Float64}}  # Ροές
    solve_time::Float64                           # Χρόνος επίλυσης
    total_time::Float64                           # Συνολικός χρόνος
    solver_name::String                           # Όνομα επιλυτή
    message::String                               # Μήνυμα
end
\end{minted}
\end{samepage}


% =============================================================================
\section{Μοντελοποίηση Δικτύου}
\label{sec:network-modeling}

Η μοντελοποίηση των περιορισμών δικτύου υλοποιείται στο module \texttt{Network.jl}. Το module παρέχει δύο προσεγγίσεις: τη βασισμένη σε γραμμές μεταφοράς (\texttt{NetworkTopology}) και τη βασισμένη σε ζώνες προσφοράς (\texttt{TransferCapacity}). Η δεύτερη προσέγγιση είναι η συνιστώμενη για χρήση με δεδομένα ENTSO-E.

\subsection{Δομή Διαθέσιμης Ικανότητας Μεταφοράς}
\label{subsec:transfer-capacity-struct}

Η δομή \texttt{TransferCapacity} αναπαριστά την ικανότητα μεταφοράς μεταξύ ζωνών προσφοράς:

\begin{samepage}
\begin{minted}[fontsize=\small,linenos,frame=lines]{julia}
struct TransferCapacity
    bidding_zones::Vector{String}    # Όλες οι ζώνες
    time_periods::Vector{String}     # Χρονικές περίοδοι
    # Ικανότητα στην κατεύθυνση source → sink
    capacity_forward::Dict{Tuple{String,String,String},Float64}
    # Ικανότητα στην κατεύθυνση sink → source
    capacity_backward::Dict{Tuple{String,String,String},Float64}
end
\end{minted}
\end{samepage}

\subsection{Ανάκτηση Δεδομένων ENTSO-E}
\label{subsec:entsoe-data-retrieval}

Η δημιουργία της δομής \texttt{TransferCapacity} από δεδομένα ENTSO-E γίνεται με την ακόλουθη συνάρτηση:

\begin{samepage}
\begin{minted}[fontsize=\small,linenos,frame=lines]{julia}
function create_transfer_capacity_from_entsoe(date::Date, 
                                              bidding_zones::Vector{String}=String[])
    zone_filter = isempty(bidding_zones) ? "" :
        "AND (out_map_code IN ('$(join(bidding_zones, "','"))') " *
        "OR in_map_code IN ('$(join(bidding_zones, "','"))'))"
    
    query = """
    SELECT
        out_map_code as source_zone,
        in_map_code as sink_zone,
        EXTRACT(HOUR FROM date_time_utc) + 1 as time_period,
        capacity_mw as capacity
    FROM entsoe.offered_transfer_capacities_implicit
    WHERE DATE(date_time_utc) = \$1
    $zone_filter
    ORDER BY out_map_code, in_map_code, date_time_utc
    """
    
    df = sql2df(query, [date])
    
    if nrow(df) == 0
        error("No transfer capacity data found for $date")
    end
    
    return build_transfer_capacity_from_dataframe(df)
end
\end{minted}
\end{samepage}

\subsection{Προσθήκη Περιορισμών στο Μοντέλο}
\label{subsec:adding-network-constraints}

Οι περιορισμοί ικανότητας μεταφοράς προστίθενται στο μοντέλο JuMP ως εξής:

\begin{samepage}
\begin{minted}[fontsize=\small,linenos,frame=lines]{julia}
function add_transfer_capacity_constraints!(model::Model, 
                                            transfer_capacity::TransferCapacity, 
                                            FLOW)
    zones = transfer_capacity.bidding_zones
    periods = transfer_capacity.time_periods
    cap_fwd = transfer_capacity.capacity_forward
    cap_bwd = transfer_capacity.capacity_backward
    
    # Περιορισμοί για κάθε ζεύγος ζωνών και περίοδο
    # Θετική ροή: source → sink, περιορίζεται από capacity_forward
    # Αρνητική ροή: sink → source, περιορίζεται από capacity_backward
    @constraint(model, 
        [source in zones, sink in zones, t in periods; source != sink],
        -get(cap_bwd, (source, sink, t), 0.0) <= 
            FLOW[source, sink, t] <= 
            get(cap_fwd, (source, sink, t), 0.0)
    )
    
    return model
end
\end{minted}
\end{samepage}

\subsection{Βοηθητικές Συναρτήσεις}
\label{subsec:network-utilities}

Το module παρέχει βοηθητικές συναρτήσεις για την εξαγωγή πληροφοριών συνδεσιμότητας:

\begin{samepage}
\begin{minted}[fontsize=\small,linenos,frame=lines]{julia}
# Επιστροφή ζευγών ζωνών με διαθέσιμη ικανότητα μεταφοράς
function get_zone_pairs(tc::TransferCapacity)
    pairs = Set{Tuple{String,String}}()
    for (source, sink, _) in keys(tc.capacity_forward)
        if get(tc.capacity_forward, (source, sink, "1"), 0.0) > 0
            push!(pairs, (source, sink))
        end
    end
    return collect(pairs)
end

# Επιστροφή συνδεδεμένων ζωνών για μια δεδομένη ζώνη
function get_connected_zones(tc::TransferCapacity, zone::String)
    connected = Set{String}()
    for (source, sink, _) in keys(tc.capacity_forward)
        if source == zone
            push!(connected, sink)
        elseif sink == zone
            push!(connected, source)
        end
    end
    return collect(connected)
end
\end{minted}
\end{samepage}


% =============================================================================
\section{Ολοκλήρωση και Κύρια Ροή Εργασιών}
\label{sec:integration}

Η ολοκλήρωση όλων των συνιστωσών επιτυγχάνεται μέσω δύο κύριων συναρτήσεων: η \texttt{generate\_energy\_prices} για μονοζωνική εκκαθάριση και η \texttt{run\_multi\_zone\_market\_clearing} για πολυζωνική εκκαθάριση με περιορισμούς διασυνοριακής μεταφοράς.

\subsection{Μονοζωνική Εκκαθάριση}
\label{subsec:single-zone-clearing}

Η συνάρτηση \texttt{generate\_energy\_prices} υλοποιεί την πλήρη ροή εργασιών για μια ζώνη προσφοράς:

\begin{samepage}
\begin{minted}[fontsize=\small,linenos,frame=lines]{julia}
function generate_energy_prices(bidding_zone::String, date::Date;
                                order_method::Symbol=:uc_based,
                                markup_factor::Float64=1.1,
                                save_to_db::Bool=true)
    
    # Βήμα 1: Δημιουργία βιβλίου εντολών
    if order_method == :uc_based
        # Επίλυση UC και μετατροπή σε εντολές
        uc_result = solve_unit_commitment(bidding_zone, date)
        bids_result = generate_market_orders_from_uc(
            bidding_zone, date;
            markup_factor=markup_factor,
            bidding_strategy=:committed_only
        )
        order_book = create_typed_order_book(bids_result, bidding_zone, date)
    elseif order_method == :alternative
        # Εναλλακτική μέθοδος (ταχύτερη)
        result = create_adjusted_order_book(bidding_zone, date)
        order_book = result.order_book
    end
    
    # Βήμα 2: Εκκαθάριση αγοράς
    mpcc_result = solve_mpcc_market_clearing(order_book)
    
    # Βήμα 3: Αποθήκευση αποτελεσμάτων
    if save_to_db && mpcc_result.status == :optimal
        save_energy_prices(mpcc_result, bidding_zone, date)
    end
    
    return mpcc_result
end
\end{minted}
\end{samepage}

\subsection{Πολυζωνική Εκκαθάριση}
\label{subsec:multi-zone-clearing}

Η συνάρτηση \texttt{run\_multi\_zone\_market\_clearing} εκτελεί ταυτόχρονη εκκαθάριση πολλαπλών ζωνών με περιορισμούς ATC:

\begin{samepage}
\begin{minted}[fontsize=\small,linenos,frame=lines]{julia}
function run_multi_zone_market_clearing(date::Date;
                                        zones::Vector{String}=String[],
                                        order_method::Symbol=:alternative,
                                        save_to_db::Bool=true)
    
    # Ανακάλυψη διαθέσιμων ζωνών αν δεν καθορίστηκαν
    if isempty(zones)
        zones = get_zones_with_transfer_capacity(date)
    end
    
    println("🌍 Running multi-zone clearing for $(length(zones)) zones")
    
    # Δημιουργία πολυζωνικού βιβλίου εντολών
    order_book = create_multi_zone_order_book(date, zones; 
                                              order_method=order_method)
    
    # Εκκαθάριση αγοράς
    result = solve_mpcc_market_clearing(order_book)
    
    # Αποθήκευση αποτελεσμάτων
    if save_to_db && result.status == :optimal
        for zone in zones
            save_energy_prices(result, zone, date)
        end
        save_transmission_flows(result, date)
    end
    
    # Εκτύπωση αποτελεσμάτων
    println("💰 Zonal Clearing Prices (average):")
    for zone in zones
        prices = values(result.market_prices[zone])
        avg = sum(prices) / length(prices)
        println("   $zone: €$(round(avg, digits=2))/MWh")
    end
    
    return result
end
\end{minted}
\end{samepage}

\paragraph{Παράλληλη Εκτέλεση UC}
Κατά την πολυζωνική εκκαθάριση με τη μέθοδο \texttt{:uc\_based}, η επίλυση Unit Commitment εκτελείται ανεξάρτητα για κάθε ζώνη. Η παράμετρος \texttt{parallel=true} ενεργοποιεί την ταυτόχρονη εκτέλεση αυτών των ανεξάρτητων υπολογισμών μέσω της συνάρτησης \texttt{pmap} του module \texttt{Distributed.jl}. Κάθε ζώνη ανατίθεται σε έναν ξεχωριστό εργάτη (worker process), ο οποίος εκτελεί πλήρως τον κύκλο ανάκτησης δεδομένων, επίλυσης UC και μετατροπής σε εντολές αγοράς. Η μέθοδος \texttt{pmap} εξασφαλίζει δυναμική κατανομή εργασιών: μόλις ένας εργάτης ολοκληρώσει μια ζώνη, αναλαμβάνει την επόμενη διαθέσιμη.

Σε περίπτωση που δεν υπάρχουν διαθέσιμοι εργάτες (\texttt{nworkers() == 0}), το σύστημα υποβαθμίζεται αυτόματα σε σειριακή εκτέλεση χωρίς σφάλμα. Η ασφάλεια νημάτων (thread safety) εξασφαλίζεται από το γεγονός ότι κάθε εργάτης λειτουργεί σε ξεχωριστό address space και τα κλειδιά κρυφής μνήμης είναι ειδικά ανά ζώνη, ενώ οι εγγραφές στη βάση δεδομένων χρησιμοποιούν συναλλαγές PostgreSQL. Η εκκαθάριση MPCC εκτελείται πάντα σειριακά μετά την ολοκλήρωση όλων των UC, καθώς απαιτεί το ενοποιημένο βιβλίο εντολών όλων των ζωνών.

\subsection{Αποθήκευση Αποτελεσμάτων}
\label{subsec:saving-results}

Τα αποτελέσματα αποθηκεύονται στη βάση δεδομένων μέσω των συναρτήσεων \texttt{save\_energy\_prices} και \texttt{save\_transmission\_flows}:

\begin{samepage}
\begin{minted}[fontsize=\small,linenos,frame=lines]{julia}
function save_energy_prices(result::MPCCResult, bidding_zone::String, date::Date)
    # Εξασφάλιση ύπαρξης πίνακα
    ensure_energy_prices_table()
    
    # Προετοιμασία δεδομένων
    prices = result.market_prices[bidding_zone]
    
    withdb() do conn
        for (period, price) in prices
            hour = parse(Int, period)
            execute(conn, """
                INSERT INTO simulations.energy_prices 
                    (bidding_zone, date, hour, price, created_at)
                VALUES (\$1, \$2, \$3, \$4, NOW())
                ON CONFLICT (bidding_zone, date, hour) 
                DO UPDATE SET price = EXCLUDED.price, created_at = NOW()
            """, [bidding_zone, date, hour, price])
        end
    end
end
\end{minted}
\end{samepage}

\subsection{Επεξεργασία Κατά Παρτίδες}
\label{subsec:batch-processing}

Για την επεξεργασία πολλαπλών ημερομηνιών, παρέχονται οι συναρτήσεις \texttt{generate\_energy\_prices\_for\_date\_range} και \texttt{run\_multi\_zone\_for\_date\_range}:

\begin{samepage}
\begin{minted}[fontsize=\small,linenos,frame=lines]{julia}
function run_multi_zone_for_date_range(start_date::Date, end_date::Date;
                                       zones::Vector{String}=String[],
                                       order_method::Symbol=:alternative,
                                       force_rerun::Bool=false,
                                       parallel::Bool=false)
    results = Dict{Date,MPCCResult}()

    for date in start_date:Day(1):end_date
        try
            result = run_multi_zone_market_clearing(date;
                zones=zones, order_method=order_method,
                force_rerun=force_rerun, parallel=parallel)
            results[date] = result
        catch e
            @warn "Failed to process $date: $e"
        end
    end
    return results
end
\end{minted}
\end{samepage}

\subsection{Παράδειγμα Χρήσης}
\label{subsec:usage-example}

Ένα πλήρες παράδειγμα χρήσης του συστήματος για τον υπολογισμό τιμών στην Ελλάδα:

\begin{samepage}
\begin{minted}[fontsize=\small,linenos,frame=lines]{julia}
using Euphemia
using Dates

# Μονοζωνική εκκαθάριση για την Ελλάδα
prices = generate_energy_prices("GR", Date(2024, 6, 15);
    order_method = :uc_based,
    markup_factor = 1.1,
    save_to_db = true
)

# Πολυζωνική εκκαθάριση Ελλάδας-Βουλγαρίας-Ιταλίας
result = run_multi_zone_market_clearing(Date(2024, 6, 15);
    zones = ["GR", "BG", "IT"],
    order_method = :alternative,
    save_to_db = true
)

# Εκτύπωση διασυνοριακών ροών
for (flow_id, flows) in result.transmission_flows
    avg_flow = sum(values(flows)) / length(flows)
    println("$flow_id: $(round(avg_flow, digits=1)) MW average")
end
\end{minted}
\end{samepage}


% =============================================================================
\section{Αυτοματοποιημένη Εκτέλεση}
\label{sec:automated-execution}

Για τη συστηματική παραγωγή τιμών σε ημερήσια βάση, αναπτύχθηκε υποδομή αυτοματοποίησης βασισμένη σε GitHub Actions. Δύο ροές εργασιών (workflows) εκτελούνται σε αυτο-φιλοξενούμενο μηχάνημα (self-hosted runner):

\begin{itemize}
\item \texttt{generate-energy-prices.yml}: Μονοζωνική εκκαθάριση, προγραμματισμένη στις 02:00 UTC ημερησίως.
\item \texttt{generate-multi-zone-prices.yml}: Πολυζωνική εκκαθάριση με διασυνοριακούς περιορισμούς, προγραμματισμένη στις 03:00 UTC ημερησίως.
\end{itemize}

Η πολυζωνική ροή εργασιών εκτελεί τα εξής βήματα: εγκατάσταση εξαρτήσεων Julia, ρύθμιση σύνδεσης βάσης δεδομένων, εκκίνηση παράλληλων εργατών (workers), εκτέλεση του σεναρίου \texttt{bin/multi\_zone\_main.jl} με τις κατάλληλες παραμέτρους, και αποστολή ειδοποιήσεων μέσω Slack σε περίπτωση επιτυχίας ή αποτυχίας. Κατά τη χρήση του εμπορικού επιλυτή Gurobi, ο αριθμός παράλληλων εργατών περιορίζεται αυτόματα σε 2, λόγω του ορίου ταυτόχρονων συνεδριών της άδειας WLS (Web License Service). Κάθε ροή εργασιών υποστηρίζει τόσο χειροκίνητη ενεργοποίηση (\texttt{workflow\_dispatch}) με παραμετροποιήσιμο εύρος ημερομηνιών, όσο και αυτόματη εκτέλεση μέσω χρονοπρογραμματισμού (\texttt{cron}).
